%% parts/p01/symbol-registry.tex --- scoped notation and lookup tables.
\chapter{Symbol Registry, Binding, and Precedence}
\chaptermark{Symbol Registry}
\label{sec:p01-symbol-registry}

The registry is normative for this handbook, not a claim that the literature
uses one notation.  Each row fixes a glyph, its formal or grammatical
standing, its scope, and (when it is an operator) its precedence.  Normally a
glyph receives another role only in a disjoint scope declared before use.  A
named primary calculus, logic, model, or method instead keeps its published
syntax: every overlapping local role must be typed and registered honestly,
not renamed to fit handbook notation.
\index{symbol registry}\index{notation!scope}

The editorial defaults are vertical bars for set cardinality,
\(\dom\mathcal M\) for a structure domain,
\(\mathcal M\equiv\mathcal N\) for elementary equivalence,
\(\varphi\Leftrightarrow\psi\) for formula equivalence,
\(\defeq\) for a definitional sign, and \(\forces\), with negation
\(\nforces\), for the declared forcing relation.  A source presentation's
explicitly typed local notation governs that presentation.  For example,
the inherited propositional semantics uses \(\varphi\equiv\psi\) for
equality of truth values under every valuation, whereas the displayed
structure operands in \(\mathcal M\equiv\mathcal N\) give the
elementary-equivalence role.  The register records both roles; neither
declaration renames the other presentation.

In the inherited CPL semantics the truth-value set is
\(\mathbb B=\{\mathbf F,\mathbf T\}\).
For numerical Boolean tables, the coding sends \(\mathbf F\) to \(0\)
and \(\mathbf T\) to \(1\); a valuation written with numerical values
uses this coding.  The coding preserves the declared truth functions and
does not identify the two sets literally.  The Boolean-table section uses
the numerical presentation; the source CPL valuation and substitution
sections retain the labelled values.

\section{Generated master register}
\label{sec:p01-generated-symbols}

The following scope tables are generated from
\texttt{registry/symbols.csv}.  Their role and scope columns are part of each
notation's declaration; the prose notes record local restrictions and
typed collision boundaries across model theory, computability and arithmetic,
set theory, and nonclassical logic.

% GENERATED by tools/symbols2tex.py from handbook/registry/symbols.csv -- do not hand-edit.
\begin{landscape}
\Needspace{.4\textheight}
\subsection*{Global scope}
\begingroup\scriptsize\setlength{\tabcolsep}{3pt}
\begin{longtable}{@{}p{0.9in}p{1.3in}p{1.0in}p{0.5in}p{0.3in}p{0.4in}p{0.35in}p{1.65in}@{}}
\toprule
\textbf{Glyph} & \textbf{Macro} & \textbf{Role} & \textbf{Scope} & \textbf{Prec} & \textbf{First} & \textbf{Reg. page} & \textbf{Principal uses / notes} \\
\midrule
\endfirsthead
\toprule
\textbf{Glyph} & \textbf{Macro} & \textbf{Role} & \textbf{Scope} & \textbf{Prec} & \textbf{First} & \textbf{Reg. page} & \textbf{Principal uses / notes} \\
\midrule
\endhead
\midrule\multicolumn{8}{r}{\emph{continued}}\\\endfoot
\bottomrule\endlastfoot
\phantomsection\label{entry:symbol-registry-00002}\(\mathcal L\) & \texttt{\detokenize{\Lang}} & signature & global & n/a & P2.2 & \thepage & Nonlogical symbols with declared arities or sorts; governs text after a language declaration. \\
\phantomsection\label{entry:symbol-registry-00003}\(x,y,z\) & \texttt{\detokenize{x,}\allowbreak\detokenize{y,}\allowbreak\detokenize{z}} & object-language variables & global & n/a & P2.2 & \thepage & Variables occurring in terms and formulas. \\
\phantomsection\label{entry:symbol-registry-00004}\(r,t\) & \texttt{\detokenize{r,}\allowbreak\detokenize{t}} & term metavariables & global & n/a & P2.2 & \thepage & Members of \(\mathrm{Term}(\mathcal L)\). \\
\phantomsection\label{entry:symbol-registry-00005}\(\varphi,\psi,\chi\) & \texttt{\detokenize{\varphi,}\allowbreak\detokenize{\psi,}\allowbreak\detokenize{\chi}} & formula metavariables & global & n/a & P2.1 & \thepage & Members of the active formula set.  A source-bound presentation may instead declare its own formula metavariables locally; the source declaration governs that presentation. \\
\phantomsection\label{entry:symbol-registry-00006}\(\neg\varphi\) & \texttt{\detokenize{\neg}} & object-language negation & global & 3 & P2.1 & \thepage & Primitive unary connective; prefix and governs the immediately following formula. \\
\phantomsection\label{entry:symbol-registry-00007}\(\varphi\land\psi\) & \texttt{\detokenize{\land}} & object-language conjunction & global & 4 & P2.1 & \thepage & Primitive or explicitly defined binary connective in the active language; left associative in handbook prose. \\
\phantomsection\label{entry:symbol-registry-00008}\(\varphi\lor\psi\) & \texttt{\detokenize{\lor}} & object-language disjunction & global & 5 & P2.1 & \thepage & Primitive or explicitly defined binary connective in the active language; left associative in handbook prose. \\
\phantomsection\label{entry:symbol-registry-00009}\(\varphi\to\psi\) & \texttt{\detokenize{\to}} & object-language implication & global & 6 & P2.1 & \thepage & Primitive or explicitly defined binary connective in the active language; right associative in handbook prose. \\
\phantomsection\label{entry:symbol-registry-00010}\(\sigma,\tau\) & \texttt{\detokenize{\sigma,}\allowbreak\detokenize{\tau}} & sentence metavariables & global & n/a & P2.2 & \thepage & Members of \(\mathrm{Sent}(\mathcal L)\). \\
\phantomsection\label{entry:symbol-registry-00011}\(\forall x\,\varphi\) & \texttt{\detokenize{\forall}} & universal quantifier & global & 3 & P2.2 & \thepage & Binds occurrences of x in its declared scope; prefix at the unary rank. \\
\phantomsection\label{entry:symbol-registry-00012}\(\exists x\,\varphi\) & \texttt{\detokenize{\exists}} & existential quantifier & global & 3 & P2.2 & \thepage & Binds occurrences of x in its declared scope; prefix at the unary rank. \\
\phantomsection\label{entry:symbol-registry-00014}\(\Gamma,\Delta,\Sigma,\Pi\) & \texttt{\detokenize{\Gamma,}\allowbreak\detokenize{\Delta,}\allowbreak\detokenize{\Sigma,}\allowbreak\detokenize{\Pi}} & formula contexts & global & n/a & P2.1 & \thepage & Finite sets, multisets, or sequences exactly as declared by the active calculus.  Source-bound locally typed uses such as the Buchi alphabet Sigma are recorded separately and are not contexts. \\
\phantomsection\label{entry:symbol-registry-00021}\(T,U\) & \texttt{\detokenize{T,}\allowbreak\detokenize{U}} & theories & global & n/a & P2.2 & \thepage & Sets of sentences in the fixed signature. \\
\phantomsection\label{entry:symbol-registry-00022}\(\mathcal M,\mathcal N\) & \texttt{\detokenize{\mathcal M,}\allowbreak\detokenize{\mathcal N}} & structures & global & n/a & P2.2 & \thepage & Structures for the active signature with nonempty domains unless locally declared otherwise; in exact Part V, \(\mathcal N\) is the standard structure for the fixed arithmetic signature. \\
\phantomsection\label{entry:symbol-registry-00023}\(\operatorname{dom}X\) & \texttt{\detokenize{\dom}} & domain operator & global & n/a & P2.2 & \thepage & Operand type fixes the structure-domain, function-domain, or name-domain reading; vertical bars are never used for domains. \\
\phantomsection\label{entry:symbol-registry-00024}\(v\) & \texttt{\detokenize{v}} & valuation & global & n/a & P2.1 & \thepage & Map from propositional variables to the active truth-value set in propositional and finite-matrix semantics; its homomorphic extension \(\widehat v\) is fixed by the active semantics. \\
\phantomsection\label{entry:symbol-registry-00025}\(s\) & \texttt{\detokenize{s}} & variable assignment & global & n/a & P2.2 & \thepage & Map from first-order variables to \(\operatorname{dom}\mathcal M\). \\
\phantomsection\label{entry:symbol-registry-00026}\(\mathcal M,s\models\varphi\) & \texttt{\detokenize{\entails}} & satisfaction & global & meta & P2.2 & \thepage & \(\varphi\) is satisfied by \(\mathcal M\) under \(s\). \\
\phantomsection\label{entry:symbol-registry-00027}\(\mathcal M\models\sigma\) & \texttt{\detokenize{\entails}} & truth in a structure & global & meta & P2.2 & \thepage & Sentence \(\sigma\) is true in \(\mathcal M\). \\
\phantomsection\label{entry:symbol-registry-00028}\(\Gamma\models\varphi\) & \texttt{\detokenize{\entails}} & semantic consequence & global & meta & P2.1 & \thepage & Every interpretation and assignment in the declared class satisfying \(\Gamma\) satisfies \(\varphi\). \\
\phantomsection\label{entry:symbol-registry-00029}\(\Gamma\vdash_{S}\varphi\) & \texttt{\detokenize{\derives_{S}}} & derivability (per-system turnstile) & global & meta & P2.1 & \thepage & A derivation exists in the named calculus \(S\); the subscript is mandatory outside a declared local abbreviation. \\
\phantomsection\label{entry:symbol-registry-00030}\(\varphi\Leftrightarrow\psi\) & \texttt{\detokenize{\Leftrightarrow}} & formula equivalence & global & meta & P2.1 & \thepage & Metalinguistic biconditionals are written in prose as iff; \(\Leftrightarrow\) relates formulas and \(\equiv\) is reserved for elementary equivalence. \\
\phantomsection\label{entry:symbol-registry-00031}\(\varphi\leftrightarrow\psi\) & \texttt{\detokenize{\leftrightarrow}} & object-language biconditional & global & 7 & P2.1 & \thepage & Object-language constructor; rank 7 of the extended precedence ladder. \\
\phantomsection\label{entry:symbol-registry-00032}\(a\mathrel{:\equiv}b\) & \texttt{\detokenize{\defeq}} & definitional abbreviation & global & meta & P1.2 & \thepage & The left expression is defined by the right expression; this sign is metalinguistic. \\
\phantomsection\label{entry:symbol-registry-00033}\(\varphi[t/x]\) & \texttt{\detokenize{[t/}\allowbreak\detokenize{x]}} & capture-avoiding substitution & global & n/a & P2.2 & \thepage & Substitution of \(t\) for free occurrences of \(x\) in \(\varphi\). \\
\phantomsection\label{entry:symbol-registry-00034}\(\operatorname{FV}(\varphi),\operatorname{FV}(\Gamma),\operatorname{FV}(M),\operatorname{FV}(\alpha)\) & \texttt{\detokenize{\FV}} & free-variable operation & global & n/a & P2.2 & \thepage & Variables with a free occurrence in the locally declared formula, context, term, or type syntax; context action takes the union over members. \\
\phantomsection\label{entry:symbol-registry-00035}\(\operatorname{Mod}(T)\) & \texttt{\detokenize{\Mod}} & model class & global & n/a & P2.2 & \thepage & Structures satisfying every sentence of \(T\). \\
\phantomsection\label{entry:symbol-registry-00036}\(\operatorname{Th}(\mathcal M)\) & \texttt{\detokenize{\Th}} & theory of a structure & global & n/a & P2.2 & \thepage & Sentences true in \(\mathcal M\). \\
\phantomsection\label{entry:symbol-registry-00037}\(X^*\) & \texttt{\detokenize{{}^*}} & finite iteration or closure & global & n/a & P1.3 & \thepage & For an alphabet this is finite-word Kleene star; for a relation it is reflexive-transitive closure. Operand type fixes the reading. \\
\phantomsection\label{entry:symbol-registry-00038}\(\bot,\top\) & \texttt{\detokenize{\bottom,}\allowbreak\detokenize{\topv}} & falsum and verum or declared CLL nullary units & global & n/a & P2.1 & \thepage & Outside the P8.6 CLL grammar these are bottom and the defined top abbreviation; in that grammar \(\bot\) is the par unit and \(\top\) is the with unit, with the ambient syntax fixing the role. \\
\phantomsection\label{entry:symbol-registry-00039}\(\mathbb N\) & \texttt{\detokenize{\N}} & set of natural numbers & global & n/a & P1.3 & \thepage & Contains 0 under the Part-1 natural-number convention. \\
\phantomsection\label{entry:symbol-registry-00040}\(\mathcal P(A)\) & \texttt{\detokenize{\mathcal P}} & power set & global & n/a & P1.3 & \thepage & Set of all subsets of the set A; set existence comes from the ambient foundation. \\
\phantomsection\label{entry:symbol-registry-00041}\(|X|\) & \texttt{\detokenize{\abs}} & cardinality & global & n/a & P1.3 & \thepage & Vertical bars mean cardinality only; domains use \(\operatorname{dom}\). \\
\phantomsection\label{entry:symbol-registry-00042}\(\mathcal M\equiv\mathcal N\) & \texttt{\detokenize{\equiv}} & elementary equivalence & global & meta & P2.2 & \thepage & The structures satisfy the same sentences; this is the sole global role of \(\equiv\). \\
\phantomsection\label{entry:symbol-registry-00043}\(\mathcal M\preccurlyeq\mathcal N\) & \texttt{\detokenize{\preccurlyeq}} & elementary substructure & global & meta & P2.2 & \thepage & Every formula with parameters from \(\mathcal M\) has the same truth value in \(\mathcal M\) and \(\mathcal N\). \\
\phantomsection\label{entry:symbol-registry-00048}\(\mathcal M\cong\mathcal N\) & \texttt{\detokenize{\cong}} & isomorphism & global & meta & P2.2 & \thepage & Structure-preserving bijection; \(\simeq\) remains separately reserved for Kleene equality. \\
\phantomsection\label{entry:symbol-registry-00075}\(\Gamma\Rightarrow\Delta\) & \texttt{\detokenize{\Rightarrow}} & sequent & global & meta & P2.1 & \thepage & Metalevel punctuation between antecedent and succedent contexts. \\
\phantomsection\label{entry:symbol-registry-00089}\(>\) & \texttt{\detokenize{>}} & strict ordering & global & 2 & P1.2 & \thepage & Contract list; strict dual of \(\leq\); ``positive'' means \(>0\) by the natural-numbers convention. \\
\phantomsection\label{entry:symbol-registry-00114}\(\Gamma\models_{S}\varphi\) & \texttt{\detokenize{\entails_{S}}} & semantic consequence (per-system turnstile) & global & meta & P1.3 & \thepage & The subscript names the semantics (\(\models_{\mathcal M}\) matrix, \(\models_{\mathrm{K3}}\), \(\models_{\mathrm{LP}}\), \(\models_{\mathsf{S4}}\), ...). \\
\phantomsection\label{entry:symbol-registry-00122}\(\mathcal M\models T\) & \texttt{\detokenize{\entails}} & model of a theory & global & meta & P2.2 & \thepage & \(\mathcal M\) satisfies every member of \(T\). \\
\phantomsection\label{entry:symbol-registry-00149}\(\mathsf{Var}\) & \texttt{\detokenize{\mathsf{Var}}} & first-order variable alphabet & global & n/a & P2.2 & \thepage & The fixed countably infinite object-language variable set; distinct from the occurrence-set operator \(\operatorname{Var}\). \\
\phantomsection\label{entry:symbol-registry-00150}\(\mathcal F,\mathcal R\) & \texttt{\detokenize{\mathcal F,}\allowbreak\detokenize{\mathcal R}} & function- and relation-symbol sets & global & n/a & P2.2 & \thepage & The disjoint nonlogical-symbol families of a first-order signature; constants are zero-ary function symbols. \\
\phantomsection\label{entry:symbol-registry-00151}\(\operatorname{ar}_{\mathcal F},\operatorname{ar}_{\mathcal R}\) & \texttt{\detokenize{\operatorname{ar}_{\mathcal F},}\allowbreak\detokenize{\operatorname{ar}_{\mathcal R}}} & signature arity maps & global & n/a & P2.2 & \thepage & Maps the function- and relation-symbol families to their finite arities. \\
\phantomsection\label{entry:symbol-registry-00152}\(\operatorname{Term}(\Lang),\operatorname{Fm}(\Lang),\operatorname{Sent}(\Lang)\) & \texttt{\detokenize{\operatorname{Term},}\allowbreak\detokenize{\operatorname{Fm},}\allowbreak\detokenize{\operatorname{Sent}}} & first-order term formula and sentence sets & global & n/a & P2.2 & \thepage & Least syntactic sets generated by the active first-order grammar; sentences are formulas with no free variables. \\
\phantomsection\label{entry:symbol-registry-00153}\(\operatorname{Var}(e)\) & \texttt{\detokenize{\operatorname{Var}}} & occurring-variable set & global & n/a & P2.2 & \thepage & All variables occurring in the declared term, formula, literal, clause, or finite clause set \(e\), whether free or bound; for a finite clause set the operator is the union of the occurrence sets of its clauses. \\
\phantomsection\label{entry:symbol-registry-00154}\(\varphi=_{\alpha}\psi\) & \texttt{\detokenize{=_{\alpha}}} & alpha-equivalence & global & meta & P2.2 & \thepage & Least constructor congruence generated by capture-avoiding bound-variable renaming; distinct from identity formula equivalence and elementary equivalence. \\
\phantomsection\label{entry:symbol-registry-00155}\(s[x\mapsto a]\) & \texttt{\detokenize{[x\mapsto a]}} & assignment update & global & n/a & P2.2 & \thepage & Agrees with assignment \(s\) except that variable \(x\) is sent to \(a\). \\
\phantomsection\label{entry:symbol-registry-00156}\(\operatorname{val}^{\mathcal M}_s(t)\) & \texttt{\detokenize{\operatorname{val}}} & term value & global & n/a & P2.2 & \thepage & Recursive denotation of the first-order term \(t\) in structure \(\mathcal M\) under assignment \(s\). \\
\phantomsection\label{entry:symbol-registry-00157}\(X\mathbin{\upharpoonright}Y\) & \texttt{\detokenize{\upharpoonright}} & restriction or reduct operator & global & n/a & P2.2 & \thepage & Operand types distinguish function or set restriction, a structure reduct, and the P8.3 public-announcement submodel induced by the surviving worlds. \\
\phantomsection\label{entry:symbol-registry-00158}\(\operatorname{Cn}(T)\) & \texttt{\detokenize{\operatorname{Cn}}} & semantic consequence closure & global & n/a & P2.2 & \thepage & Sentences of the active signature semantically entailed by theory \(T\). \\
\phantomsection\label{entry:symbol-registry-00182}\(\mathbb R\) & \texttt{\detokenize{\mathbb R}} & real ordered field & global & n/a & P2.3 & \thepage & The ordinary field of real numbers; the P2.3 categoricity claim names its ordered-field structure. \\
\phantomsection\label{entry:symbol-registry-00183}\(\mathbb Q\) & \texttt{\detokenize{\mathbb Q}} & rational ordered field & global & n/a & P5.1 & \thepage & The ordinary field of rational numbers; used as an order type in P5 and as the scalar domain of linear rational arithmetic in P12.2. \\
\phantomsection\label{entry:symbol-registry-00225}\(\mathbb Z\) & \texttt{\detokenize{\mathbb Z}} & set of integers & global & n/a & P3.1 & \thepage & Integers; used for tape positions and integer-coefficient polynomials. \\
\phantomsection\label{entry:symbol-registry-00226}\(f:A\rightharpoonup B\) & \texttt{\detokenize{\rightharpoonup}} & partial-function type arrow & global & n/a & P1.3 & \thepage & The function may be undefined outside a declared subset of \(A\). \\
\end{longtable}
\endgroup
\Needspace{.4\textheight}
\subsection*{Part 1 scopes}
\begingroup\scriptsize\setlength{\tabcolsep}{3pt}
\begin{longtable}{@{}p{0.9in}p{1.3in}p{1.0in}p{0.5in}p{0.3in}p{0.4in}p{0.35in}p{1.65in}@{}}
\toprule
\textbf{Glyph} & \textbf{Macro} & \textbf{Role} & \textbf{Scope} & \textbf{Prec} & \textbf{First} & \textbf{Reg. page} & \textbf{Principal uses / notes} \\
\midrule
\endfirsthead
\toprule
\textbf{Glyph} & \textbf{Macro} & \textbf{Role} & \textbf{Scope} & \textbf{Prec} & \textbf{First} & \textbf{Reg. page} & \textbf{Principal uses / notes} \\
\midrule
\endhead
\midrule\multicolumn{8}{r}{\emph{continued}}\\\endfoot
\bottomrule\endlastfoot
\phantomsection\label{entry:symbol-registry-00079}\(\leq\) & \texttt{\detokenize{\leq}} & canonical ordering & P1-P4.2 & 2 & P1.2 & \thepage & Ambient order on declared numerical and order-theoretic objects; P4.3 is excluded because bare \(\leq\) there is object-language arithmetic syntax. \\
\phantomsection\label{entry:symbol-registry-00085}\(<\) & \texttt{\detokenize{<}} & canonical strict ordering & P1-P4.2 & 2 & P1.2 & \thepage & Ambient strict order on declared numerical and order-theoretic objects; P4.3 is excluded because bare \(<\) there is object-language arithmetic syntax. \\
\end{longtable}
\endgroup
\Needspace{.4\textheight}
\subsection*{Part 2 scopes}
\begingroup\scriptsize\setlength{\tabcolsep}{3pt}
\begin{longtable}{@{}p{0.9in}p{1.3in}p{1.0in}p{0.5in}p{0.3in}p{0.4in}p{0.35in}p{1.65in}@{}}
\toprule
\textbf{Glyph} & \textbf{Macro} & \textbf{Role} & \textbf{Scope} & \textbf{Prec} & \textbf{First} & \textbf{Reg. page} & \textbf{Principal uses / notes} \\
\midrule
\endfirsthead
\toprule
\textbf{Glyph} & \textbf{Macro} & \textbf{Role} & \textbf{Scope} & \textbf{Prec} & \textbf{First} & \textbf{Reg. page} & \textbf{Principal uses / notes} \\
\midrule
\endhead
\midrule\multicolumn{8}{r}{\emph{continued}}\\\endfoot
\bottomrule\endlastfoot
\phantomsection\label{entry:symbol-registry-00013}\(\theta,\rho,\eta\) & \texttt{\detokenize{\theta,}\allowbreak\detokenize{\rho,}\allowbreak\detokenize{\eta}} & simultaneous term substitutions & P2.2-P12.3 & n/a & P2.2 & \thepage & Finite-support variable-to-term maps with declared action and composition. \\
\phantomsection\label{entry:symbol-registry-00045}\(\operatorname{ElDiag}(\mathcal M)\) & \texttt{\detokenize{\operatorname{ElDiag}}} & elementary diagram & P2.2-P3 & n/a & P2.2 & \thepage & All parameter sentences true in \(\mathcal M\). \\
\phantomsection\label{entry:symbol-registry-00076}\(\models_{\mathrm{full}}\) & \texttt{\detokenize{\entails_{\mathrm{full}}}} & full second-order consequence & P2.3 & meta & P2.3 & \thepage & Consequence over full standard structures. \\
\phantomsection\label{entry:symbol-registry-00077}\(\models_{\mathrm H}\) & \texttt{\detokenize{\entails_{\mathrm H}}} & Henkin consequence & P2.3 & meta & P2.3 & \thepage & Consequence over Henkin structures. \\
\phantomsection\label{entry:symbol-registry-00121}\(v\models\varphi\) & \texttt{\detokenize{\entails}} & valuation satisfaction & P2.1 & meta & P2.1 & \thepage & The valuation's homomorphic extension designates \(\varphi\). \\
\phantomsection\label{entry:symbol-registry-00123}\(\mathsf{Prop}\) & \texttt{\detokenize{\mathsf{Prop}}} & set of propositional variables & P2.1 & n/a & P2.1 & \thepage & Exactly \(\{p_n:n\in\mathbb N\}\) in the fixed countable CPL language. \\
\phantomsection\label{entry:symbol-registry-00125}\(p_n; P,Q,R,S\) & \texttt{\detokenize{p_n; P,}\allowbreak\detokenize{Q,}\allowbreak\detokenize{R,}\allowbreak\detokenize{S}} & object-language propositional variables & P2.1 & n/a & P2.1 & \thepage & \(p_n\) is the indexed variable schema; uppercase letters denote specified members of \(\mathsf{Prop}\) in examples and diagrams, never formula metavariables. \\
\phantomsection\label{entry:symbol-registry-00126}\(\mathsf{Fm}_{\mathsf{CPL}}\) & \texttt{\detokenize{\mathsf{Fm}_{\mathsf{CPL}}}} & set of CPL formulas & P2.1 & n/a & P2.1 & \thepage & Least set generated by the declared propositional grammar. \\
\phantomsection\label{entry:symbol-registry-00127}\(\mathbb B\) & \texttt{\detokenize{\mathbb B}} & Boolean truth-value set & P2.1 & n/a & P2.1 & \thepage & The two-element set \(\{0,1\}\); distinct from a Boolean algebra unless one is separately declared. \\
\phantomsection\label{entry:symbol-registry-00128}\(c(\varphi)\) & \texttt{\detokenize{c}} & formula complexity & P2.1-P2.2 & n/a & P2.1 & \thepage & Recursively counts constructor occurrences under the active propositional or first-order grammar. \\
\phantomsection\label{entry:symbol-registry-00129}\(\operatorname{Sub}(\varphi)\) & \texttt{\detokenize{\operatorname{Sub}}} & subformula-set operation & P2.1 & n/a & P2.1 & \thepage & Set of subformulas fixed by the recursive CPL definition. \\
\phantomsection\label{entry:symbol-registry-00130}\(\upsilon\) & \texttt{\detokenize{\upsilon}} & uniform propositional substitution & P2.1 & n/a & P2.1 & \thepage & A metalinguistic map \(\mathsf{Prop}\to\mathsf{Fm}_{\mathsf{CPL}}\); not the globally registered sentence metavariables \(\sigma,\tau\). \\
\phantomsection\label{entry:symbol-registry-00131}\(\varphi[\upsilon]\) & \texttt{\detokenize{[\upsilon]}} & uniform-substitution instance & P2.1 & meta & P2.1 & \thepage & Recursive action of the declared uniform substitution on a CPL formula. \\
\phantomsection\label{entry:symbol-registry-00132}\(\overline L\) & \texttt{\detokenize{\overline L}} & literal complement & P2.1-P7.8 & meta & P2.1 & \thepage & Metalinguistic involution exchanging \(p_n\) and \(\neg p_n\); P7.8 reuses it on the extended-resolution sheet; disjoint from numeral bars. \\
\phantomsection\label{entry:symbol-registry-00133}\(\operatorname{code}(L)\) & \texttt{\detokenize{\operatorname{code}}} & propositional-literal ordering code & P2.1 & n/a & P2.1 & \thepage & \(\operatorname{code}(p_n)=2n\) and \(\operatorname{code}(\neg p_n)=2n+1\); fixes one formula string for each clause. \\
\phantomsection\label{entry:symbol-registry-00134}\(\operatorname{form}_{\mathsf{CPL}}(C)\) & \texttt{\detokenize{\operatorname{form}_{\mathsf{CPL}}}} & formula interpretation of a clause & P2.1 & n/a & P2.1 & \thepage & Increasing-code disjunction of a finite clause; the empty clause maps to falsum. \\
\phantomsection\label{entry:symbol-registry-00135}\(\alpha^{\mathrm d}\) & \texttt{\detokenize{^{\mathrm d}}} & NNF dual operation & P2.1 & meta & P2.1 & \thepage & Metalinguistic involution fixing literals and exchanging conjunction with disjunction and falsum with verum. \\
\phantomsection\label{entry:symbol-registry-00136}\(\bar v\) & \texttt{\detokenize{\bar v}} & complemented valuation & P2.1 & n/a & P2.1 & \thepage & Defined by \(\bar v(p)=1-v(p)\); used only in the propositional duality theorem. \\
\phantomsection\label{entry:symbol-registry-00137}\(\langle C\rangle\) & \texttt{\detokenize{\langle C\rangle}} & clone generated by truth functions & P2.1 & meta & P2.1 & \thepage & Intersection of all positive-arity Boolean clones containing C. \\
\phantomsection\label{entry:symbol-registry-00138}\(\mathsf T_0,\mathsf T_1,\mathsf M,\mathsf L,\mathsf S\) & \texttt{\detokenize{\mathsf T_0,}\allowbreak\detokenize{\mathsf T_1,}\allowbreak\detokenize{\mathsf M,}\allowbreak\detokenize{\mathsf L,}\allowbreak\detokenize{\mathsf S}} & Post maximal clones & P2.1 & n/a & P2.1 & \thepage & Respectively 0-preserving, 1-preserving, monotone, affine, and self-dual positive-arity Boolean functions. \\
\phantomsection\label{entry:symbol-registry-00139}\(\mathbb F_2\) & \texttt{\detokenize{\mathbb F_2}} & two-element field & P2.1 & n/a & P2.1 & \thepage & Field used for affine Boolean functions and XOR constraints. \\
\phantomsection\label{entry:symbol-registry-00140}\(\mathrm{NAND}(\varphi,\psi),\mathrm{NOR}(\varphi,\psi)\) & \texttt{\detokenize{\mathrm{NAND},}\allowbreak\detokenize{\mathrm{NOR}}} & metalinguistic Boolean operator names & P2.1 & n/a & P2.1 & \thepage & Formula abbreviations for negated conjunction and negated disjunction; not active binary glyphs. \\
\phantomsection\label{entry:symbol-registry-00141}\(\mathsf{ND}_{\mathsf{CPL}},\mathsf{LK}_{\mathsf{CPL}},\mathsf H_{\mathsf{CPL}}\) & \texttt{\detokenize{\mathsf{ND}_{\mathsf{CPL}},}\allowbreak\detokenize{\mathsf{LK}_{\mathsf{CPL}},}\allowbreak\detokenize{\mathsf H_{\mathsf{CPL}}}} & CPL calculus identifiers & P2.1 & n/a & P2.1 & \thepage & Names the registered natural-deduction sequent and Hilbert calculi; derivability still carries the calculus as a turnstile subscript. \\
\phantomsection\label{entry:symbol-registry-00142}\(\operatorname{supp}(\Gamma)\) & \texttt{\detokenize{\operatorname{supp}}} & support of a finite multiset & P2.1-P7.1 & n/a & P2.1 & \thepage & Formulas having positive multiplicity in the declared propositional or first-order finite-multiset context; reused in P7.1. \\
\phantomsection\label{entry:symbol-registry-00143}\(\operatorname{tr}_{\mathsf H}(\varphi)\) & \texttt{\detokenize{\operatorname{tr}_{\mathsf H}}} & full-language-to-Hilbert translation & P2.1 & n/a & P2.1 & \thepage & Metalinguistic recursive map into the negation-implication fragment; avoids the globally reserved sentence metavariable \(\tau\). \\
\phantomsection\label{entry:symbol-registry-00144}\(\mathsf P\) & \texttt{\detokenize{\mathsf P}} & deterministic polynomial-time language class & P2.1 & n/a & P2.1 & \thepage & Local propositional-complexity standing. \\
\phantomsection\label{entry:symbol-registry-00145}\(\mathsf{NP}\) & \texttt{\detokenize{\mathsf{NP}}} & nondeterministic polynomial-time language class & P2.1-P12.1 & n/a & P2.1 & \thepage & Equivalently polynomially witnessed languages under the declared encoding; reused in P4.6 P7.8 and the P12.1 proof-certificate comparison. \\
\phantomsection\label{entry:symbol-registry-00146}\(\mathsf{coNP}\) & \texttt{\detokenize{\mathsf{coNP}}} & complements of NP languages & P2.1-P7.8 & n/a & P2.1 & \thepage & Decision-language class under the declared encoding; reused in P4.6 and P7.8. \\
\phantomsection\label{entry:symbol-registry-00147}\(\mathsf{SAT}\) & \texttt{\detokenize{\mathsf{SAT}}} & satisfiable propositional encodings & P2.1 & n/a & P2.1 & \thepage & Malformed strings are excluded by the local encoding convention. \\
\phantomsection\label{entry:symbol-registry-00148}\(\mathsf{TAUT}\) & \texttt{\detokenize{\mathsf{TAUT}}} & valid propositional encodings & P2.1-P7.8 & n/a & P2.1 & \thepage & Malformed strings are excluded by the fixed encoding convention; reused in P4.6 and P7.8. \\
\phantomsection\label{entry:symbol-registry-00159}\(\iota_n\) & \texttt{\detokenize{\iota_n}} & at-least-n sentence & P2.2 & n/a & P2.2 & \thepage & The first-order sentence asserting the existence of \(n\) pairwise distinct elements. \\
\phantomsection\label{entry:symbol-registry-00160}\(\Lang_R,\delta_R\) & \texttt{\detokenize{\Lang_R,}\allowbreak\detokenize{\delta_R}} & definitional-extension signature and axiom & P2.2 & n/a & P2.2 & \thepage & The signature obtained by adjoining relation symbol \(R\) and its exact defining biconditional axiom. \\
\phantomsection\label{entry:symbol-registry-00161}\(R^{\mathcal M}_{\varphi}\) & \texttt{\detokenize{R^{\mathcal M}_{\varphi}}} & relation defined by a formula in a structure & P2.2 & n/a & P2.2 & \thepage & The tuple relation cut out by \(\varphi\) in \(\mathcal M\); deliberately not \(R^*\), which is reserved for reflexive-transitive closure. \\
\phantomsection\label{entry:symbol-registry-00162}\(\mathsf{ND}_{\mathsf{FOL}},\mathsf{LK},\mathsf H_{\mathsf{FOL}}\) & \texttt{\detokenize{\mathsf{ND}_{\mathsf{FOL}},}\allowbreak\detokenize{\mathsf{LK},}\allowbreak\detokenize{\mathsf H_{\mathsf{FOL}}}} & first-order calculus identifiers & P2.2-P8.1 & n/a & P2.2 & \thepage & Names the registered natural-deduction sequent and equality-free Hilbert presentations; canonical LK recurs in P7 and the P8.1 negative-translation statements; turnstiles retain exact subscripts. \\
\phantomsection\label{entry:symbol-registry-00163}\(\Lang_{\Gamma},\Lang_{\Gamma}^{*}\) & \texttt{\detokenize{\Lang_{\Gamma},}\allowbreak\detokenize{\Lang_{\Gamma}^{*}}} & active tableau signatures & P2.2 & n/a & P2.2 & \thepage & The finite root subsignature and its countable fresh-parameter expansion used by the fair first-order tableau construction. \\
\phantomsection\label{entry:symbol-registry-00164}\(B^{+}_{\mathrm{at}},\mathcal M_{B^{+}_{\mathrm{at}}}\) & \texttt{\detokenize{B^{+}_{\mathrm{at}},}\allowbreak\detokenize{\mathcal M_{B^{+}_{\mathrm{at}}}}} & positive atomic branch and tableau countermodel & P2.2 & n/a & P2.2 & \thepage & The positive ground atoms on a fair saturated branch and the induced active-language Herbrand structure. \\
\phantomsection\label{entry:symbol-registry-00165}\(\mathcal M_B\) & \texttt{\detokenize{\mathcal M_B}} & Herbrand structure determined by ground atoms & P2.2 & n/a & P2.2 & \thepage & Ground-term structure whose positive relation atoms are exactly the members of the declared Herbrand-base subset \(B\). \\
\phantomsection\label{entry:symbol-registry-00166}\(\operatorname{nnf}(\varphi),\operatorname{nnf}^{\neg}(\varphi)\) & \texttt{\detokenize{\operatorname{nnf},}\allowbreak\detokenize{\operatorname{nnf}^{\neg}}} & first-order NNF translations & P2.2 & n/a & P2.2 & \thepage & Simultaneously recursive maps equivalent to \(\varphi\) and \(\neg\varphi\), respectively. \\
\phantomsection\label{entry:symbol-registry-00167}\(\operatorname{rect}(\varphi,V),\operatorname{pren}(\varphi)\) & \texttt{\detokenize{\operatorname{rect},}\allowbreak\detokenize{\operatorname{pren}}} & rectification and prenex maps & P2.2 & n/a & P2.2 & \thepage & Metalinguistic constructions; the forbidden-variable parameter belongs only to rectification. \\
\phantomsection\label{entry:symbol-registry-00168}\(\operatorname{Sk}(\sigma)\) & \texttt{\detokenize{\operatorname{Sk}}} & Skolem form & P2.2 & n/a & P2.2 & \thepage & Universal Skolem sentence in the expanded signature; equisatisfiability and reduct directions retain their separate choice standing. \\
\phantomsection\label{entry:symbol-registry-00169}\(\operatorname{form}_{\Lang}(C),\operatorname{ucl}_{\Lang}(C)\) & \texttt{\detokenize{\operatorname{form}_{\Lang},}\allowbreak\detokenize{\operatorname{ucl}_{\Lang}}} & clause formula and universal closure & P2.2-P12.2 & n/a & P2.2 & \thepage & Fixed literal-order disjunction for clause \(C\) and its universal closure over every occurring variable. \\
\phantomsection\label{entry:symbol-registry-00170}\(\operatorname{supp}_{\mathrm{var}}(\theta)\) & \texttt{\detokenize{\operatorname{supp}_{\mathrm{var}}}} & variable support of a term substitution & P2.2-P12.3 & n/a & P2.2 & \thepage & Variables not fixed by \(\theta\); the subscript prevents collision with finite-multiset support. \\
\phantomsection\label{entry:symbol-registry-00171}\(t\theta,\theta\eta\) & \texttt{\detokenize{t\theta,}\allowbreak\detokenize{\theta\eta}} & simultaneous substitution action and composition & P2.2-P12.3 & meta & P2.2 & \thepage & Action extends homomorphically and composition satisfies \((\theta\eta)(x)=(\theta(x))\eta\). \\
\phantomsection\label{entry:symbol-registry-00172}\(\theta\preceq_{\mathrm{inst}}\rho\) & \texttt{\detokenize{\preceq_{\mathrm{inst}}}} & substitution instance preorder & P2.2-P12.3 & meta & P2.2 & \thepage & Every instance of \(\rho\) is an instance of \(\theta\); equivalently \(\rho=\theta\eta\) for some substitution \(\eta\). \\
\phantomsection\label{entry:symbol-registry-00173}\(r\doteq t\) & \texttt{\detokenize{\doteq}} & unification equation & P2.2-P12.3 & meta & P2.2 & \thepage & Metalinguistic equation constraint processed by unification; not object-language identity. \\
\phantomsection\label{entry:symbol-registry-00174}\(H_{x,\varphi,c}\) & \texttt{\detokenize{H_{x,\varphi,c}}} & Henkin witness implication & P2.2 & n/a & P2.2 & \thepage & Metalinguistic abbreviation for \(\exists x\,\varphi\to\varphi[c/x]\) with fresh constant \(c\). \\
\phantomsection\label{entry:symbol-registry-00175}\(\Lang^H\) & \texttt{\detokenize{\Lang^H}} & Henkin constant expansion & P2.2 & n/a & P2.2 & \thepage & Signature obtained by adjoining the fresh constants used by the Henkin witness construction. \\
\phantomsection\label{entry:symbol-registry-00176}\(\sim_{\Delta},[t],\mathcal M_{\Delta},s_0\) & \texttt{\detokenize{\sim_{\Delta},}\allowbreak\detokenize{[t],}\allowbreak\detokenize{\mathcal M_{\Delta},}\allowbreak\detokenize{s_0}} & Henkin quotient-model notation & P2.2 & n/a & P2.2 & \thepage & Provable-equality relation term class quotient term structure and canonical variable assignment for a maximal consistent witness set \(\Delta\). \\
\phantomsection\label{entry:symbol-registry-00177}\(\Lang(\mathcal M),\mathcal M_{\mathcal M}\) & \texttt{\detokenize{\Lang(\mathcal M),}\allowbreak\detokenize{\mathcal M_{\mathcal M}}} & parameter signature and canonical expansion & P2.2-P3 & n/a & P2.2 & \thepage & Adds one constant for each domain element and interprets every such constant by the element it names. \\
\phantomsection\label{entry:symbol-registry-00178}\(\mathcal N_{\mathrm{ar}}\) & \texttt{\detokenize{\mathcal N_{\mathrm{ar}}}} & subscripted alias for the standard first-order arithmetic structure & P2.2-P4 & n/a & P2.2 & \thepage & Earlier separately authored material uses this explicit subscript; exact inherited Part V writes \(\mathcal N\) for the same structure in the fixed arithmetic signature. \\
\phantomsection\label{entry:symbol-registry-00179}\(\mathrm{FO}^{2},\mathrm{GF},\mathrm{EPR}\) & \texttt{\detokenize{\mathrm{FO}^{2},}\allowbreak\detokenize{\mathrm{GF},}\allowbreak\detokenize{\mathrm{EPR}}} & decidable first-order fragment identifiers & P2.2 & n/a & P2.2 & \thepage & Respectively the two-variable relational fragment guarded relational fragment and Bernays--Schoenfinkel effectively propositional class; equality and function scopes are row-local. \\
\phantomsection\label{entry:symbol-registry-00180}\(\mathsf{NEXPTIME},\mathsf{coNEXPTIME},\mathsf{EXPTIME},2\mathsf{EXPTIME}\) & \texttt{\detokenize{\mathsf{NEXPTIME},}\allowbreak\detokenize{\mathsf{coNEXPTIME},}\allowbreak\detokenize{\mathsf{EXPTIME},}\allowbreak\detokenize{2\mathsf{EXPTIME}}} & decision-problem complexity classes & P2.2-P4 & n/a & P2.2 & \thepage & Every bound applies only to the named decision problem. \\
\phantomsection\label{entry:symbol-registry-00181}\(\exists^{*}\forall^{*}\) & \texttt{\detokenize{\exists^{*}\forall^{*}}} & finite quantifier-block pattern & P2.2 & n/a & P2.2 & \thepage & A finite existential block followed by a finite universal block; stars are metalinguistic repetition rather than object-language operators. \\
\phantomsection\label{entry:symbol-registry-00184}\(\Lang_2,\mathsf{RelVar}_n\) & \texttt{\detokenize{\Lang_2,}\allowbreak\detokenize{\mathsf{RelVar}_n}} & relational second-order signature and variable stocks & P2.3 & n/a & P2.3 & \thepage & For every positive n the second component is the countably infinite stock of n-ary relation variables. \\
\phantomsection\label{entry:symbol-registry-00185}\(X^n,Y^n,U^n,F^{n+1}\) & \texttt{\detokenize{X^n,}\allowbreak\detokenize{Y^n,}\allowbreak\detokenize{U^n,}\allowbreak\detokenize{F^{n+1}}} & relation-variable metavariables & P2.3 & n/a & P2.3 & \thepage & Superscripts record arity; Y and F carry freshness conditions and U is an admissible relation-variable or predicate-constant instance. \\
\phantomsection\label{entry:symbol-registry-00186}\(P^n\) & \texttt{\detokenize{P^n}} & predicate-constant metavariable & P2.3 & n/a & P2.3 & \thepage & An n-ary nonlogical predicate constant in the active relational second-order signature. \\
\phantomsection\label{entry:symbol-registry-00187}\(\varphi[U^n/X^n]\) & \texttt{\detokenize{[U^n/}\allowbreak\detokenize{X^n]}} & capture-avoiding relation-variable substitution & P2.3 & meta & P2.3 & \thepage & Metalinguistic substitution of a relation variable or predicate constant of the same arity; distinct from term substitution. \\
\phantomsection\label{entry:symbol-registry-00188}\(\mathsf{PA}_2\) & \texttt{\detokenize{\mathsf{PA}_2}} & full second-order Peano theory & P2.3 & n/a & P2.3 & \thepage & The finite source theory displayed in Section sec:second-order-logic-03; the expanded signature includes the four displayed recursion axioms. \\
\phantomsection\label{entry:symbol-registry-00189}\(\mathsf D\) & \texttt{\detokenize{\mathsf D}} & effective calculus metavariable & P2.3 & n/a & P2.3 & \thepage & A calculus whose finite derivations are computably checkable; local to the non-enumerability theorem. \\
\phantomsection\label{entry:symbol-registry-00190}\(\mathcal H,(\mathcal R_n)_{n\geq1}\) & \texttt{\detokenize{\mathcal H,}\allowbreak\detokenize{\mathcal R_n}} & relational Henkin structure and relation domains & P2.3 & n/a & P2.3 & \thepage & The object domain is M and each nonempty relation domain is a selected subset of the corresponding powerset. \\
\phantomsection\label{entry:symbol-registry-00191}\(\mathsf H_{\mathsf{CC}}\) & \texttt{\detokenize{\mathsf H_{\mathsf{CC}}}} & matched comprehension-choice Henkin model class & P2.3 & n/a & P2.3 & \thepage & Models satisfy exactly the eligible comprehension and graph-choice schemas; distinct from arbitrary Henkin structures. \\
\phantomsection\label{entry:symbol-registry-00192}\(\mathsf{Comp}_{\chi}^{n},\mathsf{Choice}_{\chi}^{n}\) & \texttt{\detokenize{\mathsf{Comp}_{\chi}^{n},}\allowbreak\detokenize{\mathsf{Choice}_{\chi}^{n}}} & matched second-order schema abbreviations & P2.3 & n/a & P2.3 & \thepage & Metalinguistic names for the universally closed eligible formulas displayed on sheet:sol. \\
\phantomsection\label{entry:symbol-registry-00193}\(\Gamma\models_{\mathrm H}^{\mathsf{CC}}\varphi\) & \texttt{\detokenize{\entails_{\mathrm H}^{\mathsf{CC}}}} & matched Henkin consequence & P2.3 & meta & P2.3 & \thepage & Truth preservation over exactly the matched comprehension-choice Henkin model class and matching assignments. \\
\phantomsection\label{entry:symbol-registry-00194}\(\operatorname{ucl}(\varphi)\) & \texttt{\detokenize{\operatorname{ucl}}} & formula-level universal closure & P2.3 & n/a & P2.3 & \thepage & Universally closes exactly the displayed free object and relation parameters; distinct from the clause-language subscripted closure. \\
\phantomsection\label{entry:symbol-registry-00195}\(D,S_n,E_n,\widehat P\) & \texttt{\detokenize{D,}\allowbreak\detokenize{S_n,}\allowbreak\detokenize{E_n,}\allowbreak\detokenize{\widehat P}} & incidence-translation signature data & P2.3 & n/a & P2.3 & \thepage & \(D\) is the object sort; \(S_n\) the n-ary relation sort; \(E_n\) incidence; and \(\widehat P\) the relation-sort constant representing predicate \(P\). \\
\phantomsection\label{entry:symbol-registry-00196}\((\cdot)^\sharp\) & \texttt{\detokenize{{}^\sharp}} & incidence translation on formulas contexts structures and assignments & P2.3 & n/a & P2.3 & \thepage & A metalinguistic tagging and translation operation; the operand determines its eliminable clause. \\
\phantomsection\label{entry:symbol-registry-00197}\(T_{\mathsf{CC}}^\sharp\) & \texttt{\detokenize{T_{\mathsf{CC}}^\sharp}} & translated matched-schema theory & P2.3 & n/a & P2.3 & \thepage & The many-sorted first-order translations of exactly all eligible comprehension and graph-choice instances. \\
\phantomsection\label{entry:symbol-registry-00198}\(\mathsf{MSFOL},\mathsf{SOL}_{\mathrm H}^{\mathsf{CC}}\) & \texttt{\detokenize{\mathsf{MSFOL},}\allowbreak\detokenize{\mathsf{SOL}_{\mathrm H}^{\mathsf{CC}}}} & translation-defined calculus identifiers & P2.3 & n/a & P2.3 & \thepage & Eliminable metalinguistic derivability abbreviations; neither names an independently generated native rule set. \\
\phantomsection\label{entry:symbol-registry-00199}\(\operatorname{Reach}(a,b)\) & \texttt{\detokenize{\operatorname{Reach}}} & second-order reachability formula & P2.3 & n/a & P2.3 & \thepage & The full-semantic closure formula for finite directed reachability; its Henkin converse can fail. \\
\phantomsection\label{entry:symbol-registry-00200}\(\alpha,\beta,\gamma\) & \texttt{\detokenize{\alpha,}\allowbreak\detokenize{\beta,}\allowbreak\detokenize{\gamma}} & type metavariables & P2.3 & n/a & P2.3 & \thepage & Local to the simple-type grammar; disjoint from ordinal metavariables and the separate P7.4 type scope. \\
\phantomsection\label{entry:symbol-registry-00201}\(\iota,o\) & \texttt{\detokenize{\iota,}\allowbreak\detokenize{o}} & individual and truth-value base types & P2.3 & n/a & P2.3 & \thepage & Base types of the representative simple type theory; iota here is a type and not the P2.4 description binder. \\
\phantomsection\label{entry:symbol-registry-00202}\(\alpha\to_{\mathrm{ty}}\beta\) & \texttt{\detokenize{\to_{\mathrm{ty}}}} & object-language function-type former & P2.3-P7.6 & 6 & P2.3 & \thepage & Right-associative type constructor in the active typed calculus; distinct from object-language implication. \\
\phantomsection\label{entry:symbol-registry-00203}\(A,B,C,F,P,R\) & \texttt{\detokenize{A,}\allowbreak\detokenize{B,}\allowbreak\detokenize{C,}\allowbreak\detokenize{F,}\allowbreak\detokenize{P,}\allowbreak\detokenize{R}} & typed-term metavariables & P2.3 & n/a & P2.3 & \thepage & Used only under local type declarations in P2.3; disjoint from P8 formula metavariables and algebraic carriers. \\
\phantomsection\label{entry:symbol-registry-00204}\(A:\alpha\) & \texttt{\detokenize{:}} & type ascription & P2.3-P7.6 & meta & P2.3 & \thepage & The active typed calculus fixes the displayed term and type categories; distinct from the P7.7 labelled-formula colon. \\
\phantomsection\label{entry:symbol-registry-00205}\(\lambda x{:}\alpha.B\) & \texttt{\detokenize{\lambda}} & typed lambda abstraction & P2.3-P7.6 & 3 & P2.3 & \thepage & Object-language term former binding x in the displayed body; used by STLC System F and System T with capture-avoiding substitution. \\
\phantomsection\label{entry:symbol-registry-00206}\(\mathsf{Eq}_{\alpha},\forall_{\alpha}\) & \texttt{\detokenize{\mathsf{Eq}_{\alpha},}\allowbreak\detokenize{\forall_{\alpha}}} & typed equality and universal-quantifier constants & P2.3 & 3 & P2.3 & \thepage & Object-language typed constants; Eq-alpha is deliberately distinct from global metalinguistic alpha-equivalence. \\
\phantomsection\label{entry:symbol-registry-00207}\(D_\alpha\) & \texttt{\detokenize{D_\alpha}} & semantic domain at type alpha & P2.3 & n/a & P2.3 & \thepage & Nonempty type domain; standard function types contain all functions and Henkin domains may select a closure-stable subset. \\
\phantomsection\label{entry:symbol-registry-00208}\(\lbrack\!\lbrack A\rbrack\!\rbrack_s^{\mathcal M}\) & \texttt{\detokenize{\lbrack\!\lbrack\cdot\rbrack\!\rbrack_s^{\mathcal M}}} & typed-term denotation & P2.3 & n/a & P2.3 & \thepage & Recursive semantic value of a well-typed term under typed assignment s. \\
\phantomsection\label{entry:symbol-registry-00209}\(B[A/x]\) & \texttt{\detokenize{[A/}\allowbreak\detokenize{x]}} & capture-avoiding typed substitution & P2.3 & meta & P2.3 & \thepage & Substitutes the typed term A for free occurrences of x in B after the required alpha-renaming. \\
\phantomsection\label{entry:symbol-registry-00210}\(\mathcal A\) & \texttt{\detokenize{\mathcal A}} & auxiliary translated structure & P2.3-P2.4 & n/a & P2.3 & \thepage & A one-sorted or strict many-sorted model in the translation proofs; scope is disjoint from P8 algebraic structures. \\
\phantomsection\label{entry:symbol-registry-00211}\(\operatorname{ext}_n(r),s^{\mathrm{ext}}\) & \texttt{\detokenize{\operatorname{ext}_n,}\allowbreak\detokenize{s^{\mathrm{ext}}}} & extensional-collapse map and induced assignment & P2.3 & n/a & P2.3 & \thepage & Maps a relation-sort object to its n-ary extension and transports relation-variable assignments through that map. \\
\phantomsection\label{entry:symbol-registry-00212}\(\Lang_{\mathrm{MS}},\mathsf{Sort},\operatorname{tp}\) & \texttt{\detokenize{\Lang_{\mathrm{MS}},}\allowbreak\detokenize{\mathsf{Sort},}\allowbreak\detokenize{\operatorname{tp}}} & strict many-sorted signature data & P2.4 & n/a & P2.4 & \thepage & The nonempty sort set and symbol-typing map determine the strict many-sorted language. \\
\phantomsection\label{entry:symbol-registry-00213}\(\mathsf s,\mathsf t,\mathsf s_i;x^{\mathsf s},M_{\mathsf s}\) & \texttt{\detokenize{\mathsf s,}\allowbreak\detokenize{\mathsf t,}\allowbreak\detokenize{\mathsf s_i;x^{\mathsf s},}\allowbreak\detokenize{M_{\mathsf s}}} & sort metavariables sorted variables and carriers & P2.4 & n/a & P2.4 & \thepage & Sans-serif sort letters avoid the canonical assignment s; strict carriers are nonempty and pairwise disjoint. \\
\phantomsection\label{entry:symbol-registry-00214}\(\Lang_{\mathrm{MS}}^{\flat},Q_{\mathsf s}\) & \texttt{\detokenize{\Lang_{\mathrm{MS}}^{\flat},}\allowbreak\detokenize{Q_{\mathsf s}}} & one-sorted reduct signature and sort predicates & P2.4 & n/a & P2.4 & \thepage & The flat signature forgets type declarations and \(Q_{\mathsf s}\) cuts out the component for sort \(\mathsf s\). \\
\phantomsection\label{entry:symbol-registry-00215}\(\operatorname{tr}_{\mathrm{MS}},\mathsf{Ax}_{\mathrm{sort}}\) & \texttt{\detokenize{\operatorname{tr}_{\mathrm{MS}},}\allowbreak\detokenize{\mathsf{Ax}_{\mathrm{sort}}}} & many-sorted guard translation and sort theory & P2.3-P2.4 & meta & P2.3 & \thepage & The recursive translation guards bound variables; the theory asserts sort nonemptiness typing closure and disjointness. \\
\phantomsection\label{entry:symbol-registry-00216}\(\mathsf{VarGuard}_{\mathrm{MS}}(\Lang_{\mathrm{MS}})\) & \texttt{\detokenize{\mathsf{VarGuard}_{\mathrm{MS}}}} & total-assignment sort-guard set & P2.3-P2.4 & meta & P2.3 & \thepage & Metalinguistic one-sorted atoms constraining every object and relation variable to its declared sort; nonredundant with bound-variable guards and sort axioms, and choice-free for total assignments. \\
\phantomsection\label{entry:symbol-registry-00217}\(D,D^{\mathrm{in}},D^{\mathrm{out}}\) & \texttt{\detokenize{D,}\allowbreak\detokenize{D^{\mathrm{in}},}\allowbreak\detokenize{D^{\mathrm{out}}}} & free-logic outer and inner domains & P2.4-P8.4 & n/a & P2.4 & \thepage & Local outer-domain presentations; the inner existing domain is nonempty and contained in the outer domain. \\
\phantomsection\label{entry:symbol-registry-00218}\(E!(t)\) & \texttt{\detokenize{E!}} & free-logic existence predicate & P2.4-P8.4 & 3 & P2.4 & \thepage & True exactly when the term value lies in the declared inner existing domain of the displayed free-logic semantics. \\
\phantomsection\label{entry:symbol-registry-00219}\(\iota x\,\varphi\) & \texttt{\detokenize{\iota}} & definite-description term former & P2.4-P8.4 & 3 & P2.4 & \thepage & Binds x in the description formula; a local Russellian elimination or primitive partial-term policy fixes its standing. \\
\phantomsection\label{entry:symbol-registry-00220}\(L_{\omega_1\omega}\) & \texttt{\detokenize{L_{\omega_1\omega}}} & countably infinitary first-order language & P2.4 & n/a & P2.4 & \thepage & Allows countable Boolean combinations but only finite quantifier strings. \\
\phantomsection\label{entry:symbol-registry-00221}\(\bigwedge_{i\in I}\varphi_i\) & \texttt{\detokenize{\bigwedge}} & countable infinitary conjunction & P2.4 & 4 & P2.4 & \thepage & I is countable; the connective retains the ordinary conjunction precedence rank. \\
\phantomsection\label{entry:symbol-registry-00222}\(\bigvee_{i\in I}\varphi_i\) & \texttt{\detokenize{\bigvee}} & countable infinitary disjunction & P2.4 & 5 & P2.4 & \thepage & I is countable; the connective retains the ordinary disjunction precedence rank. \\
\phantomsection\label{entry:symbol-registry-00223}\(\mathsf{FOL}(Q)\) & \texttt{\detokenize{\mathsf{FOL}(Q)}} & generalized-quantifier extension placeholder & P2.4 & n/a & P2.4 & \thepage & Pointer notation only; Part 3 must define the supplied isomorphism-invariant quantifier family and its exact arity. \\
\phantomsection\label{entry:symbol-registry-00224}\(\Sigma^1_n,\Pi^1_n,\Delta^1_n,\boldsymbol\Sigma^1_n,\boldsymbol\Pi^1_n\) & \texttt{\detokenize{\Sigma^1_n,}\allowbreak\detokenize{\Pi^1_n,}\allowbreak\detokenize{\Delta^1_n,}\allowbreak\detokenize{\boldsymbol\Sigma^1_n,}\allowbreak\detokenize{\boldsymbol\Pi^1_n}} & lightface and boldface analytical hierarchy classes & P2.4-P4.5 & n/a & P2.4 & \thepage & Lightface classes have effective matrices and \(\Delta^1_n\) is their intersection; boldface classes allow real or set parameters and must not be identified with them. \\
\end{longtable}
\endgroup
\Needspace{.4\textheight}
\subsection*{Part 3 scopes}
\begingroup\scriptsize\setlength{\tabcolsep}{3pt}
\begin{longtable}{@{}p{0.9in}p{1.3in}p{1.0in}p{0.5in}p{0.3in}p{0.4in}p{0.35in}p{1.65in}@{}}
\toprule
\textbf{Glyph} & \textbf{Macro} & \textbf{Role} & \textbf{Scope} & \textbf{Prec} & \textbf{First} & \textbf{Reg. page} & \textbf{Principal uses / notes} \\
\midrule
\endfirsthead
\toprule
\textbf{Glyph} & \textbf{Macro} & \textbf{Role} & \textbf{Scope} & \textbf{Prec} & \textbf{First} & \textbf{Reg. page} & \textbf{Principal uses / notes} \\
\midrule
\endhead
\midrule\multicolumn{8}{r}{\emph{continued}}\\\endfoot
\bottomrule\endlastfoot
\phantomsection\label{entry:symbol-registry-00044}\(\operatorname{Diag}(\mathcal M)\) & \texttt{\detokenize{\operatorname{Diag}}} & atomic diagram & P3 & n/a & P3.1 & \thepage & Atomic sentences with parameters true in \(\mathcal M\). \\
\phantomsection\label{entry:symbol-registry-00046}\(S_n(T)\) & \texttt{\detokenize{S_n(T)}} & type space & P3 & n/a & P3.2 & \thepage & Complete \(n\)-types consistent with \(T\) under the clopen topology. \\
\phantomsection\label{entry:symbol-registry-00047}\(\prod_{i\in I}\mathcal M_i/\mathcal U\) & \texttt{\detokenize{\prod_{i\in I}\mathcal M_i/}\allowbreak\detokenize{\mathcal U}} & ultraproduct & P3.3 & n/a & P3.3 & \thepage & Quotient of the direct product by equality on a set in the ultrafilter \(\mathcal U\). \\
\phantomsection\label{entry:symbol-registry-00439}\(\PA\) & \texttt{\detokenize{\PA}} & Peano arithmetic & P3-P7.5 & n/a & P3.3 & \thepage & The fixed first-order theory defined in P5.1; earlier table and later metatheoretic uses denote that same theory. \\
\end{longtable}
\endgroup
\Needspace{.4\textheight}
\subsection*{Part 4 scopes}
\begingroup\scriptsize\setlength{\tabcolsep}{3pt}
\begin{longtable}{@{}p{0.9in}p{1.3in}p{1.0in}p{0.5in}p{0.3in}p{0.4in}p{0.35in}p{1.65in}@{}}
\toprule
\textbf{Glyph} & \textbf{Macro} & \textbf{Role} & \textbf{Scope} & \textbf{Prec} & \textbf{First} & \textbf{Reg. page} & \textbf{Principal uses / notes} \\
\midrule
\endfirsthead
\toprule
\textbf{Glyph} & \textbf{Macro} & \textbf{Role} & \textbf{Scope} & \textbf{Prec} & \textbf{First} & \textbf{Reg. page} & \textbf{Principal uses / notes} \\
\midrule
\endhead
\midrule\multicolumn{8}{r}{\emph{continued}}\\\endfoot
\bottomrule\endlastfoot
\phantomsection\label{entry:symbol-registry-00015}\(\mathcal A_{\mathrm{in}},\mathcal A_{\mathrm{tape}}\) & \texttt{\detokenize{\mathcal A_{\mathrm{in}},}\allowbreak\detokenize{\mathcal A_{\mathrm{tape}}}} & input and tape alphabets & P4.1 & n/a & P4.1 & \thepage & No formula-context standing in this scope. \\
\phantomsection\label{entry:symbol-registry-00016}\(u,v,w\) & \texttt{\detokenize{u,}\allowbreak\detokenize{v,}\allowbreak\detokenize{w}} & finite words & P4.1 & n/a & P4.1 & \thepage & Finite words over the declared alphabet; later P4 chapters use the letters numerically. \\
\phantomsection\label{entry:symbol-registry-00017}\(\Sigma^0_n,\Pi^0_n\) & \texttt{\detokenize{\Sigma^0_n,}\allowbreak\detokenize{\Pi^0_n}} & syntactic arithmetic formula classes & P4-P7.5 & n/a & P4.3 & \thepage & Classes fixed by the declared normalization and unbounded-quantifier alternation; P7.5 uses the same hierarchy in induction-fragment names. \\
\phantomsection\label{entry:symbol-registry-00018}\(\boldsymbol{\Sigma}^0_n,\boldsymbol{\Pi}^0_n,\boldsymbol{\Delta}^0_n\) & \texttt{\detokenize{\boldsymbol{\Sigma}^0_n,}\allowbreak\detokenize{\boldsymbol{\Pi}^0_n,}\allowbreak\detokenize{\boldsymbol{\Delta}^0_n}} & arithmetical definability classes of relations & P4-P5 & n/a & P4.3 & \thepage & Bold symbols classify relations definable in \(\mathcal N_{\mathrm{ar}}\), not formula strings or Borel pointclasses. \\
\phantomsection\label{entry:symbol-registry-00049}\(\simeq\) & \texttt{\detokenize{\simeq}} & Kleene equality & P4.2-P4.5 & meta & P4.2 & \thepage & Both sides are defined with equal value or both are undefined. \\
\phantomsection\label{entry:symbol-registry-00050}\(\leq_m\) & \texttt{\detokenize{\leq_m}} & many-one reducibility & P4-P5 & meta & P4.2 & \thepage & A total computable translator preserves membership. \\
\phantomsection\label{entry:symbol-registry-00051}\(\leq_T\) & \texttt{\detokenize{\leq_T}} & Turing reducibility & P4-P5 & meta & P4.4 & \thepage & An oracle machine for the right-hand set decides the left-hand set. \\
\phantomsection\label{entry:symbol-registry-00052}\(\varphi_e\) & \texttt{\detokenize{\varphi_e}} & partial computable function with index e & P4-P5 & n/a & P4.2 & \thepage & Oracle form \(\varphi^B_e\) is relative to \(B\); finite-string form \(\varphi_e^\tau\) and stage form \(\varphi^{B_s}_{e,s}\) are defined in the two priority proofs. \\
\phantomsection\label{entry:symbol-registry-00053}\(B'\) & \texttt{\detokenize{'}} & Turing jump & P4.4 & n/a & P4.4 & \thepage & Postfix prime on an oracle set. \\
\phantomsection\label{entry:symbol-registry-00054}\(\varnothing^{(n)}\) & \texttt{\detokenize{\varnothing^{(n)}}} & iterated jump of the empty set & P4.4-P5 & n/a & P4.4 & \thepage & First defined with the general jump iteration. \\
\phantomsection\label{entry:symbol-registry-00056}\(\bar n\) & \texttt{\detokenize{\bar n}} & numeral & P4.3-P5 & n/a & P4.3 & \thepage & Closed arithmetic term \(S^n0\); bars are reserved here for numerals, while tuples use vectors. \\
\phantomsection\label{entry:symbol-registry-00080}\(\leq\) & \texttt{\detokenize{\leq}} & object-language arithmetic order abbreviation & P4.3 & 2 & P4.3 & \thepage & Eliminable abbreviation \(r\leq t\defeq\exists z\,(z+r=t)\), with the least variable fresh for the displayed terms; marked bounded quantifiers inherit it. \\
\phantomsection\label{entry:symbol-registry-00081}\(\leq\) & \texttt{\detokenize{\leq}} & canonical ordering & P4.4-P4.6 & 2 & P4.4 & \thepage & Ambient order on degrees, natural numbers, and resource bounds; not arithmetic object-language syntax. \\
\phantomsection\label{entry:symbol-registry-00084}\(a\leq_{\N}b\) & \texttt{\detokenize{\leq_{\N}}} & metalinguistic natural-number order & P4.3-P5.1 & 2 & P4.3 & \thepage & Used where bare \(\leq\) simultaneously denotes arithmetic-language syntax; operands are standard naturals or codes, never arithmetic terms. \\
\phantomsection\label{entry:symbol-registry-00086}\(<\) & \texttt{\detokenize{<}} & object-language arithmetic strict-order abbreviation & P4.3 & 2 & P4.3 & \thepage & Eliminable abbreviation \(r<t\defeq\exists z\,(r+Sz=t)\), with the least variable fresh for the displayed terms. \\
\phantomsection\label{entry:symbol-registry-00087}\(<\) & \texttt{\detokenize{<}} & canonical strict ordering & P4.4-P14 & 2 & P4.4 & \thepage & Ambient strict order on declared numerical and order-theoretic objects; not arithmetic object-language syntax. \\
\phantomsection\label{entry:symbol-registry-00088}\(a<_{\N}b\) & \texttt{\detokenize{<_{\N}}} & metalinguistic natural-number strict order & P4.3 & 2 & P4.3 & \thepage & Used where bare \(<\) simultaneously denotes arithmetic-language syntax; operands are standard naturals or codes, never arithmetic terms. \\
\phantomsection\label{entry:symbol-registry-00227}\(\vec a,\vec n\) & \texttt{\detokenize{\vec a,}\allowbreak\detokenize{\vec n}} & finite tuples of natural-number parameters & P4-P5 & n/a & P4.1 & \thepage & Vector arrows distinguish tuples from numeral bars. \\
\phantomsection\label{entry:symbol-registry-00228}\(\vec x,\vec y,\vec z,\vec u,\vec v\) & \texttt{\detokenize{\vec x,}\allowbreak\detokenize{\vec y,}\allowbreak\detokenize{\vec z,}\allowbreak\detokenize{\vec u,}\allowbreak\detokenize{\vec v}} & finite tuples of numerical variables & P4-P5 & n/a & P4.1 & \thepage & Vector arrows distinguish variable blocks from numeral bars. \\
\phantomsection\label{entry:symbol-registry-00229}\(\overline{\vec a}\) & \texttt{\detokenize{\overline{\vec a}}} & tuple of arithmetic numerals & P4.3-P5 & n/a & P4.3 & \thepage & Componentwise numeral formation from a declared tuple of standard naturals. \\
\phantomsection\label{entry:symbol-registry-00230}\(\mathcal T\) & \texttt{\detokenize{\mathcal T}} & deterministic Turing-machine metavariable & P4.1-P4.2 & n/a & P4.1 & \thepage & Machine under the fixed single-tape convention. \\
\phantomsection\label{entry:symbol-registry-00231}\(\sqcup\) & \texttt{\detokenize{\sqcup}} & blank tape symbol & P4.1-P4.2 & n/a & P4.1 & \thepage & Distinguished tape-alphabet symbol. \\
\phantomsection\label{entry:symbol-registry-00232}\(\delta\) & \texttt{\detokenize{\delta}} & Turing-machine transition function & P4.1-P4.2 & n/a & P4.1 & \thepage & Typed by the fixed machine tuple. \\
\phantomsection\label{entry:symbol-registry-00233}\(\longrightarrow_{\mathcal T}\) & \texttt{\detokenize{\longrightarrow_{\mathcal T}}} & one-step machine transition & P4.1-P4.2 & meta & P4.1 & \thepage & Relation on configurations of the displayed machine. \\
\phantomsection\label{entry:symbol-registry-00234}\(\operatorname{bin}(n)\) & \texttt{\detokenize{\operatorname{bin}}} & canonical binary word & P4.1-P4.2 & n/a & P4.1 & \thepage & Zero is represented by the one-symbol word zero. \\
\phantomsection\label{entry:symbol-registry-00235}\(\operatorname{bin}_m(i)\) & \texttt{\detokenize{\operatorname{bin}_m}} & fixed-width binary address & P4.6 & n/a & P4.6 & \thepage & Length-m expansion padded with leading zeros for succinct-instance bit addressing. \\
\phantomsection\label{entry:symbol-registry-00236}\(\operatorname{pair}_{\N}(a,b)\) & \texttt{\detokenize{\operatorname{pair}_{\N}}} & numerical pairing function & P4.1-P4.3 & n/a & P4.1 & \thepage & Distinct from the set-theoretic ordered pair. \\
\phantomsection\label{entry:symbol-registry-00237}\(\operatorname{code}_k(\vec n)\) & \texttt{\detokenize{\operatorname{code}_k}} & fixed-arity tuple code & P4.1-P4.3 & n/a & P4.1 & \thepage & Defined recursively from the numerical pairing function. \\
\phantomsection\label{entry:symbol-registry-00238}\(\chi_A\) & \texttt{\detokenize{\chi_A}} & characteristic function of a set & P4.1-P4.4 & n/a & P4.1 & \thepage & Source-bound notation: one exactly on members of \(A\). \\
\phantomsection\label{entry:symbol-registry-00239}\(\chi_R\) & \texttt{\detokenize{\chi_R}} & characteristic function of a relation & P4.1-P4.4 & n/a & P4.1 & \thepage & Source-bound notation: one exactly on tuples in \(R\). \\
\phantomsection\label{entry:symbol-registry-00240}\(\mu g\) & \texttt{\detokenize{\mu}} & unbounded minimization & P4.1 & n/a & P4.1 & \thepage & Definedness-sensitive least-zero operation; relational least-true use is the explicit zero-test abbreviation in P4.1. \\
\phantomsection\label{entry:symbol-registry-00241}\(\mathsf T_k(e,\vec x,y)\) & \texttt{\detokenize{\mathsf T_k}} & primitive-recursive accepting-computation-code relation & P4.1 & n/a & P4.1 & \thepage & Kleene normal-form relation. \\
\phantomsection\label{entry:symbol-registry-00242}\(\mathsf U(y)\) & \texttt{\detokenize{\mathsf U}} & primitive-recursive normal-form output extractor & P4.1 & n/a & P4.1 & \thepage & Distinct from the universal program evaluator. \\
\phantomsection\label{entry:symbol-registry-00243}\(U(e,x)\) & \texttt{\detokenize{U}} & universal partial computable function & P4.2-P4.4 & n/a & P4.2 & \thepage & Exact source notation; oracle superscript \(U^B\) gives relativization. \\
\phantomsection\label{entry:symbol-registry-00244}\(\mathsf{par}_n^m\) & \texttt{\detokenize{\mathsf{par}_n^m}} & parameter-specialization compiler & P4.2-P4.4 & n/a & P4.2 & \thepage & Total computable compiler for the fixed numbering. \\
\phantomsection\label{entry:symbol-registry-00245}\(K\) & \texttt{\detokenize{K}} & self-application domain set & P4.2-P4.4 & n/a & P4.2 & \thepage & Indices \(e\) for which \(\varphi_e(e)\) is defined. \\
\phantomsection\label{entry:symbol-registry-00246}\(H\) & \texttt{\detokenize{H}} & binary application domain set & P4.2 & n/a & P4.2 & \thepage & Paired indices and inputs with defined computation. \\
\phantomsection\label{entry:symbol-registry-00247}\(\Lang_A\) & \texttt{\detokenize{\Lang_A}} & first-order language of arithmetic & P4.3-P5 & n/a & P4.3 & \thepage & Signature containing zero successor addition and multiplication. \\
\phantomsection\label{entry:symbol-registry-00248}\(F(\vec x)\mathord\downarrow\) & \texttt{\detokenize{\downarrow}} & convergence of a partial numerical expression & P4.3-P4.5 & meta & P4.3 & \thepage & Metalinguistic definedness predicate. \\
\phantomsection\label{entry:symbol-registry-00249}\(F(\vec x)\mathord\uparrow\) & \texttt{\detokenize{\uparrow}} & divergence of a partial numerical expression & P4.3-P4.5 & meta & P4.3 & \thepage & Metalinguistic undefinedness predicate. \\
\phantomsection\label{entry:symbol-registry-00250}\(\mathsf P(x,y,z)\) & \texttt{\detokenize{\mathsf P}} & bounded arithmetic pairing-graph formula & P4.3 & n/a & P4.3 & \thepage & Scope is disjoint from complexity class \(\mathsf P\). \\
\phantomsection\label{entry:symbol-registry-00251}\(\mathsf B(w,i,v)\) & \texttt{\detokenize{\mathsf B}} & bounded Goedel-beta graph formula & P4.3 & n/a & P4.3 & \thepage & Object-language arithmetic formula. \\
\phantomsection\label{entry:symbol-registry-00252}\(\beta(w,i)\) & \texttt{\detokenize{\beta}} & Goedel beta function & P4.3 & n/a & P4.3 & \thepage & Numerical finite-sequence decoder. \\
\phantomsection\label{entry:symbol-registry-00253}\(\mathsf{Comp}(e,x,s)\) & \texttt{\detokenize{\mathsf{Comp}}} & bounded accepting-history formula & P4.3 & n/a & P4.3 & \thepage & The bounded formula supplied by the inherited verification theorem; no canonical first formula is asserted. \\
\phantomsection\label{entry:symbol-registry-00254}\(\mathsf{TOT}\) & \texttt{\detokenize{\mathsf{TOT}}} & totality index set & P4.3 & n/a & P4.3 & \thepage & Indices of programs accepting with valid output on every input. \\
\phantomsection\label{entry:symbol-registry-00255}\(\mathcal C\) & \texttt{\detokenize{\mathcal C}} & locally declared class of sets or encoded languages & P4.3-P4.6 & n/a & P4.3 & \thepage & The active class and member type are declared locally. \\
\phantomsection\label{entry:symbol-registry-00256}\(N!\) & \texttt{\detokenize{!}} & factorial & P4.3 & 1 & P4.3 & \thepage & Scope disjoint from the Part-8 exponential. \\
\phantomsection\label{entry:symbol-registry-00257}\(B^{(n)}\) & \texttt{\detokenize{B^{(n)}}} & n-fold Turing jump & P4.4-P5 & n/a & P4.4 & \thepage & Iteration from an arbitrary oracle set. \\
\phantomsection\label{entry:symbol-registry-00258}\(\equiv_m\) & \texttt{\detokenize{\equiv_m}} & many-one equivalence & P4.4 & meta & P4.4 & \thepage & Exact source notation for mutual many-one reducibility. \\
\phantomsection\label{entry:symbol-registry-00259}\(\sim_T\) & \texttt{\detokenize{\sim_T}} & Turing equivalence & P4.4 & meta & P4.4 & \thepage & Mutual Turing reducibility. \\
\phantomsection\label{entry:symbol-registry-00260}\(\deg_T(A)\) & \texttt{\detokenize{\deg_T}} & Turing degree of a set & P4.4 & n/a & P4.4 & \thepage & Equivalence class under Turing equivalence. \\
\phantomsection\label{entry:symbol-registry-00261}\([A]_T\) & \texttt{\detokenize{[A]_T}} & Turing-equivalence class & P4.4 & n/a & P4.4 & \thepage & Alternate display of the degree of \(A\). \\
\phantomsection\label{entry:symbol-registry-00262}\(A\oplus B\) & \texttt{\detokenize{\oplus}} & computable join representative & P4.4 & n/a & P4.4 & \thepage & Even-odd coded union. \\
\phantomsection\label{entry:symbol-registry-00263}\(\mathbf0\) & \texttt{\detokenize{\mathbf0}} & computable Turing degree & P4.4 & n/a & P4.4 & \thepage & Degree of the empty set. \\
\phantomsection\label{entry:symbol-registry-00264}\(\mathbf0'\) & \texttt{\detokenize{\mathbf0'}} & jump of the computable degree & P4.4 & n/a & P4.4 & \thepage & Degree of the empty-set jump. \\
\phantomsection\label{entry:symbol-registry-00265}\(R_e\) & \texttt{\detokenize{R_e}} & priority requirement against computing A from B & P4.4 & n/a & P4.4 & \thepage & Finite-injury requirement. \\
\phantomsection\label{entry:symbol-registry-00266}\(S_e\) & \texttt{\detokenize{S_e}} & priority requirement against computing B from A & P4.4 & n/a & P4.4 & \thepage & Finite-injury requirement. \\
\phantomsection\label{entry:symbol-registry-00267}\(W_e\) & \texttt{\detokenize{W_e}} & c.e. set induced by program index e & P4.4 & n/a & P4.4 & \thepage & Domain under the acceptable numbering. \\
\phantomsection\label{entry:symbol-registry-00268}\(\leq_1\) & \texttt{\detokenize{\leq_1}} & one-one reducibility & P4.4 & meta & P4.4 & \thepage & The translator is total computable and injective. \\
\phantomsection\label{entry:symbol-registry-00269}\(D_d\) & \texttt{\detokenize{D_d}} & finite set decoded from d & P4.4 & n/a & P4.4 & \thepage & Positions of one-bits in \(d\). \\
\phantomsection\label{entry:symbol-registry-00270}\(I_{\mathcal C}\) & \texttt{\detokenize{I_{\mathcal C}}} & extensional index set & P4.4 & n/a & P4.4 & \thepage & Indices \(e\) such that \(W_e\) belongs to the declared class. \\
\phantomsection\label{entry:symbol-registry-00271}\(\mathsf{Fin}_{\mathcal C}\) & \texttt{\detokenize{\mathsf{Fin}_{\mathcal C}}} & c.e. set of finite-witness codes & P4.4 & n/a & P4.4 & \thepage & Rice-Shapiro finite positive witnesses. \\
\phantomsection\label{entry:symbol-registry-00272}\(\alpha_i\) & \texttt{\detokenize{\alpha_i}} & function variable in the analytical hierarchy & P4.5 & n/a & P4.5 & \thepage & Scope disjoint from ordinal and type metavariables. \\
\phantomsection\label{entry:symbol-registry-00273}\(\mathcal O\) & \texttt{\detokenize{\mathcal O}} & Kleene system of ordinal notations & P4.5 & n/a & P4.5 & \thepage & The cited standard system fixed in the method entry. \\
\phantomsection\label{entry:symbol-registry-00274}\(<_{\mathcal O}\) & \texttt{\detokenize{<_{\mathcal O}}} & Kleene notation order & P4.5 & 2 & P4.5 & \thepage & Order belonging to the fixed cited notation system. \\
\phantomsection\label{entry:symbol-registry-00275}\(H_a\) & \texttt{\detokenize{H_a}} & effective jump-hierarchy set & P4.5 & n/a & P4.5 & \thepage & Indexed by a Kleene-O notation. \\
\phantomsection\label{entry:symbol-registry-00276}\(\mathsf{HYP}\) & \texttt{\detokenize{\mathsf{HYP}}} & hyperarithmetical subsets of the naturals & P4.5 & n/a & P4.5 & \thepage & Lightface class. \\
\phantomsection\label{entry:symbol-registry-00277}\(\mathsf{LO}\) & \texttt{\detokenize{\mathsf{LO}}} & valid computable linear-order codes & P4.5 & n/a & P4.5 & \thepage & Universe within which \(\mathsf{WF}\) and \(\mathsf{IF}\) are complementary. \\
\phantomsection\label{entry:symbol-registry-00278}\(R_e\) & \texttt{\detokenize{R_e}} & computable relation decoded from e & P4.5 & n/a & P4.5 & \thepage & Scope disjoint from the P4.4 priority-requirement row. \\
\phantomsection\label{entry:symbol-registry-00279}\(\mathsf{WF}\) & \texttt{\detokenize{\mathsf{WF}}} & codes of computable well-orders & P4.5 & n/a & P4.5 & \thepage & Subset of \(\mathsf{LO}\). \\
\phantomsection\label{entry:symbol-registry-00280}\(\mathsf{IF}\) & \texttt{\detokenize{\mathsf{IF}}} & codes of ill-founded computable linear orders & P4.5 & n/a & P4.5 & \thepage & Complement of \(\mathsf{WF}\) within \(\mathsf{LO}\). \\
\phantomsection\label{entry:symbol-registry-00281}\(\mathsf U^1_{\Sigma_n}\) & \texttt{\detokenize{\mathsf U^1_{\Sigma_n}}} & universal lightface Sigma-one-n relation & P4.5 & n/a & P4.5 & \thepage & Pairs a valid definition code with an input satisfying that definition. \\
\phantomsection\label{entry:symbol-registry-00282}\(\mathsf U^1_{\Pi_n}\) & \texttt{\detokenize{\mathsf U^1_{\Pi_n}}} & universal lightface Pi-one-n relation & P4.5 & n/a & P4.5 & \thepage & Complementary universal relation under the fixed coding. \\
\phantomsection\label{entry:symbol-registry-00283}\(\operatorname{len}(x)\) & \texttt{\detokenize{\operatorname{len}}} & bit length of a finite binary string & P4.5-P6.2 & n/a & P4.5 & \thepage & Replaces forbidden vertical-bar length notation. \\
\phantomsection\label{entry:symbol-registry-00284}\(C_V(x)\) & \texttt{\detokenize{C_V}} & plain Kolmogorov complexity & P4.5 & n/a & P4.5 & \thepage & Relative to decompressor \(V\). \\
\phantomsection\label{entry:symbol-registry-00285}\(K_V(x)\) & \texttt{\detokenize{K_V}} & prefix-free Kolmogorov complexity & P4.5 & n/a & P4.5 & \thepage & Relative to prefix-free decompressor \(V\). \\
\phantomsection\label{entry:symbol-registry-00286}\(U_C\) & \texttt{\detokenize{U_C}} & optimal plain decompressor & P4.5 & n/a & P4.5 & \thepage & Distinguished from theory metavariables. \\
\phantomsection\label{entry:symbol-registry-00287}\(U_K\) & \texttt{\detokenize{U_K}} & optimal prefix-free machine & P4.5 & n/a & P4.5 & \thepage & Used throughout the Chaitin entries instead of bare \(U\). \\
\phantomsection\label{entry:symbol-registry-00288}\(\Omega_{U_K}\) & \texttt{\detokenize{\Omega_{U_K}}} & halting probability of the fixed optimal prefix-free machine & P4.5 & n/a & P4.5 & \thepage & Not attached to an arbitrary plain machine. \\
\phantomsection\label{entry:symbol-registry-00289}\(\mathsf{Run}_{U_K}(p,x,t)\) & \texttt{\detokenize{\mathsf{Run}_{U_K}}} & bounded prefix-machine computation formula & P4.5 & n/a & P4.5 & \thepage & Fixed-machine/string analogue of the inherited Comp formula; the finite trace codings are primitively recursively intertranslatable. \\
\phantomsection\label{entry:symbol-registry-00290}\(\operatorname{BB}_{\mathrm{step}}(n)\) & \texttt{\detokenize{\operatorname{BB}_{\mathrm{step}}}} & Busy Beaver step function & P4.5 & n/a & P4.5 & \thepage & Maximum halting transition count under the displayed machine convention. \\
\phantomsection\label{entry:symbol-registry-00291}\(\operatorname{BB}_{1}(n)\) & \texttt{\detokenize{\operatorname{BB}_{1}}} & Busy Beaver one-count function & P4.5 & n/a & P4.5 & \thepage & Maximum number of one-filled cells at halt under the displayed convention. \\
\phantomsection\label{entry:symbol-registry-00292}\(A\leq_m^pB\) & \texttt{\detokenize{\leq_m^p}} & polynomial-time many-one reducibility & P4.6 & meta & P4.6 & \thepage & A total deterministic polynomial-time translator preserves membership. \\
\phantomsection\label{entry:symbol-registry-00293}\(\mathsf P\) & \texttt{\detokenize{\mathsf P}} & deterministic polynomial-time language class & P4.6 & n/a & P4.6 & \thepage & Scope disjoint from the P4.3 pairing formula. \\
\phantomsection\label{entry:symbol-registry-00294}\(\mathsf{SAT}\) & \texttt{\detokenize{\mathsf{SAT}}} & satisfiable propositional encodings & P4.6 & n/a & P4.6 & \thepage & Cook-Levin target under the declared encoding. \\
\phantomsection\label{entry:symbol-registry-00295}\(3\mathsf{SAT}\) & \texttt{\detokenize{3\mathsf{SAT}}} & 3-CNF satisfiability language & P4.6 & n/a & P4.6 & \thepage & Encodings of satisfiable 3-CNF formulas whose clauses contain at most three literals. \\
\phantomsection\label{entry:symbol-registry-00296}\(\mathsf{DTIME}(f(n))\) & \texttt{\detokenize{\mathsf{DTIME}}} & deterministic time-class constructor & P4.6 & n/a & P4.6 & \thepage & Under the fixed machine convention. \\
\phantomsection\label{entry:symbol-registry-00297}\(\mathsf{NTIME}(f(n))\) & \texttt{\detokenize{\mathsf{NTIME}}} & nondeterministic time-class constructor & P4.6 & n/a & P4.6 & \thepage & Under the fixed machine convention. \\
\phantomsection\label{entry:symbol-registry-00298}\(\mathsf{SPACE}(f(n))\) & \texttt{\detokenize{\mathsf{SPACE}}} & deterministic space-class constructor & P4.6 & n/a & P4.6 & \thepage & Counts work-tape cells. \\
\phantomsection\label{entry:symbol-registry-00299}\(\mathsf{NSPACE}(f(n))\) & \texttt{\detokenize{\mathsf{NSPACE}}} & nondeterministic space-class constructor & P4.6 & n/a & P4.6 & \thepage & Counts work-tape cells. \\
\phantomsection\label{entry:symbol-registry-00300}\(\mathsf{PSPACE}\) & \texttt{\detokenize{\mathsf{PSPACE}}} & polynomial-space language class & P4.6 & n/a & P4.6 & \thepage & Decision-language class. \\
\phantomsection\label{entry:symbol-registry-00301}\(\mathsf{NPSPACE}\) & \texttt{\detokenize{\mathsf{NPSPACE}}} & nondeterministic polynomial-space language class & P4.6 & n/a & P4.6 & \thepage & Decision-language class. \\
\phantomsection\label{entry:symbol-registry-00302}\(\mathsf{EXPSPACE}\) & \texttt{\detokenize{\mathsf{EXPSPACE}}} & exponential-space language class & P4.6 & n/a & P4.6 & \thepage & Decision-language class. \\
\phantomsection\label{entry:symbol-registry-00303}\(O(f(n))\) & \texttt{\detokenize{O}} & asymptotic upper bound & P4.6 & n/a & P4.6 & \thepage & Resource-bound asymptotics. \\
\phantomsection\label{entry:symbol-registry-00304}\(o(g(n))\) & \texttt{\detokenize{o}} & strict asymptotic domination & P4.6 & n/a & P4.6 & \thepage & Hierarchy-gap hypothesis. \\
\phantomsection\label{entry:symbol-registry-00305}\(\Sigma^p_k\) & \texttt{\detokenize{\Sigma^p_k}} & existential-start polynomial-hierarchy class & P4.6 & n/a & P4.6 & \thepage & Bounded alternating string quantifiers. \\
\phantomsection\label{entry:symbol-registry-00306}\(\Pi^p_k\) & \texttt{\detokenize{\Pi^p_k}} & universal-start polynomial-hierarchy class & P4.6 & n/a & P4.6 & \thepage & Bounded alternating string quantifiers. \\
\phantomsection\label{entry:symbol-registry-00307}\(\mathsf{PH}\) & \texttt{\detokenize{\mathsf{PH}}} & polynomial hierarchy & P4.6 & n/a & P4.6 & \thepage & Union of its \(\Sigma^p_k\) levels. \\
\phantomsection\label{entry:symbol-registry-00308}\(\mathsf{TQBF}\) & \texttt{\detokenize{\mathsf{TQBF}}} & true closed quantified-Boolean-formula language & P4.6 & n/a & P4.6 & \thepage & Unbounded alternation. \\
\phantomsection\label{entry:symbol-registry-00309}\(\Sigma_k\text{-}\mathsf{SAT}\) & \texttt{\detokenize{\Sigma_k\text{-}\mathsf{SAT}}} & existential-start fixed-alternation QBF language & P4.6 & n/a & P4.6 & \thepage & Fixed \(k\). \\
\phantomsection\label{entry:symbol-registry-00310}\(\Pi_k\text{-}\mathsf{SAT}\) & \texttt{\detokenize{\Pi_k\text{-}\mathsf{SAT}}} & universal-start fixed-alternation QBF language & P4.6 & n/a & P4.6 & \thepage & Fixed \(k\). \\
\phantomsection\label{entry:symbol-registry-00311}\(\mathfrak L\) & \texttt{\detokenize{\mathfrak L}} & logic metavariable in a capture statement & P4.6 & n/a & P4.6 & \thepage & Distinct from \(\Lang\), which globally denotes a signature. \\
\phantomsection\label{entry:symbol-registry-00312}\(A\leq_m^{\log}B\) & \texttt{\detokenize{\leq_m^{\log}}} & logspace many-one reducibility & P4.6 & meta & P4.6 & \thepage & A deterministic logarithmic-work-space translator preserves membership. \\
\phantomsection\label{entry:symbol-registry-00313}\(\mathsf{CVP}\) & \texttt{\detokenize{\mathsf{CVP}}} & circuit value problem & P4.6 & n/a & P4.6 & \thepage & Valid acyclic Boolean circuit encodings with output value one. \\
\phantomsection\label{entry:symbol-registry-00314}\(\mathsf{SCVP}\) & \texttt{\detokenize{\mathsf{SCVP}}} & succinct circuit value problem & P4.6 & n/a & P4.6 & \thepage & Circuit-described exponentially large circuit value instances. \\
\phantomsection\label{entry:symbol-registry-00315}\(\mathsf{STQBF}\) & \texttt{\detokenize{\mathsf{STQBF}}} & succinct true quantified Boolean formula problem & P4.6 & n/a & P4.6 & \thepage & Circuit-described exponentially long closed QBF instances. \\
\phantomsection\label{entry:symbol-registry-00316}\([\operatorname{lfp}_{R,\bar x}\varphi](\bar t)\) & \texttt{\detokenize{\operatorname{lfp}}} & least fixed-point formula constructor & P4.6-P9 & n/a & P4.6 & \thepage & Requires positive occurrence of the relation variable and denotes the least monotone fixed point. \\
\end{longtable}
\endgroup
\Needspace{.4\textheight}
\subsection*{Part 5 scopes}
\begingroup\scriptsize\setlength{\tabcolsep}{3pt}
\begin{longtable}{@{}p{0.9in}p{1.3in}p{1.0in}p{0.5in}p{0.3in}p{0.4in}p{0.35in}p{1.65in}@{}}
\toprule
\textbf{Glyph} & \textbf{Macro} & \textbf{Role} & \textbf{Scope} & \textbf{Prec} & \textbf{First} & \textbf{Reg. page} & \textbf{Principal uses / notes} \\
\midrule
\endfirsthead
\toprule
\textbf{Glyph} & \textbf{Macro} & \textbf{Role} & \textbf{Scope} & \textbf{Prec} & \textbf{First} & \textbf{Reg. page} & \textbf{Principal uses / notes} \\
\midrule
\endhead
\midrule\multicolumn{8}{r}{\emph{continued}}\\\endfoot
\bottomrule\endlastfoot
\phantomsection\label{entry:symbol-registry-00055}\(\ulcorner\varphi\urcorner\) & \texttt{\detokenize{\ulcorner\cdot\urcorner}} & Goedel number & P5 & n/a & P5.1 & \thepage & Corner quotes around an expression denote its fixed effective code. \\
\phantomsection\label{entry:symbol-registry-00057}\(\mathsf{Proof}_T\) & \texttt{\detokenize{\mathsf{Proof}_T}} & proof predicate & P5 & n/a & P5.1 & \thepage & The fixed Sigma-zero-one formula \(\mathsf{Proof}_T(p,x)\) representing the primitive-recursive relation that p codes a T-proof with conclusion code x. \\
\phantomsection\label{entry:symbol-registry-00058}\(\mathsf{Prov}_T\) & \texttt{\detokenize{\mathsf{Prov}_T}} & provability predicate & P5 & n/a & P5.1 & \thepage & The exact inherited abbreviation \(\exists p\,\mathsf{Proof}_T(p,x)\); formula codes rather than formulas are its arguments. \\
\phantomsection\label{entry:symbol-registry-00059}\(\mathsf{Prov}_{\PA}\) & \texttt{\detokenize{\mathsf{Prov}_{\PA}}} & standard PA provability predicate & P5 & n/a & P5.6 & \thepage & The instance of the fixed inherited provability predicate used by the exact PA arithmetical realization in Solovay completeness. \\
\phantomsection\label{entry:symbol-registry-00060}\(\operatorname{ProofRel}_T,\operatorname{RProofRel}_T,\operatorname{NegCode}\) & \texttt{\detokenize{\operatorname{ProofRel}_T,}\allowbreak\detokenize{\operatorname{RProofRel}_T,}\allowbreak\detokenize{\operatorname{NegCode}}} & metalinguistic primitive-recursive proof relations and negation-code function & P5 & meta & P5.6 & \thepage & Relations and function on natural-number codes used before choosing object-language representatives; none is a term or predicate of the fixed arithmetic language. \\
\phantomsection\label{entry:symbol-registry-00061}\(\mathsf{RProof}_T\) & \texttt{\detokenize{\mathsf{RProof}_T}} & Rosser proof formula & P5 & n/a & P5.6 & \thepage & A fixed Sigma-zero-one formula numeralwise representing the primitive-recursive Rosser proof relation. \\
\phantomsection\label{entry:symbol-registry-00062}\(\mathsf{RProv}_T\) & \texttt{\detokenize{\mathsf{RProv}_T}} & Rosser provability predicate & P5 & n/a & P5.6 & \thepage & A proof of the coded sentence exists whose code is strictly smaller than every proof code for its negation; distinct from ordinary provability. \\
\phantomsection\label{entry:symbol-registry-00063}\(\mathsf{Con}(T)\) & \texttt{\detokenize{\mathsf{Con}}} & consistency sentence & P5-P7.6 & n/a & P5.2 & \thepage & \(\neg\mathsf{Prov}_T(\overline{\ulcorner\bottom_A\urcorner})\); P7 names the theory argument explicitly. \\
\phantomsection\label{entry:symbol-registry-00064}\(\mathsf{Con}_{\mathsf{Proof}_T}(T)\) & \texttt{\detokenize{\mathsf{Con}_{\mathsf{Proof}_T}}} & presentation-indexed consistency sentence & P5 & n/a & P5.5 & \thepage & Exposes the chosen proof predicate; for the fixed inherited predicate it is exactly \(\mathsf{Con}(T)\). \\
\phantomsection\label{entry:symbol-registry-00078}\(\varepsilon_0\) & \texttt{\detokenize{\varepsilon_0}} & least epsilon ordinal & P5-P7 & n/a & P5.4 & \thepage & Least fixed point of ordinal exponentiation base \(\omega\); calibrates PA in P7.5. \\
\phantomsection\label{entry:symbol-registry-00082}\(\leq\) & \texttt{\detokenize{\leq}} & object-language arithmetic order abbreviation & P5.1 & 2 & P5.1 & \thepage & The P4.3 abbreviation imported for the Q and Rosser chapters. \\
\phantomsection\label{entry:symbol-registry-00083}\(\leq\) & \texttt{\detokenize{\leq}} & canonical ordering & P5.2-P14 & 2 & P6.1 & \thepage & Ambient order on declared numerical and order-theoretic objects, including ordinals, cardinals, degrees, and frame orders. \\
\phantomsection\label{entry:symbol-registry-00091}\(\Box\) & \texttt{\detokenize{\Box}} & necessity operator & P5-P8.3 & 3 & P5.3 & \thepage & Prefix unary at the negation rank; modal and provability readings are fixed by the active system. \\
\phantomsection\label{entry:symbol-registry-00093}\(\Diamond\) & \texttt{\detokenize{\Diamond}} & possibility operator & P5-P8.3 & 3 & P5.3 & \thepage & Prefix unary at the negation rank; \(\neg\Box\neg\) where defined. \\
\phantomsection\label{entry:symbol-registry-00317}\(\operatorname{otp}(A,<)\) & \texttt{\detokenize{\operatorname{otp}}} & order type of a declared linear order & P5.4-P7 & n/a & P5.4 & \thepage & Metalinguistic isomorphism type of the displayed ordered set; not an ordinal unless the order is well founded. \\
\phantomsection\label{entry:symbol-registry-00318}\(\mathcal I,\varphi_{\mathrm{dom}},\psi^{\mathcal I}\) & \texttt{\detokenize{\mathcal I,}\allowbreak\detokenize{\varphi_{\mathrm{dom}},}\allowbreak\detokenize{\psi^{\mathcal I}}} & relative-interpretation data and formula translation & P5.4 & n/a & P5.4 & \thepage & A one-dimensional parameter-free interpretation; the domain formula and recursive translation are fixed at def:p05-relative-interpretation. \\
\phantomsection\label{entry:symbol-registry-00319}\(E(x,y),F_f(\vec x,y)\) & \texttt{\detokenize{E,}\allowbreak\detokenize{F_f}} & interpreted equality and function-graph formulas & P5.4 & n/a & P5.4 & \thepage & The target-language equality surrogate and total functional graph data fixed at def:p05-relative-interpretation; constants are nullary functions. \\
\phantomsection\label{entry:symbol-registry-00443}\(\mathsf Z_2\) & \texttt{\detokenize{\mathsf Z_2}} & full second-order arithmetic & P5.7-P7.5 & n/a & P5.7 & \thepage & RCA-zero augmented by full induction and full comprehension in the fixed two-sorted language; ambient theory for the displayed reverse-mathematics subsystems. \\
\phantomsection\label{entry:symbol-registry-00444}\(\mathsf{RCA}_0\) & \texttt{\detokenize{\mathsf{RCA}_0}} & recursive-comprehension base theory & P5.7-P7.5 & n/a & P5.7 & \thepage & The exact two-sorted base with Q number axioms, set extensionality, Sigma-zero-one induction, and Delta-zero-one comprehension is fixed in Part V. \\
\phantomsection\label{entry:symbol-registry-00445}\(\mathsf{WKL}_0\) & \texttt{\detokenize{\mathsf{WKL}_0}} & weak-Konig-lemma subsystem & P5.7-P7.5 & n/a & P5.7 & \thepage & RCA-zero plus the assertion that every infinite coded binary tree has an infinite path. \\
\phantomsection\label{entry:symbol-registry-00446}\(\mathsf{ACA}_0\) & \texttt{\detokenize{\mathsf{ACA}_0}} & arithmetical-comprehension theory & P5.7-P7.5 & n/a & P5.7 & \thepage & RCA-zero plus comprehension for every arithmetical formula; its first-order consequences are exactly those of PA. \\
\phantomsection\label{entry:symbol-registry-00447}\(\mathsf{ACA}\) & \texttt{\detokenize{\mathsf{ACA}}} & full-induction arithmetical-comprehension theory & P5.5 & n/a & P5.5 & \thepage & ACA-zero strengthened by induction for every formula of the two-sorted language; used to calibrate the uniform Paris-Harrington principle. \\
\phantomsection\label{entry:symbol-registry-00448}\(\mathsf{ATR}_0\) & \texttt{\detokenize{\mathsf{ATR}_0}} & arithmetical-transfinite-recursion theory & P5.7-P7.5 & n/a & P5.7 & \thepage & RCA-zero plus the arithmetical-transfinite-recursion scheme along coded well-orders. \\
\phantomsection\label{entry:symbol-registry-00449}\(\Pi^1_1\text{-}\mathsf{CA}_0\) & \texttt{\detokenize{\Pi^1_1\text{-}\mathsf{CA}_0}} & Pi-one-one comprehension theory & P5.7-P7.5 & n/a & P5.7 & \thepage & RCA-zero plus comprehension for formulas with universal set quantifiers over an arithmetical matrix. \\
\end{longtable}
\endgroup
\Needspace{.4\textheight}
\subsection*{Part 6 scopes}
\begingroup\scriptsize\setlength{\tabcolsep}{3pt}
\begin{longtable}{@{}p{0.9in}p{1.3in}p{1.0in}p{0.5in}p{0.3in}p{0.4in}p{0.35in}p{1.65in}@{}}
\toprule
\textbf{Glyph} & \textbf{Macro} & \textbf{Role} & \textbf{Scope} & \textbf{Prec} & \textbf{First} & \textbf{Reg. page} & \textbf{Principal uses / notes} \\
\midrule
\endfirsthead
\toprule
\textbf{Glyph} & \textbf{Macro} & \textbf{Role} & \textbf{Scope} & \textbf{Prec} & \textbf{First} & \textbf{Reg. page} & \textbf{Principal uses / notes} \\
\midrule
\endhead
\midrule\multicolumn{8}{r}{\emph{continued}}\\\endfoot
\bottomrule\endlastfoot
\phantomsection\label{entry:symbol-registry-00019}\(\alpha,\beta\) & \texttt{\detokenize{\alpha,}\allowbreak\detokenize{\beta}} & ordinal metavariables & P6 & n/a & P6.1 & \thepage & Ordinals of the ambient set theory; recurs under a local declaration in P7.5 ordinal analysis; disjoint from the type-metavariable scope. \\
\phantomsection\label{entry:symbol-registry-00065}\(V_\alpha\) & \texttt{\detokenize{V_\alpha}} & cumulative-hierarchy level & P6 & n/a & P6.1 & \thepage & \(V=\bigcup_\alpha V_\alpha\) is class notation, registered with its level notation. \\
\phantomsection\label{entry:symbol-registry-00066}\(\operatorname{rk}(x)\) & \texttt{\detokenize{\rank}} & set-theoretic rank & P6 & n/a & P6.1 & \thepage & Least \(\alpha\) with \(x\subseteq V_\alpha\). \\
\phantomsection\label{entry:symbol-registry-00067}\(L_\alpha\) & \texttt{\detokenize{L_\alpha}} & constructible level & P6.3 & n/a & P6.3 & \thepage & The \(\alpha\)-th level of the constructible hierarchy. \\
\phantomsection\label{entry:symbol-registry-00068}\(\mathbb{P}\) & \texttt{\detokenize{\mathbb{P}}} & forcing poset & P6.3-P6.4 & n/a & P6.3 & \thepage & \(q\leq p\) reads \(q\) is stronger, per the fixed forcing convention. \\
\phantomsection\label{entry:symbol-registry-00069}\(V^{\mathbb P}\) & \texttt{\detokenize{V^{\mathbb P}}} & class of P-names & P6.3 & n/a & P6.3 & \thepage & Name metavariables \(\nu\) range over \(M\cap V^{\mathbb P}\). \\
\phantomsection\label{entry:symbol-registry-00070}\(\check x\) & \texttt{\detokenize{\check x}} & check name & P6.3 & n/a & P6.3 & \thepage & Canonical forcing name of a ground-model set. \\
\phantomsection\label{entry:symbol-registry-00071}\(\nu^G\) & \texttt{\detokenize{\nu^G}} & name evaluation & P6.3 & n/a & P6.3 & \thepage & Interpretation by the generic filter \(G\) of a locally declared forcing name \(\nu,\nu_f\), or \(\zeta\). \\
\phantomsection\label{entry:symbol-registry-00072}\(M[G]\) & \texttt{\detokenize{M[G]}} & generic extension & P6.3-P6.6 & n/a & P6.3 & \thepage & The forcing extension obtained by adjoining the \(M\)-generic filter \(G\). \\
\phantomsection\label{entry:symbol-registry-00073}\(p\Vdash^{\mathrm{sem}}_{M,\mathbb P}\varphi\) & \texttt{\detokenize{\forces^{\mathrm{sem}}_{M,\mathbb P}}} & semantic set-theoretic forcing & P6.3 & meta & P6.3 & \thepage & Exact external relation: every M-generic filter containing p yields an extension satisfying the formula at the interpreted names. \\
\phantomsection\label{entry:symbol-registry-00074}\(p\Vdash_{M,\mathbb P}\varphi\) & \texttt{\detokenize{\forces_{M,\mathbb P}}} & imported recursive forcing assertion & P6.3 & meta & P6.3 & \thepage & Formula-by-formula recursive assertion imported by the cited forcing theorem; equivalent there to the separately registered semantic relation. Bare \(p\Vdash\varphi\) is its local abbreviation only under the renewed ccc/name-counting convention. \\
\phantomsection\label{entry:symbol-registry-00320}\(\nu,\nu_f,\zeta\) & \texttt{\detokenize{\nu,}\allowbreak\detokenize{\nu_f,}\allowbreak\detokenize{\zeta}} & forcing-name metavariables & P6.3 & n/a & P6.3 & \thepage & Members of the active P-name hierarchy; subscripts identify a name by its intended value and never turn a name into that value. \\
\phantomsection\label{entry:symbol-registry-00321}\(\mathbb B\) & \texttt{\detokenize{\mathbb B}} & complete Boolean algebra used for Boolean-valued models & P6.3 & n/a & P6.3 & \thepage & Scope is disjoint from the P2.1 Boolean truth-value set; completeness supplies all joins and meets used by the Boolean-valued recursion. \\
\phantomsection\label{entry:symbol-registry-00322}\(V^{\mathbb B}_\alpha,V^{\mathbb B}\) & \texttt{\detokenize{V^{\mathbb B}_\alpha,}\allowbreak\detokenize{V^{\mathbb B}}} & ambient Boolean-valued name hierarchy & P6.3 & n/a & P6.3 & \thepage & Names are Boolean-valued functions on earlier levels for an ambient-complete algebra; distinct from the P-name hierarchy and the ground-model-computed variant. \\
\phantomsection\label{entry:symbol-registry-00323}\(\operatorname{Bool}_{\mathbb B}(\varphi)\) & \texttt{\detokenize{\operatorname{Bool}_{\mathbb B}}} & ambient Boolean truth value of a formula & P6.3 & n/a & P6.3 & \thepage & Value in an ambient-complete Boolean algebra; this is not the two-valued satisfaction relation. \\
\phantomsection\label{entry:symbol-registry-00324}\((V^{\mathbb B})^M,\operatorname{Bool}_{\mathbb B}^M,S_{\varphi,\bar{\mathbf v}}^M\) & \texttt{\detokenize{(V^{\mathbb B})^M,}\allowbreak\detokenize{\operatorname{Bool}_{\mathbb B}^M,}\allowbreak\detokenize{S_{\varphi,\bar{\mathbf v}}^M}} & ground-model-computed Boolean names truth and attained-value set & P6.3 & n/a & P6.3 & \thepage & All recursions, Separation, and joins are computed in M when M regards B as complete; ambient completeness is not assumed. \\
\phantomsection\label{entry:symbol-registry-00325}\(b^{\complement_{\mathbb B}}\) & \texttt{\detokenize{{}^{\complement_{\mathbb B}}}} & Boolean-algebra complement & P6.3 & 1 & P6.3 & \thepage & Postfix operation in the declared algebra; the subscript prevents collision with object-language negation. \\
\phantomsection\label{entry:symbol-registry-00326}\(b\wedge_{\mathbb B}c\) & \texttt{\detokenize{\wedge_{\mathbb B}}} & Boolean-algebra meet & P6.3 & 4 & P6.3 & \thepage & Operation in the declared algebra; distinct from the globally registered object-language conjunction. \\
\phantomsection\label{entry:symbol-registry-00327}\(b\vee_{\mathbb B}c\) & \texttt{\detokenize{\vee_{\mathbb B}}} & Boolean-algebra join & P6.3 & 5 & P6.3 & \thepage & Operation in the declared algebra; distinct from the globally registered object-language disjunction. \\
\phantomsection\label{entry:symbol-registry-00328}\(\mathbf u,\mathbf v,\mathbf w\) & \texttt{\detokenize{\mathbf u,}\allowbreak\detokenize{\mathbf v,}\allowbreak\detokenize{\mathbf w}} & Boolean-valued name metavariables & P6.3 & n/a & P6.3 & \thepage & Members of \(V^{\mathbb B}\); bold forms avoid the globally reserved valuation and term metavariables. \\
\phantomsection\label{entry:symbol-registry-00329}\(\mathbf u\sim_G\mathbf v\) & \texttt{\detokenize{\sim_G}} & generic-ultrafilter equivalence of Boolean names & P6.3 & meta & P6.3 & \thepage & Defined by Boolean equality having value in \(G\); the quotient relation is not ordinary set equality. \\
\phantomsection\label{entry:symbol-registry-00330}\(\mathcal G,\mathfrak F\) & \texttt{\detokenize{\mathcal G,}\allowbreak\detokenize{\mathfrak F}} & permutation group and normal subgroup filter & P6.3 & n/a & P6.3 & \thepage & Typed components of a permutation-model datum; \(\mathfrak F\) contains subgroups of \(\mathcal G\). \\
\phantomsection\label{entry:symbol-registry-00331}\(\operatorname{Sym}(A),\operatorname{Stab}_{\mathcal G}(x)\) & \texttt{\detokenize{\operatorname{Sym},}\allowbreak\detokenize{\operatorname{Stab}_{\mathcal G}}} & full permutation group and stabilizer & P6.3 & n/a & P6.3 & \thepage & The stabilizer is a subgroup of the declared permutation group; support and symmetry are defined from it. \\
\phantomsection\label{entry:symbol-registry-00332}\(\omega^\omega\) & \texttt{\detokenize{\omega^\omega}} & ordinal exponentiation & P6.1 & n/a & P6.1 & \thepage & In the ordinal-arithmetic chapter this is the ordinal power obtained by right-recursion. \\
\phantomsection\label{entry:symbol-registry-00333}\(\omega^\omega\) & \texttt{\detokenize{\omega^\omega}} & Baire function space & P6.6-P6.8 & n/a & P6.6 & \thepage & The set of functions from \(\omega\) to \(\omega\); Chapter 6.7 adds the product topology. Scope is disjoint from the ordinal-power role. \\
\phantomsection\label{entry:symbol-registry-00334}\(\boldsymbol{\Sigma}^{0}_{\alpha},\boldsymbol{\Pi}^{0}_{\alpha},\boldsymbol{\Delta}^{0}_{\alpha}\) & \texttt{\detokenize{\boldsymbol{\Sigma}^{0}_{\alpha},}\allowbreak\detokenize{\boldsymbol{\Pi}^{0}_{\alpha},}\allowbreak\detokenize{\boldsymbol{\Delta}^{0}_{\alpha}}} & boldface Borel pointclasses & P6.7 & n/a & P6.7 & \thepage & Countable-ordinal Borel levels on a Polish space; distinct from the P4-P5 bold arithmetical-definability classes. \\
\phantomsection\label{entry:symbol-registry-00335}\(\boldsymbol{\Sigma}^{1}_{n},\boldsymbol{\Pi}^{1}_{n},\boldsymbol{\Delta}^{1}_{n}\) & \texttt{\detokenize{\boldsymbol{\Sigma}^{1}_{n},}\allowbreak\detokenize{\boldsymbol{\Pi}^{1}_{n},}\allowbreak\detokenize{\boldsymbol{\Delta}^{1}_{n}}} & boldface projective pointclasses & P6.6-P6.8 & n/a & P6.6 & \thepage & Real-parameter projective classes; distinct from lightface analytical classes and from the Levy hierarchy. \\
\phantomsection\label{entry:symbol-registry-00336}\(\operatorname{Bor}(X)\) & \texttt{\detokenize{\operatorname{Bor}}} & Borel sigma-algebra of a Polish space & P6.7 & n/a & P6.7 & \thepage & The least sigma-algebra containing the open subsets of \(X\). \\
\phantomsection\label{entry:symbol-registry-00337}\(\operatorname{crit}(j)\) & \texttt{\detokenize{\operatorname{crit}}} & critical point of an elementary embedding & P6.6-P6.8 & n/a & P6.6 & \thepage & Least ordinal moved by \(j\); subscripted embedding symbols such as \(j_U\) use the same operation. \\
\phantomsection\label{entry:symbol-registry-00338}\(G(A)\) & \texttt{\detokenize{G(A)}} & Gale--Stewart payoff game & P6.8 & n/a & P6.8 & \thepage & Length-omega perfect-information game whose payoff set is \(A\subseteq\omega^\omega\). \\
\phantomsection\label{entry:symbol-registry-00339}\(\mathsf{AD},\mathsf{PD}\) & \texttt{\detokenize{\mathsf{AD},}\allowbreak\detokenize{\mathsf{PD}}} & determinacy assertions & P6.8 & n/a & P6.8 & \thepage & AD quantifies over all payoff sets of reals; PD restricts to boldface projective payoff sets. \\
\phantomsection\label{entry:symbol-registry-00340}\(\mathsf{ZF},\mathsf{ZFC}\) & \texttt{\detokenize{\mathsf{ZF},}\allowbreak\detokenize{\mathsf{ZFC}}} & first-order set theories without and with Choice & P6 & n/a & P6.1 & \thepage & ZF is the displayed axiom basis without Choice; ZFC adds the displayed object-language Axiom of Choice. \\
\phantomsection\label{entry:symbol-registry-00341}\(\mathsf Z,\mathsf{ZC},\mathsf{KP},\mathsf{KP}^{\mathcal P}\) & \texttt{\detokenize{\mathsf Z,}\allowbreak\detokenize{\mathsf{ZC},}\allowbreak\detokenize{\mathsf{KP},}\allowbreak\detokenize{\mathsf{KP}^{\mathcal P}}} & set-theory system abbreviations & P6.2 & n/a & P6.2 & \thepage & Z omits Replacement from the stated ZF basis; ZC adds Choice; KP is the stated admissible-set basis and its power-set variant is explicitly superscripted. \\
\phantomsection\label{entry:symbol-registry-00342}\(\mathsf{GB},\mathsf{NBG},\mathsf{MK},\mathsf{KM}\) & \texttt{\detokenize{\mathsf{GB},}\allowbreak\detokenize{\mathsf{NBG},}\allowbreak\detokenize{\mathsf{MK},}\allowbreak\detokenize{\mathsf{KM}}} & set-and-class theory abbreviations & P6.2 & n/a & P6.2 & \thepage & The comparison table fixes GB as choice-free NBG and treats MK and KM as names for the declared impredicative class theory convention. \\
\phantomsection\label{entry:symbol-registry-00343}\(\mathsf{NF},\mathsf{TTT}\) & \texttt{\detokenize{\mathsf{NF},}\allowbreak\detokenize{\mathsf{TTT}}} & New Foundations and Tangled Type Theory abbreviations & P6.2 & n/a & P6.2 & \thepage & Distinct theories linked only by the cited consistency reduction; neither is presented as a subtheory of ZF. \\
\phantomsection\label{entry:symbol-registry-00344}\(\mathsf{ZFA}\) & \texttt{\detokenize{\mathsf{ZFA}}} & set theory with atoms & P6.3 & n/a & P6.3 & \thepage & Theory of the local permutation-model example; the cited transfer is required before drawing a pure-ZF conclusion. \\
\phantomsection\label{entry:symbol-registry-00345}\(\mathsf{AC}\) & \texttt{\detokenize{\mathsf{AC}}} & Axiom of Choice & P6 & n/a & P6.1 & \thepage & Object-language set-theory axiom; occurrences in metatheoretic dependency tables name that exact principle. \\
\phantomsection\label{entry:symbol-registry-00346}\(\mathsf{DC},\mathsf{AC}_{\omega},\mathsf{AC}_{\omega}^{\mathrm{fin}},\mathsf{BPI}\) & \texttt{\detokenize{\mathsf{DC},}\allowbreak\detokenize{\mathsf{AC}_{\omega},}\allowbreak\detokenize{\mathsf{AC}_{\omega}^{\mathrm{fin}},}\allowbreak\detokenize{\mathsf{BPI}}} & restricted or weaker choice principles & P6.2-P6.4 & n/a & P6.2 & \thepage & Dependent Choice, Countable Choice, Countable Choice for finite sets, and the Boolean Prime Ideal Theorem; implication arrows are stated separately. \\
\phantomsection\label{entry:symbol-registry-00347}\(\kappa,\lambda,\mu\) & \texttt{\detokenize{\kappa,}\allowbreak\detokenize{\lambda,}\allowbreak\detokenize{\mu}} & ordinal or cardinal metavariables as explicitly quantified & P6 & n/a & P6.1 & \thepage & The local quantifier fixes the standing: limit ordinals and conditional hierarchy markers occur before the cardinal-arithmetic chapters. \\
\phantomsection\label{entry:symbol-registry-00348}\(\gamma,\delta,\xi\) & \texttt{\detokenize{\gamma,}\allowbreak\detokenize{\delta,}\allowbreak\detokenize{\xi}} & additional ordinal metavariables & P6 & n/a & P6.1 & \thepage & Ordinals in recursive, limit, club, forcing, and hierarchy clauses; complements the existing alpha-beta ordinal row. \\
\phantomsection\label{entry:symbol-registry-00349}\(\kappa^+\) & \texttt{\detokenize{{}^+}} & successor-cardinal operation & P6.3-P6.8 & 1 & P6.3 & \thepage & Least cardinal strictly above kappa; distinct from ordinal successor alpha+1. \\
\phantomsection\label{entry:symbol-registry-00350}\(\kappa^\lambda,\lambda^{<\kappa}\) & \texttt{\detokenize{\kappa^\lambda,}\allowbreak\detokenize{\lambda^{<\kappa}}} & cardinal exponentiation and bounded-exponent supremum & P6.3-P6.6 & n/a & P6.3 & \thepage & Cardinal function-space size and the supremum over exponents below kappa; the full entry is Chapter 6.5. \\
\phantomsection\label{entry:symbol-registry-00351}\(\aleph_\alpha,\beth_\alpha\) & \texttt{\detokenize{\aleph_\alpha,}\allowbreak\detokenize{\beth_\alpha}} & aleph and beth cardinal hierarchies & P6.3-P6.8 & n/a & P6.3 & \thepage & Aleph enumerates infinite initial ordinals and beth iterates powerset cardinality; limit stages take suprema. \\
\phantomsection\label{entry:symbol-registry-00352}\(\operatorname{cf}(\alpha),\operatorname{cf}(P,\leq)\) & \texttt{\detokenize{\operatorname{cf}}} & cofinality of an ordinal or ordered set & P6.3-P6.8 & n/a & P6.3 & \thepage & For an ordinal: least cofinal order type. For a nonempty ordered set: least cardinality of a cofinal family. Regularity is cf(kappa)=kappa. \\
\phantomsection\label{entry:symbol-registry-00353}\(\mathsf{CH},\mathsf{GCH}\) & \texttt{\detokenize{\mathsf{CH},}\allowbreak\detokenize{\mathsf{GCH}}} & continuum hypotheses & P6.3-P6.6 & n/a & P6.3 & \thepage & CH is the aleph-one equation at the continuum; GCH is its all-aleph scheme. \\
\phantomsection\label{entry:symbol-registry-00354}\(\mathfrak c\) & \texttt{\detokenize{\mathfrak c}} & cardinality of the continuum & P6.4 & n/a & P6.4 & \thepage & Defined as two to the aleph-zero; used only to bound the Martin-axiom scheme. \\
\phantomsection\label{entry:symbol-registry-00355}\(\mathsf{Con}(\mathsf{ZFC})^\in\) & \texttt{\detokenize{\mathsf{Con}(\mathsf{ZFC})^\in}} & membership-language translation of the arithmetical ZFC-consistency sentence & P6.3 & n/a & P6.3 & \thepage & Obtained by translating the fixed arithmetical no-proof sentence along the canonical interpretation of arithmetic by the finite ordinals; the separate countable-transitive-model hypothesis is stronger and must not be conflated with this sentence. \\
\phantomsection\label{entry:symbol-registry-00356}\(2^{<\omega},2^\omega\) & \texttt{\detokenize{2^{<\omega},}\allowbreak\detokenize{2^\omega}} & finite and infinite binary-sequence spaces & P6.2-P6.4 & n/a & P6.2 & \thepage & Finite strings and omega-sequences over two; the first carries prefix order and the second supplies branches or reals as locally declared. \\
\phantomsection\label{entry:symbol-registry-00357}\(\mathcal T\) & \texttt{\detokenize{\mathcal T}} & tree metavariable & P6.2-P6.4 & n/a & P6.2 & \thepage & A binary or set-theoretic tree as locally declared; scope is disjoint from the P4.1-P4.2 Turing-machine role. \\
\phantomsection\label{entry:symbol-registry-00358}\(\mathbf s\) & \texttt{\detokenize{\mathbf s}} & finite binary-sequence metavariable & P6.2 & n/a & P6.2 & \thepage & Node of the canonical binary tree; bold form avoids the globally reserved assignment metavariable s. \\
\phantomsection\label{entry:symbol-registry-00359}\(\operatorname{Fn}(I,J,{<}\delta)\) & \texttt{\detokenize{\operatorname{Fn}}} & bounded-size partial-function poset & P6.3 & n/a & P6.3 & \thepage & Partial functions from I to J of size below delta, ordered by reverse inclusion. \\
\phantomsection\label{entry:symbol-registry-00360}\(\operatorname{Add}(\omega,\kappa)\) & \texttt{\detokenize{\operatorname{Add}}} & Cohen forcing adding kappa many reals & P6.3-P6.4 & n/a & P6.3 & \thepage & The explicitly displayed Fn poset; preservation and continuum claims retain their separate hypotheses. \\
\phantomsection\label{entry:symbol-registry-00361}\(\operatorname{Coll}(\kappa,\lambda)\) & \texttt{\detokenize{\operatorname{Coll}}} & Levy collapse forcing & P6.3 & n/a & P6.3 & \thepage & Partial functions from kappa to lambda of size below kappa under the catalog convention; collapse claims are cited. \\
\phantomsection\label{entry:symbol-registry-00362}\(1_{\mathbb P}\) & \texttt{\detokenize{1_{\mathbb P}}} & weakest forcing condition & P6.3 & n/a & P6.3 & \thepage & Top element under the convention that q less-than-or-equal p means q is stronger. \\
\phantomsection\label{entry:symbol-registry-00363}\(\mathop{\triangle}_{\xi<\kappa}C_\xi\) & \texttt{\detokenize{\mathop{\triangle}}} & diagonal intersection & P6.4 & n/a & P6.4 & \thepage & Set of \(\alpha<\kappa\) belonging to every earlier \(C_\xi\); distinct from symmetric difference. \\
\phantomsection\label{entry:symbol-registry-00364}\([X]^\ell\) & \texttt{\detokenize{[X]^\ell}} & ell-element subsets of a set & P6.4-P6.6 & n/a & P6.4 & \thepage & Domain of finite-color partition relations; ell is a positive finite integer. \\
\phantomsection\label{entry:symbol-registry-00365}\(\kappa\longrightarrow(\lambda)^\ell_\mu\) & \texttt{\detokenize{\longrightarrow}} & partition relation & P6.4-P6.6 & meta & P6.4 & \thepage & Every mu-coloring of the ell-element subsets of kappa has a homogeneous subset of cardinality lambda; not metalinguistic implication. \\
\phantomsection\label{entry:symbol-registry-00366}\(\beth_n(\kappa)\) & \texttt{\detokenize{\beth_n(\kappa)}} & finite beth iteration above a cardinal & P6.4 & n/a & P6.4 & \thepage & Starts at kappa and iterates powerset cardinality n times; distinct from the global beth hierarchy indexed from aleph-zero. \\
\phantomsection\label{entry:symbol-registry-00367}\(\mathcal T_\alpha,\operatorname{ht}_{\mathcal T}(\mathbf t)\) & \texttt{\detokenize{\mathcal T_\alpha,}\allowbreak\detokenize{\operatorname{ht}_{\mathcal T}}} & tree level and node height & P6.4 & n/a & P6.4 & \thepage & Height is the predecessor order type and the alpha-level contains exactly nodes of that height. \\
\phantomsection\label{entry:symbol-registry-00368}\(\mathbf t\) & \texttt{\detokenize{\mathbf t}} & set-theoretic tree-node metavariable & P6.4 & n/a & P6.4 & \thepage & Node of the locally declared tree; bold form avoids the globally reserved term metavariable t. \\
\phantomsection\label{entry:symbol-registry-00369}\(\mathsf{SH}\) & \texttt{\detokenize{\mathsf{SH}}} & Suslin hypothesis & P6.4 & n/a & P6.4 & \thepage & Assertion that no Suslin tree exists under the chapter's tree convention. \\
\phantomsection\label{entry:symbol-registry-00370}\(\diamondsuit_{\omega_1}\) & \texttt{\detokenize{\diamondsuit_{\omega_1}}} & diamond principle at omega-one & P6.4 & n/a & P6.4 & \thepage & Stationary-guessing sequence defined in the combinatorial-principles entry. \\
\phantomsection\label{entry:symbol-registry-00371}\(\square_\kappa\) & \texttt{\detokenize{\square_\kappa}} & square principle & P6.4 & n/a & P6.4 & \thepage & Coherent club sequence with the explicit no-thread condition; distinct from the modal necessity glyph. \\
\phantomsection\label{entry:symbol-registry-00372}\(\mathsf{MA}_\kappa,\mathsf{MA}\) & \texttt{\detokenize{\mathsf{MA}_\kappa,}\allowbreak\detokenize{\mathsf{MA}}} & Martin axiom instance and scheme & P6.4 & n/a & P6.4 & \thepage & The instance meets at most kappa dense sets in a ccc poset; the scheme ranges below the continuum. \\
\phantomsection\label{entry:symbol-registry-00373}\(\Delta_0,\Sigma_n,\Pi_n,\Delta_n\) & \texttt{\detokenize{\Delta_0,}\allowbreak\detokenize{\Sigma_n,}\allowbreak\detokenize{\Pi_n,}\allowbreak\detokenize{\Delta_n}} & Levy-hierarchy formula classes & P6.6 & n/a & P6.6 & \thepage & Set-theoretic bounded and unbounded-quantifier classes; distinct from arithmetical and projective hierarchies. \\
\phantomsection\label{entry:symbol-registry-00374}\(\mathcal U\) & \texttt{\detokenize{\mathcal U}} & kappa-complete ultrafilter witnessing measurability & P6.6 & n/a & P6.6 & \thepage & Calligraphic form avoids the globally reserved bare theory metavariable U. \\
\phantomsection\label{entry:symbol-registry-00375}\(j:V\to M\) & \texttt{\detokenize{j:V\to M}} & elementary class embedding & P6.6-P6.8 & n/a & P6.6 & \thepage & Metalinguistic embedding read formula by formula; neither V nor M is added as a class object to ZFC. \\
\phantomsection\label{entry:symbol-registry-00376}\({}^{\lambda}M\subseteq M\) & \texttt{\detokenize{{}^{\lambda}M\subseteq M}} & closure under ambient lambda-sequences & P6.6-P6.8 & meta & P6.6 & \thepage & Every ambient function from lambda into M belongs to M; stronger than transitivity. \\
\phantomsection\label{entry:symbol-registry-00377}\(\operatorname{Ult}(V,\mathcal U)\) & \texttt{\detokenize{\operatorname{Ult}}} & universe ultrapower by a kappa-complete ultrafilter & P6.6 & n/a & P6.6 & \thepage & Class ultrapower whose well-foundedness and transitive collapse are cited. \\
\phantomsection\label{entry:symbol-registry-00378}\(\Lang_{\in}\) & \texttt{\detokenize{\Lang_{\in}}} & first-order membership signature & P6 & n/a & P6.1 & \thepage & One-sorted signature with equality and membership as its sole nonlogical relation; classes remain metalinguistic. \\
\phantomsection\label{entry:symbol-registry-00379}\(\operatorname{Ord}\) & \texttt{\detokenize{\operatorname{Ord}}} & proper class of all ordinals & P6 & n/a & P6.1 & \thepage & Metalinguistic class notation and never a set-valued union index. \\
\phantomsection\label{entry:symbol-registry-00380}\(\omega\) & \texttt{\detokenize{\omega}} & least infinite ordinal and von Neumann natural-number set & P6-P7.5 & n/a & P6.1 & \thepage & Least inductive set under the displayed ZF construction; P7.5 uses only the ordinal standing. \\
\phantomsection\label{entry:symbol-registry-00381}\(V\) & \texttt{\detokenize{V}} & ambient universe of sets & P6 & n/a & P6.1 & \thepage & Metalinguistic class abbreviation for the union of the cumulative hierarchy; not a set or a valuation. \\
\phantomsection\label{entry:symbol-registry-00382}\(\sup A\) & \texttt{\detokenize{\sup}} & ordinal supremum & P6-P7.5 & n/a & P6.1 & \thepage & Union of a set of ordinals and the least ordinal at least every member; extended to the P7.5 ordinal towers. \\
\phantomsection\label{entry:symbol-registry-00383}\(\alpha+\beta,\alpha\cdot\beta,\alpha^\beta\) & \texttt{\detokenize{+,}\allowbreak\detokenize{\cdot,}\allowbreak\detokenize{{}^{}}} & ordinal arithmetic operations & P6.1-P7.5 & n/a & P6.1 & \thepage & Right-recursive ordinal sum product and exponentiation; distinct from cardinal arithmetic by the declared operands and chapter scope. \\
\phantomsection\label{entry:symbol-registry-00384}\(A\approx B\) & \texttt{\detokenize{\approx}} & equipotence of sets & P6.2 & meta & P6.2 & \thepage & There is a bijection from A to B; distinct from isomorphism of structures and from formula equivalence. \\
\phantomsection\label{entry:symbol-registry-00385}\(h(A)\) & \texttt{\detokenize{h(A)}} & Hartogs number of a set & P6.2 & n/a & P6.2 & \thepage & Least ordinal admitting no injection into A; its existence is a ZF theorem. \\
\phantomsection\label{entry:symbol-registry-00386}\(\operatorname{Def}(A)\) & \texttt{\detokenize{\operatorname{Def}}} & definable subsets of a membership structure & P6.3 & n/a & P6.3 & \thepage & Subsets of A first-order definable over \((A,\in)\) with finitely many parameters from A. \\
\phantomsection\label{entry:symbol-registry-00387}\(L\) & \texttt{\detokenize{L}} & constructible universe & P6.3-P6.6 & n/a & P6.3 & \thepage & Proper class union of the constructible levels; distinct from the global signature macro and from a set-sized level. \\
\phantomsection\label{entry:symbol-registry-00441}\(\mathsf{PRA}\) & \texttt{\detokenize{\mathsf{PRA}}} & primitive recursive arithmetic & P6-P7.5 & n/a & P6.6 & \thepage & Fixed for the master strength map by primitive-recursive function symbols, defining equations, and quantifier-free induction; developed systematically in P7.5. \\
\phantomsection\label{entry:symbol-registry-00442}\(\mathsf I\Sigma_n\) & \texttt{\detokenize{\mathsf I\Sigma_n}} & induction-fragment family & P6-P7.5 & n/a & P6.6 & \thepage & Robinson arithmetic plus induction exactly for Sigma-zero-n formulas, for n at least one; P6 first uses the n=1 conservativity calibration and P7.5 gives the systematic treatment. \\
\phantomsection\label{entry:symbol-registry-00690}\(\prod\mathfrak a\) & \texttt{\detokenize{\prod}} & product space of regular cardinals & P6.5 & n/a & P6.5 & \thepage & Functions with f(kappa) below kappa at every coordinate in the nonempty set a; the cardinality of this space is the cardinal product. \\
\phantomsection\label{entry:symbol-registry-00691}\(f={}_Dg,\ f\leq_Dg,\ f<_Dg\) & \texttt{\detokenize{={}_D,}\allowbreak\detokenize{\leq_D,}\allowbreak\detokenize{<_D}} & equality and comparison modulo a filter & P6.5 & meta & P6.5 & \thepage & The set of coordinates satisfying the indicated comparison belongs to D; for an ultrafilter the quotient order is linear. \\
\phantomsection\label{entry:symbol-registry-00692}\(\operatorname{pcf}(\mathfrak a)\) & \texttt{\detokenize{\operatorname{pcf}}} & possible cofinalities of a product & P6.5 & n/a & P6.5 & \thepage & Cofinalities of the ultrafilter quotients of the product of the infinite regular cardinals in a; principal ultrafilters are included. \\
\phantomsection\label{entry:symbol-registry-00693}\(\omega_4\) & \texttt{\detokenize{\omega_4}} & fourth uncountable initial ordinal & P6.5 & n/a & P6.5 & \thepage & Initial ordinal of cardinality aleph-four; used as the aleph index in Shelah's bound at aleph-omega. \\
\end{longtable}
\endgroup
\Needspace{.4\textheight}
\subsection*{Part 7 scopes}
\begingroup\scriptsize\setlength{\tabcolsep}{3pt}
\begin{longtable}{@{}p{0.9in}p{1.3in}p{1.0in}p{0.5in}p{0.3in}p{0.4in}p{0.35in}p{1.65in}@{}}
\toprule
\textbf{Glyph} & \textbf{Macro} & \textbf{Role} & \textbf{Scope} & \textbf{Prec} & \textbf{First} & \textbf{Reg. page} & \textbf{Principal uses / notes} \\
\midrule
\endfirsthead
\toprule
\textbf{Glyph} & \textbf{Macro} & \textbf{Role} & \textbf{Scope} & \textbf{Prec} & \textbf{First} & \textbf{Reg. page} & \textbf{Principal uses / notes} \\
\midrule
\endhead
\midrule\multicolumn{8}{r}{\emph{continued}}\\\endfoot
\bottomrule\endlastfoot
\phantomsection\label{entry:symbol-registry-00020}\(\alpha,\beta,\gamma\) & \texttt{\detokenize{\alpha,}\allowbreak\detokenize{\beta,}\allowbreak\detokenize{\gamma}} & type metavariables & P7.3-P7.4 & n/a & P7.3 & \thepage & Scope is disjoint from the P7.5 ordinal role and the P7.6 finite-type role. \\
\phantomsection\label{entry:symbol-registry-00388}\(\mathsf{LK}_{\mathrm d},\mathsf{LK}_{=}\) & \texttt{\detokenize{\mathsf{LK}_{\mathrm d},}\allowbreak\detokenize{\mathsf{LK}_{=}}} & defined-negation and equality sequent-calculus identifiers & P7.1-P7.2 & n/a & P7.1 & \thepage & Exact registered presentations; neither identifier is an alias for canonical primitive-negation LK. \\
\phantomsection\label{entry:symbol-registry-00389}\(\mathsf{LJ}\) & \texttt{\detokenize{\mathsf{LJ}}} & intuitionistic sequent-calculus identifier & P7.1-P8.1 & n/a & P7.1 & \thepage & Exact single-succedent calculus fixed by sheet:lj; reused for IPL and equality-free IQC with every turnstile subscripted. \\
\phantomsection\label{entry:symbol-registry-00390}\(\mathsf L\) & \texttt{\detokenize{\mathsf L}} & associative Lambek-calculus identifier & P7.1-P8.6 & n/a & P7.1 & \thepage & The cited associative ordered nonempty-antecedent Lambek calculus; \(\mathsf L^*\) and \(\mathsf L_1\) are source-named contrast extensions rather than aliases. \\
\phantomsection\label{entry:symbol-registry-00391}\(d(\varphi),d(\Gamma)\) & \texttt{\detokenize{d}} & defined-negation formula and context translation & P7.1 & n/a & P7.1 & \thepage & Recursive equality-free translation from canonical LK syntax to the defined-negation language; context action is occurrencewise. \\
\phantomsection\label{entry:symbol-registry-00392}\(\bigwedge\Gamma\) & \texttt{\detokenize{\bigwedge}} & finite context conjunction & P7.1 & 4 & P7.1 & \thepage & Metalinguistic fold of a finite context with empty value top; distinct from the P2.4 infinitary object-language connective. \\
\phantomsection\label{entry:symbol-registry-00393}\(\bigvee\Delta\) & \texttt{\detokenize{\bigvee}} & finite context disjunction & P7.1 & 5 & P7.1 & \thepage & Metalinguistic fold of a finite context with empty value bottom; distinct from the P2.4 infinitary object-language connective. \\
\phantomsection\label{entry:symbol-registry-00394}\(D,D_1,D_2\) & \texttt{\detokenize{D,}\allowbreak\detokenize{D_1,}\allowbreak\detokenize{D_2}} & derivation metavariables & P7.2-P7.6 & n/a & P7.2 & \thepage & Finite derivation trees in the locally fixed calculus; distinct from the roman Dialectica superscript. \\
\phantomsection\label{entry:symbol-registry-00395}\(\operatorname{rk}_{\mathrm{fm}}(\varphi)\) & \texttt{\detokenize{\operatorname{rk}_{\mathrm{fm}}}} & formula rank & P7.2 & n/a & P7.2 & \thepage & Number of logical symbols in the displayed formula; the subscript separates it from set-theoretic rank. \\
\phantomsection\label{entry:symbol-registry-00396}\(\operatorname{ht}(D)\) & \texttt{\detokenize{\operatorname{ht}}} & derivation height & P7.2 & n/a & P7.2 & \thepage & Greatest inference-step count on a leaf-to-root branch; distinct from tree-node height in P6.4. \\
\phantomsection\label{entry:symbol-registry-00397}\(\varphi^k\) & \texttt{\detokenize{{}^k}} & repeated Mix-formula occurrence & P7.1-P7.2 & n/a & P7.1 & \thepage & Exactly k comma-separated occurrences on the registered Mix sheets; not exponentiation or iteration. \\
\phantomsection\label{entry:symbol-registry-00398}\(\mathsf{NJ}\) & \texttt{\detokenize{\mathsf{NJ}}} & intuitionistic natural-deduction calculus identifier & P7.1-P7.3 & n/a & P7.1 & \thepage & Exact calculus fixed by sheet:nj; excludes RAA, primitive negation, equality, and induction. \\
\phantomsection\label{entry:symbol-registry-00399}\(\longrightarrow_{\mathsf{NJ}}\) & \texttt{\detokenize{\longrightarrow_{\mathsf{NJ}}}} & one-step NJ derivation reduction & P7.3 & meta & P7.3 & \thepage & Compatible closure of the exact principal reductions and permutation conversions fixed in the normalization entry. \\
\phantomsection\label{entry:symbol-registry-00400}\(\longrightarrow\) & \texttt{\detokenize{\longrightarrow}} & abstract reduction relation & P7.3 & meta & P7.3 & \thepage & Locally declared reduction relation used only in the abstract normalization definitions; not implication or a sequent arrow. \\
\phantomsection\label{entry:symbol-registry-00401}\(\mathsf{Base},\mathsf b\) & \texttt{\detokenize{\mathsf{Base},}\allowbreak\detokenize{\mathsf b}} & base-type set and base-type metavariable & P7.3-P7.4 & n/a & P7.3 & \thepage & Fixed generating set for the arrow-only STLC type grammar and one member of it. \\
\phantomsection\label{entry:symbol-registry-00402}\(\mathsf{STLC}\) & \texttt{\detokenize{\mathsf{STLC}}} & arrow-only simply typed lambda-calculus identifier & P7.3-P7.4 & n/a & P7.3 & \thepage & Exact calculus fixed by sheet:stlc and used in every corresponding typing turnstile. \\
\phantomsection\label{entry:symbol-registry-00403}\(\Theta_{\mathrm{ty}}\) & \texttt{\detokenize{\Theta_{\mathrm{ty}}}} & term-variable typing context & P7.3-P7.6 & n/a & P7.3 & \thepage & Finite list of distinct term-variable declarations carrying well-formed types in the active typed calculus. \\
\phantomsection\label{entry:symbol-registry-00404}\(M,N,P\) & \texttt{\detokenize{M,}\allowbreak\detokenize{N,}\allowbreak\detokenize{P}} & lambda-term metavariables & P7.3-P7.6 & n/a & P7.3 & \thepage & Terms of the locally fixed typed or untyped lambda grammar; the active calculus fixes the exact syntax. \\
\phantomsection\label{entry:symbol-registry-00405}\(\Theta_{\mathrm{ty}}\derives_{\mathsf{STLC}}M{:}\alpha\) & \texttt{\detokenize{\derives_{\mathsf{STLC}}}} & STLC typing judgment & P7.3-P7.4 & meta & P7.3 & \thepage & The exact arrow-only typing judgment of sheet:stlc; the context contains distinct term-variable declarations. \\
\phantomsection\label{entry:symbol-registry-00406}\(M[N/x]\) & \texttt{\detokenize{[N/}\allowbreak\detokenize{x]}} & capture-avoiding term substitution & P7.3-P7.6 & meta & P7.3 & \thepage & Substitutes term N for free occurrences of x in M after required alpha-renaming; distinct from formula substitution. \\
\phantomsection\label{entry:symbol-registry-00407}\(\lambda x.M\) & \texttt{\detokenize{\lambda}} & untyped lambda abstraction & P7.4 & 3 & P7.4 & \thepage & Term former binding x in M; separate exact-scope role from typed abstraction. \\
\phantomsection\label{entry:symbol-registry-00408}\(M\longrightarrow_{\beta}N,\ M\longrightarrow_{\beta}^{*}N\) & \texttt{\detokenize{\longrightarrow_{\beta},}\allowbreak\detokenize{\longrightarrow_{\beta}^{*}}} & one-step beta reduction and finite reflexive-transitive closure & P7.3-P7.4 & meta & P7.3 & \thepage & The active sheet fixes term beta reduction; the star denotes finite reflexive-transitive closure. \\
\phantomsection\label{entry:symbol-registry-00409}\(M=_{\beta}N\) & \texttt{\detokenize{=_{\beta}}} & beta conversion & P7.4 & meta & P7.4 & \thepage & Least congruence equivalence relation containing beta reduction; not elementary equivalence or formula equivalence. \\
\phantomsection\label{entry:symbol-registry-00410}\(M\mathrel{\mathsf{Par}}N\) & \texttt{\detokenize{\mathsf{Par}}} & parallel beta-reduction relation & P7.4 & meta & P7.4 & \thepage & Auxiliary relation on alpha-equivalence classes fixed by the four local clauses; its closure is compared with ordinary beta reduction. \\
\phantomsection\label{entry:symbol-registry-00411}\(\operatorname{dev}(M)\) & \texttt{\detokenize{\operatorname{dev}}} & complete development of a lambda term & P7.4 & n/a & P7.4 & \thepage & Metalinguistic recursion contracting all residuals exposed by the starting syntax; not a reduction relation. \\
\phantomsection\label{entry:symbol-registry-00412}\(\lambda f.\lambda x.f^n x\) & \texttt{\detokenize{\lambda f.\lambda x.f^n x}} & Church numeral encoding of n & P7.4 & n/a & P7.4 & \thepage & Closed untyped term used as the fixed numerical input-output encoding for lambda definability. \\
\phantomsection\label{entry:symbol-registry-00413}\(\mathsf S,\mathsf K,\mathsf{SK}\) & \texttt{\detokenize{\mathsf S,}\allowbreak\detokenize{\mathsf K,}\allowbreak\detokenize{\mathsf{SK}}} & S-K combinators and combinatory-calculus identifier & P7.4 & n/a & P7.4 & \thepage & Sheet:sk-combinatory fixes the two constants and exact calculus; scoped away from Post-clone and modal-system letters. \\
\phantomsection\label{entry:symbol-registry-00414}\(M\longrightarrow_{\mathsf{SK}}N\) & \texttt{\detokenize{\longrightarrow_{\mathsf{SK}}}} & one-step S-K contraction & P7.4 & meta & P7.4 & \thepage & Compatible contextual contraction fixed by sheet:sk-combinatory. \\
\phantomsection\label{entry:symbol-registry-00415}\(\Delta_{\mathrm{ty}}\) & \texttt{\detokenize{\Delta_{\mathrm{ty}}}} & System-F type-variable context & P7.4 & n/a & P7.4 & \thepage & Finite list of distinct type variables; distinct from formula and sequent contexts. \\
\phantomsection\label{entry:symbol-registry-00416}\(\mathsf F\) & \texttt{\detokenize{\mathsf F}} & Church-style System-F calculus identifier & P7.4 & n/a & P7.4 & \thepage & Exact identifier on sheet:system-f; its scope is disjoint from temporal eventually in P8.3. \\
\phantomsection\label{entry:symbol-registry-00417}\(\forall a.\alpha\) & \texttt{\detokenize{\forall}} & System-F universal type former & P7.4 & 3 & P7.4 & \thepage & Binds the type variable a in alpha; its syntactic category distinguishes it from the object-language formula quantifier. \\
\phantomsection\label{entry:symbol-registry-00418}\(\Lambda a.M\) & \texttt{\detokenize{\Lambda}} & System-F type abstraction & P7.4 & 3 & P7.4 & \thepage & Term former binding type variable a in M; distinct from ordinary term abstraction. \\
\phantomsection\label{entry:symbol-registry-00419}\(M[\alpha],\alpha[\beta/a],M[\alpha/a]\) & \texttt{\detokenize{[\alpha],}\allowbreak\detokenize{[\beta/}\allowbreak\detokenize{a],}\allowbreak\detokenize{[\alpha/}\allowbreak\detokenize{a]}} & System-F type application and capture-avoiding type substitution & P7.4 & meta & P7.4 & \thepage & The active syntactic category distinguishes term-level type application from substitution into a type or term. \\
\phantomsection\label{entry:symbol-registry-00420}\(\Delta_{\mathrm{ty}}\derives_{\mathsf F}\alpha\ \mathsf{type}\) & \texttt{\detokenize{\derives_{\mathsf F}}} & System-F type-formation judgment & P7.4 & meta & P7.4 & \thepage & Formation judgment for the least type set fixed by sheet:system-f. \\
\phantomsection\label{entry:symbol-registry-00421}\(\Delta_{\mathrm{ty}};\Theta_{\mathrm{ty}}\derives_{\mathsf F}M{:}\alpha\) & \texttt{\detokenize{\derives_{\mathsf F}}} & System-F term-typing judgment & P7.4 & meta & P7.4 & \thepage & Church-style typing judgment with separate type-variable and term-variable contexts. \\
\phantomsection\label{entry:symbol-registry-00422}\(\Gamma\;\mathsf{ctx}\) & \texttt{\detokenize{\mathsf{ctx}}} & MLTT context-formation judgment & P7.4 & meta & P7.4 & \thepage & Judgment of sheet:mltt-intensional asserting that the ordered dependent context Gamma is well formed. \\
\phantomsection\label{entry:symbol-registry-00423}\(\Gamma\vdash A\;\mathsf{type}\) & \texttt{\detokenize{\vdash,}\allowbreak\detokenize{\mathsf{type}}} & MLTT type-formation judgment & P7.4 & meta & P7.4 & \thepage & Type judgment in the selected intensional dependent calculus; distinct from System-F formation and formula derivability. \\
\phantomsection\label{entry:symbol-registry-00424}\(\Gamma\vdash a:A\) & \texttt{\detokenize{\vdash,}\allowbreak\detokenize{:}} & MLTT term-typing judgment & P7.4 & meta & P7.4 & \thepage & Dependent typing judgment on sheet:mltt-intensional; the colon is typed ascription rather than a labelled modal formula. \\
\phantomsection\label{entry:symbol-registry-00425}\(\Gamma\vdash A\equiv B\;\mathsf{type},\ \Gamma\vdash a\equiv b:A\) & \texttt{\detokenize{\equiv}} & MLTT judgmental type and term equality & P7.4 & meta & P7.4 & \thepage & Typed definitional equality generated by the displayed computations and congruence; equality reflection into identity types is absent. \\
\phantomsection\label{entry:symbol-registry-00426}\(\Pi_{x:A}B,\Sigma_{x:A}B\) & \texttt{\detokenize{\Pi,}\allowbreak\detokenize{\Sigma}} & MLTT dependent product and dependent sum type formers & P7.4 & 3 & P7.4 & \thepage & Bind x in B on sheet:mltt-intensional; distinct from analytical classes and sequent-context metavariables. \\
\phantomsection\label{entry:symbol-registry-00427}\(\operatorname{Id}_A(a,b)\) & \texttt{\detokenize{\operatorname{Id}}} & MLTT intensional identity type & P7.4 & n/a & P7.4 & \thepage & Proposition-as-type of identifications between terms a and b of A; no judgmental equality reflection is assumed. \\
\phantomsection\label{entry:symbol-registry-00428}\(\mathcal U_i\) & \texttt{\detokenize{\mathcal U_i}} & MLTT Russell universe at level i & P7.4 & n/a & P7.4 & \thepage & Cumulative predicative universe hierarchy on sheet:mltt-intensional; distinct from the P6 measurable-cardinal ultrafilter role of bare calligraphic U. \\
\phantomsection\label{entry:symbol-registry-00429}\(J_C(d,a,b,p)\) & \texttt{\detokenize{J_C}} & MLTT identity eliminator & P7.4 & n/a & P7.4 & \thepage & Typed eliminator for the displayed identity motive C with its judgmental refl computation rule. \\
\phantomsection\label{entry:symbol-registry-00430}\(\sigma,\tau:\Delta\to\Gamma\) & \texttt{\detokenize{\sigma,}\allowbreak\detokenize{\tau}} & MLTT substitutions with common domain and codomain & P7.4 & meta & P7.4 & \thepage & Capture-avoiding substitutions used by typed constructor congruence; equality premises and conclusions live in Delta. \\
\phantomsection\label{entry:symbol-registry-00431}\(\operatorname{Eq}_{\Delta}(J,J')\) & \texttt{\detokenize{\operatorname{Eq}_{\Delta}}} & multi-sorted MLTT equality premise & P7.4 & meta & P7.4 & \thepage & Eliminable metalinguistic abbreviation for type or term judgmental equality after transport to the common context Delta. \\
\phantomsection\label{entry:symbol-registry-00432}\(\mathsf 1_{\mathrm{ty}},\mathsf 0_{\mathrm{ty}}\) & \texttt{\detokenize{\mathsf 1_{\mathrm{ty}},}\allowbreak\detokenize{\mathsf 0_{\mathrm{ty}}}} & unit and empty types in the categorical STLC extension & P7.3 & n/a & P7.3 & \thepage & Typed subscripts distinguish these formers from numerical constants and the computable Turing degree. \\
\phantomsection\label{entry:symbol-registry-00433}\(\alpha\mathbin{\times_{\mathrm{ty}}}\beta,\alpha\mathbin{+_{\mathrm{ty}}}\beta\) & \texttt{\detokenize{\times_{\mathrm{ty}},}\allowbreak\detokenize{+_{\mathrm{ty}}}} & product and sum type formers in the categorical STLC extension & P7.3 & 3 & P7.3 & \thepage & Exact added type formers used to construct finite products and coproducts; distinct from arithmetic operations. \\
\phantomsection\label{entry:symbol-registry-00434}\(\star,\langle M,N\rangle,\pi_iM,\mathsf{inl}(M),\mathsf{inr}(M)\) & \texttt{\detokenize{\star,}\allowbreak\detokenize{\langle\ ,}\allowbreak\detokenize{\ \rangle,}\allowbreak\detokenize{\pi_i,}\allowbreak\detokenize{\mathsf{inl},}\allowbreak\detokenize{\mathsf{inr}}} & unit product and sum constructors for the categorical STLC extension & P7.3 & n/a & P7.3 & \thepage & Term constructors governed by the exact typing table and beta-eta equations in the Curry-Howard entry. \\
\phantomsection\label{entry:symbol-registry-00435}\(\mathsf{case}(L;x.M;y.N),\mathsf{abort}_{\alpha}(M)\) & \texttt{\detokenize{\mathsf{case},}\allowbreak\detokenize{\mathsf{abort}_{\alpha}}} & sum and empty-type eliminators for the categorical STLC extension & P7.3 & n/a & P7.3 & \thepage & Typed eliminators used only by the explicitly defined bicartesian extension. \\
\phantomsection\label{entry:symbol-registry-00436}\(r(A),h(c),o(D)\) & \texttt{\detokenize{r,}\allowbreak\detokenize{h,}\allowbreak\detokenize{o}} & formula rank Cut height and derivation ordinal measure & P7.5 & meta & P7.5 & \thepage & Metalinguistic measures for the finite-set propositional Cut strategy; o(D) is the Hessenberg natural sum displayed in its definition. \\
\phantomsection\label{entry:symbol-registry-00437}\(\operatorname{Cut}_A(D_0,D_1),\operatorname{Cut}(D),A_c\) & \texttt{\detokenize{\operatorname{Cut}_A,}\allowbreak\detokenize{\operatorname{Cut},}\allowbreak\detokenize{A_c}} & Cut node Cut-node set and node-indexed cut formula & P7.5 & meta & P7.5 & \thepage & Metalinguistic notation for one Cut with premises D0 D1, the finite set of Cut nodes of D, and the formula attached to node c in the ordinal assignment. \\
\phantomsection\label{entry:symbol-registry-00438}\(\mathop{\#}\) & \texttt{\detokenize{\#}} & Hessenberg natural sum & P7.5 & meta & P7.5 & \thepage & Commutative ordinal sum used only to combine the finitely many Cut-node contributions in o(D). \\
\phantomsection\label{entry:symbol-registry-00440}\(\mathsf{HA}\) & \texttt{\detokenize{\mathsf{HA}}} & Heyting arithmetic & P7.5-P7.6 & n/a & P7.5 & \thepage & First-order arithmetic with full induction over intuitionistic logic; not a semantic notion of constructive truth. \\
\phantomsection\label{entry:symbol-registry-00450}\(\alpha,\beta\) & \texttt{\detokenize{\alpha,}\allowbreak\detokenize{\beta}} & ordinal metavariables & P7.5 & n/a & P7.5 & \thepage & Ordinals in the local proof-theoretic calibration; disjoint from the P7.3-P7.4 and P7.6 type scopes. \\
\phantomsection\label{entry:symbol-registry-00451}\(\gamma\) & \texttt{\detokenize{\gamma}} & ordinal metavariable & P7.5 & n/a & P7.5 & \thepage & Ordinal variable in the Veblen hierarchy and the definition of Gamma-zero; locally disjoint from formula and type metavariables. \\
\phantomsection\label{entry:symbol-registry-00452}\(\omega^\omega\) & \texttt{\detokenize{\omega^\omega}} & ordinal exponentiation & P7.5 & n/a & P7.5 & \thepage & The ordinal power used in the PRA calibration; disjoint from the P6 Baire-space role. \\
\phantomsection\label{entry:symbol-registry-00453}\(e\mapsto(D_e,\prec_e)\) & \texttt{\detokenize{\mapsto}} & effective notation-system enumeration & P7.5 & meta & P7.5 & \thepage & Maps each natural-number code to its recursive domain and ordering relation in the relative ordinal calibration. \\
\phantomsection\label{entry:symbol-registry-00454}\(\mathcal C\) & \texttt{\detokenize{\mathcal C}} & formula class for transfinite induction & P7.5 & n/a & P7.5 & \thepage & Locally fixed family of formulas parameterizing the proof-theoretic calibration; disjoint from the P4 class scope. \\
\phantomsection\label{entry:symbol-registry-00455}\(\Delta^0_1\) & \texttt{\detokenize{\Delta^0_1}} & Delta-zero-one arithmetic formula class & P7.5 & n/a & P7.5 & \thepage & Formulas both Sigma-zero-one and Pi-zero-one under the fixed arithmetic syntax; used in RCA-zero comprehension. \\
\phantomsection\label{entry:symbol-registry-00456}\(\omega_n\) & \texttt{\detokenize{\omega_n}} & finite tower sequence below epsilon-zero & P7.5 & n/a & P7.5 & \thepage & Defined by omega-zero equals one and omega-(k+1) equals omega to the omega-k; used in the induction-fragment calibration. \\
\phantomsection\label{entry:symbol-registry-00457}\(\Gamma_0\) & \texttt{\detokenize{\Gamma_0}} & Feferman--Schuette ordinal & P7.5 & n/a & P7.5 & \thepage & Least positive gamma satisfying Veb-gamma of zero equals gamma under the locally defined Veblen hierarchy. \\
\phantomsection\label{entry:symbol-registry-00458}\(\operatorname{Veb}_{\alpha}(\beta)\) & \texttt{\detokenize{\operatorname{Veb}}} & Veblen hierarchy of normal ordinal functions & P7.5 & n/a & P7.5 & \thepage & Veb-zero is ordinal exponentiation by omega; later levels enumerate common fixed points of all earlier levels. \\
\phantomsection\label{entry:symbol-registry-00459}\(\alpha,\beta\) & \texttt{\detokenize{\alpha,}\allowbreak\detokenize{\beta}} & finite-type metavariables & P7.6 & n/a & P7.6 & \thepage & Finite types generated from ground type zero by the registered type arrow; disjoint from P7.5 ordinals. \\
\phantomsection\label{entry:symbol-registry-00460}\(\mathsf T_{\mathrm{tm}}\) & \texttt{\detokenize{\mathsf T_{\mathrm{tm}}}} & System-T term calculus identifier & P7.6 & n/a & P7.6 & \thepage & Exact finite-type term calculus fixed by sheet:system-t. \\
\phantomsection\label{entry:symbol-registry-00461}\(\mathsf T_{\mathrm{qf}}\) & \texttt{\detokenize{\mathsf T_{\mathrm{qf}}}} & quantifier-free System-T target theory & P7.6 & n/a & P7.6 & \thepage & Separately named equality theory used as the Dialectica matrix target; not the term calculus alone. \\
\phantomsection\label{entry:symbol-registry-00462}\(\Theta_{\mathrm{ty}}\derives_{\mathsf T_{\mathrm{tm}}}M{:}\alpha\) & \texttt{\detokenize{\derives_{\mathsf T_{\mathrm{tm}}}}} & System-T typing judgment & P7.6 & meta & P7.6 & \thepage & Typing judgment of sheet:system-t with the exact finite-type term calculus in the turnstile subscript. \\
\phantomsection\label{entry:symbol-registry-00463}\(0\) & \texttt{\detokenize{0}} & System-T ground type and zero constant & P7.6 & n/a & P7.6 & \thepage & Syntactic category distinguishes the ground type from the term constant; both are fixed by sheet:system-t. \\
\phantomsection\label{entry:symbol-registry-00464}\(S\) & \texttt{\detokenize{S}} & System-T successor constant & P7.6 & n/a & P7.6 & \thepage & Ground-type successor constant on sheet:system-t; scoped away from formula and Post-clone letters. \\
\phantomsection\label{entry:symbol-registry-00465}\(R_\alpha\) & \texttt{\detokenize{R_\alpha}} & System-T finite-type recursor & P7.6 & n/a & P7.6 & \thepage & One primitive recursor at each finite type with argument order and computation equations fixed by sheet:system-t. \\
\phantomsection\label{entry:symbol-registry-00466}\(M\longrightarrow_{\mathsf T_{\mathrm{tm}}}N\) & \texttt{\detokenize{\longrightarrow_{\mathsf T_{\mathrm{tm}}}}} & one-step System-T term reduction & P7.6 & meta & P7.6 & \thepage & Compatible reduction generated by beta contraction and the two registered recursor computation rules. \\
\phantomsection\label{entry:symbol-registry-00467}\(\varphi^{\mathrm D}\) & \texttt{\detokenize{{}^{\mathrm D}}} & Dialectica translation of a formula & P7.6 & n/a & P7.6 & \thepage & Translation into an existential witness block, universal challenge block, and quantifier-free target matrix. \\
\phantomsection\label{entry:symbol-registry-00468}\(\varphi_{\mathrm D}(\vec x,\vec y)\) & \texttt{\detokenize{{}_{\mathrm D}}} & Dialectica quantifier-free matrix & P7.6 & n/a & P7.6 & \thepage & Matrix associated with the locally fixed Dialectica translation; typed witness and challenge tuples remain explicit. \\
\phantomsection\label{entry:symbol-registry-00469}\(\vec x,\vec y,\vec u,\vec v\) & \texttt{\detokenize{\vec x,}\allowbreak\detokenize{\vec y,}\allowbreak\detokenize{\vec u,}\allowbreak\detokenize{\vec v}} & Dialectica witness and challenge tuples & P7.6 & n/a & P7.6 & \thepage & Finite tuples of finite-type variables; distinct from the P4-P5 tuples of numerical variables. \\
\phantomsection\label{entry:symbol-registry-00470}\(\vec f,\vec g,\vec t\) & \texttt{\detokenize{\vec f,}\allowbreak\detokenize{\vec g,}\allowbreak\detokenize{\vec t}} & finite tuples of System-T functionals and terms & P7.6 & n/a & P7.6 & \thepage & Typed functional witnesses in translation clauses and the extracted closed-term tuple in Dialectica soundness. \\
\phantomsection\label{entry:symbol-registry-00471}\(\pi_0,\pi_1\) & \texttt{\detokenize{\pi_0,}\allowbreak\detokenize{\pi_1}} & primitive-recursive pairing projections & P7.6 & n/a & P7.6 & \thepage & First and second projections for the fixed numerical pairing function used in number realizability. \\
\phantomsection\label{entry:symbol-registry-00472}\(\mathsf{ECT}_0\) & \texttt{\detokenize{\mathsf{ECT}_0}} & extended Church-thesis scheme & P7.6 & n/a & P7.6 & \thepage & Universal closure of the displayed almost-negative-domain scheme; not an unrestricted Church-thesis axiom. \\
\phantomsection\label{entry:symbol-registry-00473}\(\mathsf{PC}_{\varepsilon}\) & \texttt{\detokenize{\mathsf{PC}_{\varepsilon}}} & equality-free epsilon-calculus identifier & P7.7 & n/a & P7.7 & \thepage & Exact Hilbert-FOL extension fixed by sheet:epsilon; primitive quantifiers are retained and equality is excluded. \\
\phantomsection\label{entry:symbol-registry-00474}\(\mathsf{EC}\) & \texttt{\detokenize{\mathsf{EC}}} & epsilon-free quantifier-free fragment & P7.7 & n/a & P7.7 & \thepage & Eliminable metalinguistic abbreviation used in the first epsilon theorem; not a second epsilon system. \\
\phantomsection\label{entry:symbol-registry-00475}\(\varepsilon x\,\varphi(x)\) & \texttt{\detokenize{\varepsilon}} & Hilbert epsilon term former & P7.7 & 3 & P7.7 & \thepage & Capture-avoiding term binder on the equality-free epsilon sheet; distinct from epsilon-zero. \\
\phantomsection\label{entry:symbol-registry-00476}\(w{:}\varphi\) & \texttt{\detokenize{:}} & world-labelled formula & P7.7 & meta & P7.7 & \thepage & Labelled-sequent syntax only; distinct from type ascription and not a semantic satisfaction assertion. \\
\phantomsection\label{entry:symbol-registry-00477}\(wRv\) & \texttt{\detokenize{R}} & relational atom in labelled-sequent syntax & P7.7 & n/a & P7.7 & \thepage & Syntactic frame-relation atom inside a labelled derivation language; distinct from semantic accessibility in P8. \\
\phantomsection\label{entry:symbol-registry-00478}\(\mathsf P_{\mathrm{pf}},\mathsf Q_{\mathrm{pf}}\) & \texttt{\detokenize{\mathsf P_{\mathrm{pf}},}\allowbreak\detokenize{\mathsf Q_{\mathrm{pf}}}} & Cook-Reckhow proof-system metavariables & P7.8 & n/a & P7.8 & \thepage & Polynomial-time encoded surjections onto TAUT; subscripts prevent collision with complexity and formula letters. \\
\phantomsection\label{entry:symbol-registry-00479}\(\operatorname{Prf}_{\mathsf P_{\mathrm{pf}}}(\pi,\varphi)\) & \texttt{\detokenize{\operatorname{Prf}_{\mathsf P_{\mathrm{pf}}}}} & encoded propositional proof relation & P7.8 & meta & P7.8 & \thepage & Polynomial-time decidable relation whose range is exactly TAUT under the fixed malformed-input convention. \\
\phantomsection\label{entry:symbol-registry-00480}\(\pi\) & \texttt{\detokenize{\pi}} & encoded propositional proof metavariable & P7.8 & n/a & P7.8 & \thepage & Binary proof string for the fixed Cook-Reckhow encoding; no standing as a projection or permutation. \\
\phantomsection\label{entry:symbol-registry-00481}\(\operatorname{size}(a)\) & \texttt{\detokenize{\operatorname{size}}} & bit length of an encoded proof or formula object & P7.8 & n/a & P7.8 & \thepage & Defined from the fixed binary encoding; deliberately not vertical-bar cardinality notation. \\
\phantomsection\label{entry:symbol-registry-00482}\(\mathsf{PHP}^{n}_{n+1},x_{i,j}\) & \texttt{\detokenize{\mathsf{PHP}^{n}_{n+1},}\allowbreak\detokenize{x_{i,j}}} & pigeonhole-principle CNF family and its variables & P7.8 & n/a & P7.8 & \thepage & Exact weak PHP encoding with n+1 pigeons and n holes; no clauses forbid one pigeon from occupying two holes. \\
\phantomsection\label{entry:symbol-registry-00483}\(\longrightarrow\) & \texttt{\detokenize{\longrightarrow}} & proof-system function arrow & P7.8 & meta & P7.8 & \thepage & Metalinguistic total-function arrow in the Cook-Reckhow definition; disjoint from the P7.3 abstract-reduction scope. \\
\end{longtable}
\endgroup
\Needspace{.4\textheight}
\subsection*{Part 8 scopes}
\begingroup\scriptsize\setlength{\tabcolsep}{3pt}
\begin{longtable}{@{}p{0.9in}p{1.3in}p{1.0in}p{0.5in}p{0.3in}p{0.4in}p{0.35in}p{1.65in}@{}}
\toprule
\textbf{Glyph} & \textbf{Macro} & \textbf{Role} & \textbf{Scope} & \textbf{Prec} & \textbf{First} & \textbf{Reg. page} & \textbf{Principal uses / notes} \\
\midrule
\endfirsthead
\toprule
\textbf{Glyph} & \textbf{Macro} & \textbf{Role} & \textbf{Scope} & \textbf{Prec} & \textbf{First} & \textbf{Reg. page} & \textbf{Principal uses / notes} \\
\midrule
\endhead
\midrule\multicolumn{8}{r}{\emph{continued}}\\\endfoot
\bottomrule\endlastfoot
\phantomsection\label{entry:symbol-registry-00090}\(w\Vdash A\) & \texttt{\detokenize{\forces}} & Kripke forcing & P8.1 & meta & P8.1 & \thepage & The negated partner is \(\nVdash\) (macro \texttt{\detokenize{\nforces}}); modal satisfaction continues to use \(\models\). \\
\phantomsection\label{entry:symbol-registry-00092}\(\Box A\) & \texttt{\detokenize{\Box}} & Anderson--Belnap relevance abbreviation & P8.6 & 3 & P8.6 & \thepage & Exact eliminable abbreviation \((A\to A)\to A\) on the E and R sheets; it is not a primitive modality. \\
\phantomsection\label{entry:symbol-registry-00094}\(\mathsf X\) & \texttt{\detokenize{\mathsf X}} & temporal next & P8.3-P9.4 & 3 & P8.3 & \thepage & Prefix unary at the negation rank. \\
\phantomsection\label{entry:symbol-registry-00095}\(A\mathbin{\mathsf U}B\) & \texttt{\detokenize{\mathbin{\mathsf U}}} & temporal until & P8.3-P9.5 & 3.5 & P8.3 & \thepage & Binds tighter than every binary connective and looser than the unary rank; licenses the U-unfolding axiom's written form. \\
\phantomsection\label{entry:symbol-registry-00096}\(\mathsf F\) & \texttt{\detokenize{\mathsf F}} & temporal eventually & P8.3-P9.5 & 3 & P8.3 & \thepage & Defined \(\mathsf FA:\equiv\top\mathbin{\mathsf U}A\), prefix unary at the negation rank. \\
\phantomsection\label{entry:symbol-registry-00097}\(\mathsf G\) & \texttt{\detokenize{\mathsf G}} & temporal always & P8.3-P9.5 & 3 & P8.3 & \thepage & Defined \(\mathsf GA:\equiv\neg\mathsf F\neg A\), prefix unary at the negation rank. \\
\phantomsection\label{entry:symbol-registry-00098}\(K_i\) & \texttt{\detokenize{K_i}} & knowledge operator (agent i) & P8.3 & 3 & P8.3 & \thepage & Prefix unary at the negation rank. \\
\phantomsection\label{entry:symbol-registry-00099}\(\sim_i\) & \texttt{\detokenize{\sim_i}} & epistemic accessibility & P8.3 & n/a & P8.3 & \thepage & Agent-indexed accessibility; an equivalence relation for knowledge readings. \\
\phantomsection\label{entry:symbol-registry-00100}\(E_G\) & \texttt{\detokenize{E_G}} & mutual knowledge & P8.3 & 3 & P8.3 & \thepage & \(E_GA:\equiv\bigwedge_{i\in G}K_iA\), prefix unary at the negation rank. \\
\phantomsection\label{entry:symbol-registry-00101}\(C_G\) & \texttt{\detokenize{C_G}} & common knowledge & P8.3 & 3 & P8.3 & \thepage & Prefix unary at the negation rank; satisfaction via reflexive-transitive reachability under the union of the \(\sim_i\). \\
\phantomsection\label{entry:symbol-registry-00102}\(\otimes\) & \texttt{\detokenize{\otimes}} & linear tensor and its homomorphic semantic interpretation & P8.6 & 4 & P8.6 & \thepage & Formula operands give the exact IMLL or CLL connective; the IMLL algebra and CLL phase clauses interpret that same connective with their source-standard tensor operation. \\
\phantomsection\label{entry:symbol-registry-00103}\(\multimap\) & \texttt{\detokenize{\multimap}} & linear implication and its residual interpretation & P8.6 & 6 & P8.6 & \thepage & Formula operands give IMLL or CLL linear implication; the algebraic and phase clauses interpret that same connective by the source-standard residual. \\
\phantomsection\label{entry:symbol-registry-00104}\(\mathbf 1\) & \texttt{\detokenize{\mathbf 1}} & multiplicative unit & P8.6 & n/a & P8.6 & \thepage & Syntactic tensor unit in IMLL and CLL; the ambient calculus distinguishes it from a semantic monoid identity. \\
\phantomsection\label{entry:symbol-registry-00105}\(!A\) & \texttt{\detokenize{!}} & bang exponential and its phase interpretation & P8.6 & 2.5 & P8.6 & \thepage & Prefix connective on the source-bounded IMLL extension and full CLL; the phase clause for \(!F\) interprets that same connective, and neither use is factorial. \\
\phantomsection\label{entry:symbol-registry-00106}\(W\) & \texttt{\detokenize{W}} & set of worlds & P8 & n/a & P8.1 & \thepage & Nonempty world carrier for intuitionistic, modal, temporal, epistemic, and broad Routley--Meyer semantics; world letters are declared locally. \\
\phantomsection\label{entry:symbol-registry-00107}\(R\) & \texttt{\detokenize{R}} & binary accessibility relation & P8.2-P8.3 & n/a & P8.2 & \thepage & Frame relation of modal, temporal, and epistemic structures; scope is disjoint from the P8.6 ternary relation. \\
\phantomsection\label{entry:symbol-registry-00108}\(N\) & \texttt{\detokenize{N}} & normal worlds & P8.6 & n/a & P8.6 & \thepage & Nonempty distinguished subset of the broad inherited Routley--Meyer structure fixing its source-scoped validity notion. \\
\phantomsection\label{entry:symbol-registry-00109}\(R\) & \texttt{\detokenize{R}} & ternary accessibility relation & P8.6 & n/a & P8.6 & \thepage & Ternary relation in broad and matched Routley--Meyer structures; scope is disjoint from binary accessibility. \\
\phantomsection\label{entry:symbol-registry-00110}\(w^*\) & \texttt{\detokenize{^{*}}} & Routley star & P8.6 & n/a & P8.6 & \thepage & Postfix unary operation on relevance worlds used by the negation clause; distinct from closure stars and the extension marker in \(\mathsf L^*\). \\
\phantomsection\label{entry:symbol-registry-00111}\(\mathsf V\) & \texttt{\detokenize{\mathsf V}} & matrix value set & P8.5 & n/a & P8.5 & \thepage & Finite nonempty truth-value set for the generic matrix sheet; distinct from world valuation \(V\) and cumulative \(V_\alpha\). \\
\phantomsection\label{entry:symbol-registry-00112}\(D\) & \texttt{\detokenize{D}} & designated values & P8.5 & n/a & P8.5 & \thepage & Subset of the active matrix value set fixing matrix consequence. \\
\phantomsection\label{entry:symbol-registry-00113}\(f_c\) & \texttt{\detokenize{f_c}} & matrix operation for connective c & P8.5 & n/a & P8.5 & \thepage & Total n-ary operation on \(\mathsf V\) interpreting the declared connective \(c\). \\
\phantomsection\label{entry:symbol-registry-00115}\(\models^{\ell}\) & \texttt{\detokenize{\entails^{\ell}}} & local modal consequence & P8.2-P8.3 & meta & P8.2 & \thepage & Truth preservation at one world in the declared modal or epistemic frame class. \\
\phantomsection\label{entry:symbol-registry-00116}\(\models^{g}\) & \texttt{\detokenize{\entails^{g}}} & global modal consequence & P8.2-P8.3 & meta & P8.2 & \thepage & Truth preservation at all worlds in the declared modal or epistemic frame class. \\
\phantomsection\label{entry:symbol-registry-00117}\(\mathcal M,w\models A\) & \texttt{\detokenize{\entails}} & satisfaction at a world & P8 & meta & P8.2 & \thepage & Satisfaction at world \(w\) of the structure \(\mathcal M\). \\
\phantomsection\label{entry:symbol-registry-00118}\(A,B,C\) & \texttt{\detokenize{A,}\allowbreak\detokenize{B,}\allowbreak\detokenize{C}} & default formula metavariables in inherited Part-8 prose & P8 & n/a & P8.1 & \thepage & Default only where the inherited or cited presentation uses these letters.  A named source presentation's own locally declared syntax overrides this row; this row never licenses renaming source metavariables. \\
\phantomsection\label{entry:symbol-registry-00119}\(\mathcal A=(A,\leq,\otimes,1,\multimap)\) & \texttt{\detokenize{\mathcal A,}\allowbreak\detokenize{A,}\allowbreak\detokenize{\leq,}\allowbreak\detokenize{\otimes,}\allowbreak\detokenize{1,}\allowbreak\detokenize{\multimap}} & residuated pomonoid datum & P8.6 & n/a & P8.6 & \thepage & The inherited source tuple.  Within its definition A is locally typed as the algebra carrier, overriding the default formula-metavariable role. \\
\phantomsection\label{entry:symbol-registry-00120}\(V\) & \texttt{\detokenize{V}} & world valuation & P8 & n/a & P8.1 & \thepage & \(V:\mathsf{Prop}\to\mathcal P(W)\) (temporal structures: \(V:\mathbb N\to\mathcal P(\mathsf{Prop})\)); distinct from matrix \(\mathsf V\) and cumulative \(V_\alpha\). \\
\phantomsection\label{entry:symbol-registry-00124}\(\mathsf{Prop}\) & \texttt{\detokenize{\mathsf{Prop}}} & default propositional-atom set & P8 & n/a & P8.1 & \thepage & Used where the inherited or cited presentation declares this atom-set glyph.  A named source presentation may instead declare its own atom notation, which governs locally without house renaming. \\
\phantomsection\label{entry:symbol-registry-00484}\(\mathsf{IPL},\mathsf{KC},\mathsf{LC}\) & \texttt{\detokenize{\mathsf{IPL},}\allowbreak\detokenize{\mathsf{KC},}\allowbreak\detokenize{\mathsf{LC}}} & intuitionistic and intermediate logic identifiers & P8.1 & n/a & P8.1 & \thepage & Exact IPL card and the named KC and Goedel--Dummett LC extensions; IQC is omitted because no IQC identifier is printed. \\
\phantomsection\label{entry:symbol-registry-00485}\(\Gamma\derives_{\mathsf{IPL}}A\) & \texttt{\detokenize{\derives_{\mathsf{IPL}}}} & finite-support IPL derivability & P8.1-P8.2 & meta & P8.1 & \thepage & A premise-set judgment witnessed by a finite multiset subcontext derivable on sheet:lj; P8.2 reuses it only in the exact IPL-to-S4 embedding theorem. \\
\phantomsection\label{entry:symbol-registry-00486}\(\Gamma\entails_{\mathsf{IPL}}A\) & \texttt{\detokenize{\entails_{\mathsf{IPL}}}} & IPL Kripke consequence & P8.1 & meta & P8.1 & \thepage & Truth preservation at one world in every intuitionistic Kripke structure of the declared model class under all premises in Gamma. \\
\phantomsection\label{entry:symbol-registry-00487}\(\Gamma_0\) & \texttt{\detokenize{\Gamma_0}} & finite-support witness context & P8.1 & n/a & P8.1 & \thepage & Finite multiset whose members lie in the ambient IPL premise set; disjoint from the P7.5 Feferman--Schuette ordinal. \\
\phantomsection\label{entry:symbol-registry-00488}\(G(A),A^{\mathrm{Ko}},A^{\mathrm N},A_{\mathrm{Ku}},A^{\mathrm{Ku}}\) & \texttt{\detokenize{G,}\allowbreak\detokenize{A^{\mathrm{Ko}},}\allowbreak\detokenize{A^{\mathrm N},}\allowbreak\detokenize{A_{\mathrm{Ku}},}\allowbreak\detokenize{A^{\mathrm{Ku}}}} & negative-translation maps & P8.1 & meta & P8.1 & \thepage & Exact Glivenko Kolmogorov Goedel--Gentzen and inner/full Kuroda maps; superscripts act memberwise on premise contexts when printed there. \\
\phantomsection\label{entry:symbol-registry-00489}\(\mathcal H=(H,\wedge,\vee,0,1,\Rightarrow)\) & \texttt{\detokenize{\mathcal H,}\allowbreak\detokenize{H,}\allowbreak\detokenize{\wedge,}\allowbreak\detokenize{\vee,}\allowbreak\detokenize{0,}\allowbreak\detokenize{1,}\allowbreak\detokenize{\Rightarrow}} & Heyting algebra datum & P8.1 & n/a & P8.1 & \thepage & The cited presentation's bounded distributive lattice with residual.  Local element operands type these glyphs as algebraic operations rather than object-language connectives or a sequent arrow. \\
\phantomsection\label{entry:symbol-registry-00490}\(a\wedge b,a\vee b,0,1,a\Rightarrow b,\neg a\) & \texttt{\detokenize{\wedge,}\allowbreak\detokenize{\vee,}\allowbreak\detokenize{0,}\allowbreak\detokenize{1,}\allowbreak\detokenize{\Rightarrow,}\allowbreak\detokenize{\neg}} & Heyting operations and bounds & P8.1 & n/a & P8.1 & \thepage & Source operations on elements of the locally declared Heyting algebra; the operand type, not a house subscript, distinguishes them from formula connectives. \\
\phantomsection\label{entry:symbol-registry-00491}\(\nu\) & \texttt{\detokenize{\nu}} & Heyting or open-set formula valuation & P8.1 & n/a & P8.1 & \thepage & Local semantic valuation extended homomorphically from propositional atoms; disjoint from the P6.3 forcing-name metavariable. \\
\phantomsection\label{entry:symbol-registry-00492}\(\operatorname{int}_{X}\) & \texttt{\detokenize{\operatorname{int}_{X}}} & topological interior operator & P8.1 & meta & P8.1 & \thepage & Interior in the declared topological space X used by the open-set implication clause. \\
\phantomsection\label{entry:symbol-registry-00493}\(D_w\) & \texttt{\detokenize{D_w}} & world-indexed intuitionistic domain & P8.1 & n/a & P8.1 & \thepage & Nonempty domain at world w in expanding-domain IQC semantics; monotone along the intuitionistic preorder. \\
\phantomsection\label{entry:symbol-registry-00494}\(K,KD,KT,KTB,KB\) & \texttt{\detokenize{K,}\allowbreak\detokenize{KD,}\allowbreak\detokenize{KT,}\allowbreak\detokenize{KTB,}\allowbreak\detokenize{KB}} & Chellas normal modal logic identifiers & P8.2 & n/a & P8.2 & \thepage & Chellas writes KD KT and KTB for the serial reflexive and reflexive-symmetric extensions and writes KB for the symmetry-only extension.  The familiar D T and B system names occur only as Hughes--Cresswell aliases on the corresponding cards. \\
\phantomsection\label{entry:symbol-registry-00495}\(K4,KT4,S4,KD45\) & \texttt{\detokenize{K4,}\allowbreak\detokenize{KT4,}\allowbreak\detokenize{S4,}\allowbreak\detokenize{KD45}} & Chellas normal modal logic identifiers & P8.2 & n/a & P8.2 & \thepage & Chellas's exact plain math-italic names; S4 is the source-named system KT4. \\
\phantomsection\label{entry:symbol-registry-00496}\(S4.2,S4.3,S5\) & \texttt{\detokenize{S4.2,}\allowbreak\detokenize{S4.3,}\allowbreak\detokenize{S5}} & Chellas normal modal logic identifiers & P8.2 & n/a & P8.2 & \thepage & Chellas's exact plain math-italic names for the displayed extensions. \\
\phantomsection\label{entry:symbol-registry-00497}\(\mathsf{GL},\mathsf{Grz}\) & \texttt{\detokenize{\mathsf{GL},}\allowbreak\detokenize{\mathsf{Grz}}} & normal modal logic identifiers & P8.2 & n/a & P8.2 & \thepage & Exact named provability and Grzegorczyk systems outside the ordinary modal-cube route. \\
\phantomsection\label{entry:symbol-registry-00498}\(E,EM,EC,EN,EMC,EMCN\) & \texttt{\detokenize{E,}\allowbreak\detokenize{EM,}\allowbreak\detokenize{EC,}\allowbreak\detokenize{EN,}\allowbreak\detokenize{EMC,}\allowbreak\detokenize{EMCN}} & Chellas compositional modal logic identifiers & P8.2 & n/a & P8.2 & \thepage & Chellas's exact plain math-italic compositional system names.  M C and N remain schema labels rather than substitute system identifiers; normal K is EMCN. \\
\phantomsection\label{entry:symbol-registry-00499}\(\mathfrak L\) & \texttt{\detokenize{\mathfrak L}} & normal modal logic metavariable & P8.2 & n/a & P8.2 & \thepage & Locally fixed normal modal logic in derivability canonical-model and frame-class constructions; disjoint from the P4.6 capture-logic use. \\
\phantomsection\label{entry:symbol-registry-00500}\(\mathcal C\) & \texttt{\detokenize{\mathcal C}} & modal frame-class metavariable & P8.2 & n/a & P8.2 & \thepage & Locally declared class of modal frames; disjoint from the P4 set or language class and the P7 formula-class use. \\
\phantomsection\label{entry:symbol-registry-00501}\(\operatorname{ST}_x,P_p\) & \texttt{\detokenize{\operatorname{ST}_x,}\allowbreak\detokenize{P_p}} & modal standard translation and atom predicate & P8.2 & meta & P8.2 & \thepage & ST at free world variable x maps each modal atom p to the corresponding unary first-order predicate \(P_p\). \\
\phantomsection\label{entry:symbol-registry-00502}\(A^{\Box}\) & \texttt{\detokenize{{}^{\Box}}} & Goedel--McKinsey--Tarski modal translation & P8.2 & meta & P8.2 & \thepage & Recursive IPL-to-S4 formula translation; distinct from applying the object-language box to an already translated formula. \\
\phantomsection\label{entry:symbol-registry-00503}\(Z\) & \texttt{\detokenize{Z}} & modal bisimulation relation & P8.2 & meta & P8.2 & \thepage & Relation between worlds of two pointed Kripke structures satisfying atom zig and zag clauses; disjoint from P8.3 fixpoint variables. \\
\phantomsection\label{entry:symbol-registry-00504}\(\Phi,\sim_{\Phi},[w]_{\Phi}\) & \texttt{\detokenize{\Phi,}\allowbreak\detokenize{\sim_{\Phi},}\allowbreak\detokenize{[\cdot]_{\Phi}}} & filtration set equivalence and quotient world & P8.2 & meta & P8.2 & \thepage & Finite subformula-closed set Phi induces modal equivalence and its world classes in the declared filtration. \\
\phantomsection\label{entry:symbol-registry-00505}\(N:W\to\mathcal P(\mathcal P(W))\) & \texttt{\detokenize{N}} & neighborhood function & P8.2 & n/a & P8.2 & \thepage & The source's world-indexed neighborhood assignment.  Its function type distinguishes it from system EN and the later Routley--Meyer set of normal worlds. \\
\phantomsection\label{entry:symbol-registry-00506}\(\lbrack\!\lbrack A\rbrack\!\rbrack_{\mathcal M}\) & \texttt{\detokenize{\lbrack\!\lbrack\cdot\rbrack\!\rbrack_{\mathcal M}}} & neighborhood truth-set abbreviation & P8.2 & meta & P8.2 & \thepage & Set of worlds satisfying A in the declared neighborhood model; distinct from typed-term and P8.6 semantic brackets. \\
\phantomsection\label{entry:symbol-registry-00507}\(\mathcal A\) & \texttt{\detokenize{\mathcal A}} & admissible proposition family & P8.2 & n/a & P8.2 & \thepage & Boolean-and-modal-operation closed proposition family in a general frame; disjoint from the P8.3 automaton and P8.6 pomonoid. \\
\phantomsection\label{entry:symbol-registry-00508}\(A\alpha,\mathcal F\alpha\) & \texttt{\detokenize{A,}\allowbreak\detokenize{\mathcal F}} & Davies--Humberstone actuality and fixedly operators & P8.2 & 3 & P8.2 & \thepage & Exact source unary operators in the two-index S5 presentation: A resets the lower evaluation index to the upper actual-world index and script F universally varies the upper index. \\
\phantomsection\label{entry:symbol-registry-00509}\(\langle W,V\rangle\models_x^y\alpha\) & \texttt{\detokenize{\models_x^y}} & Davies--Humberstone two-index truth relation & P8.2 & meta & P8.2 & \thepage & Exact source judgment: x is the lower evaluation-world index and y is the upper actual-world perspective. \\
\phantomsection\label{entry:symbol-registry-00510}\(\alpha\) & \texttt{\detokenize{\alpha}} & Davies--Humberstone formula metavariable & P8.2 & n/a & P8.2 & \thepage & Ranges over formulas of the source two-index S5 language; disjoint from the CTL-star path-formula role in P8.3. \\
\phantomsection\label{entry:symbol-registry-00511}\(V(p,x)\) & \texttt{\detokenize{V}} & Davies--Humberstone atomic valuation & P8.2 & n/a & P8.2 & \thepage & Exact source two-place truth-value notation for atom p at world x inside the model pair; locally typed apart from set-valued world valuations. \\
\phantomsection\label{entry:symbol-registry-00512}\(\mathcal M=\langle\mathcal G,\mathcal R,\mathcal D,\mathcal I\rangle\) & \texttt{\detokenize{\mathcal M,}\allowbreak\detokenize{\mathcal G,}\allowbreak\detokenize{\mathcal R,}\allowbreak\detokenize{\mathcal D,}\allowbreak\detokenize{\mathcal I}} & Fitting quantified-modal model datum & P8.2 & n/a & P8.2 & \thepage & Exact source tuple of worlds accessibility world-domain assignment and world-relative interpretation; the local types govern overlapping glyphs. \\
\phantomsection\label{entry:symbol-registry-00513}\(\Gamma,\Delta\) & \texttt{\detokenize{\Gamma,}\allowbreak\detokenize{\Delta}} & Fitting quantified-modal worlds & P8.2 & n/a & P8.2 & \thepage & Source world metavariables in the displayed Fitting model; their local element-of-G type overrides the default context role. \\
\phantomsection\label{entry:symbol-registry-00514}\(\Phi\) & \texttt{\detokenize{\Phi}} & Fitting quantified-modal formula metavariable & P8.2 & n/a & P8.2 & \thepage & Ranges over formulas in the source presentation; locally typed apart from the filtration-set use of Phi. \\
\phantomsection\label{entry:symbol-registry-00515}\(\mathcal D(\Gamma),\mathcal I(P)(\Gamma),\mathcal I(E)(\Gamma)\) & \texttt{\detokenize{\mathcal D,}\allowbreak\detokenize{\mathcal I}} & Fitting local domain and predicate extensions & P8.2 & n/a & P8.2 & \thepage & The source world-domain and world-relative interpretation notation; the special existence predicate satisfies I(E)(Gamma)=D(Gamma). \\
\phantomsection\label{entry:symbol-registry-00516}\((\forall^E x)\Phi,(\exists^E x)\Phi\) & \texttt{\detokenize{\forall^E,}\allowbreak\detokenize{\exists^E}} & Fitting actualist quantifier abbreviations & P8.2 & 3 & P8.2 & \thepage & Exact source relativizations of the primitive possibilist quantifiers by the existence predicate E. \\
\phantomsection\label{entry:symbol-registry-00517}\(E(x)\supset\Phi\) & \texttt{\detokenize{\supset}} & Fitting object-language conditional & P8.2 & 6 & P8.2 & \thepage & Exact source conditional inside the actualist universal-quantifier abbreviation; locally typed apart from the handbook's default implication glyph. \\
\phantomsection\label{entry:symbol-registry-00518}\(\mathcal M,\Gamma\Vdash_v\Phi\) & \texttt{\detokenize{\forces_v}} & Fitting quantified-modal truth judgment & P8.2 & meta & P8.2 & \thepage & Exact source-shaped truth judgment at world Gamma under a world-independent valuation v. \\
\phantomsection\label{entry:symbol-registry-00519}\(\mathsf{LTL},\mathsf{CTL},\mathsf{CTL}^{*},\mathsf{PAL}\) & \texttt{\detokenize{\mathsf{LTL},}\allowbreak\detokenize{\mathsf{CTL},}\allowbreak\detokenize{\mathsf{CTL}^{*},}\allowbreak\detokenize{\mathsf{PAL}}} & temporal and dynamic calculus identifiers & P8.3 & n/a & P8.3 & \thepage & Exact registered LTL CTL full CTL-star and public-announcement presentations. \\
\phantomsection\label{entry:symbol-registry-00520}\(L_{\mu}\) & \texttt{\detokenize{L_{\mu}}} & Kozen modal fixed-point language and calculus identifier & P8.3 & n/a & P8.3 & \thepage & The exact italic identifier used by Kozen for the finitary propositional modal mu-calculus; the same identifier subscripts its equational derivability relation. \\
\phantomsection\label{entry:symbol-registry-00521}\(\mathcal K=(S,R,V),\pi\) & \texttt{\detokenize{\mathcal K,}\allowbreak\detokenize{S,}\allowbreak\detokenize{R,}\allowbreak\detokenize{V,}\allowbreak\detokenize{\pi}} & total transition structure and path & P8.3-P9.5 & n/a & P8.3 & \thepage & K is a total transition structure and pi is an infinite path from the declared state; this path role is disjoint from the P7.8 proof-code metavariable. \\
\phantomsection\label{entry:symbol-registry-00522}\(\mathsf{S5}_{I},\mathsf{S5}_n\) & \texttt{\detokenize{\mathsf{S5}_{I},}\allowbreak\detokenize{\mathsf{S5}_n}} & multi-agent epistemic calculus identifiers & P8.3 & n/a & P8.3 & \thepage & Exact agent-indexed S5 system and its finite-cardinality alias. \\
\phantomsection\label{entry:symbol-registry-00523}\(\mathsf{KD45}_{I},\mathsf{KD45}_n\) & \texttt{\detokenize{\mathsf{KD45}_{I},}\allowbreak\detokenize{\mathsf{KD45}_n}} & multi-agent doxastic calculus identifiers & P8.3 & n/a & P8.3 & \thepage & Exact agent-indexed KD45 system and its finite-cardinality alias. \\
\phantomsection\label{entry:symbol-registry-00524}\(\mathsf{S5}_{I}^{C},\mathsf{S5}_{n}^{C}\) & \texttt{\detokenize{\mathsf{S5}_{I}^{C},}\allowbreak\detokenize{\mathsf{S5}_{n}^{C}}} & common-knowledge calculus identifiers & P8.3 & n/a & P8.3 & \thepage & Exact common-knowledge extension and its finite-cardinality alias. \\
\phantomsection\label{entry:symbol-registry-00525}\(I\) & \texttt{\detokenize{I}} & nonempty agent-index set & P8.3 & n/a & P8.3 & \thepage & Fixed set indexing the epistemic or doxastic modalities and relations; disjoint from the P8.6 idempotent exponential phase set. \\
\phantomsection\label{entry:symbol-registry-00526}\(\mathsf E,\mathsf A\) & \texttt{\detokenize{\mathsf E,}\allowbreak\detokenize{\mathsf A}} & CTL and CTL-star path quantifiers & P8.3-P9.5 & 3 & P8.3 & \thepage & Source-defined existential and universal path quantifiers; scoped away from P8.2 non-normal E. \\
\phantomsection\label{entry:symbol-registry-00527}\(\mathsf{EX},\mathsf{AX},\mathsf{EF},\mathsf{AF}\) & \texttt{\detokenize{\mathsf{EX},}\allowbreak\detokenize{\mathsf{AX},}\allowbreak\detokenize{\mathsf{EF},}\allowbreak\detokenize{\mathsf{AF}}} & CTL next and eventuality operators & P8.3-P9.5 & 3 & P8.3 & \thepage & Source-defined paired existential and universal next and eventuality operators. \\
\phantomsection\label{entry:symbol-registry-00528}\(Z,Y,\mu Z.A,\nu Z.A\) & \texttt{\detokenize{Z,}\allowbreak\detokenize{Y,}\allowbreak\detokenize{\mu,}\allowbreak\detokenize{\nu}} & modal fixed-point variables and binders & P8.3-P9.5 & 3 & P8.3 & \thepage & Z and Y are Kozen's fixed-point variables; mu is the primitive least-fixed-point binder and nu is its capture-avoiding defined dual in the exact L-mu grammar. \\
\phantomsection\label{entry:symbol-registry-00529}\(\zeta,\operatorname{Ext}_{\zeta}(A)\) & \texttt{\detokenize{\zeta,}\allowbreak\detokenize{\operatorname{Ext}_{\zeta}}} & modal fixed-point environment and extension map & P8.3-P9.5 & meta & P8.3 & \thepage & Zeta maps fixed-point variables to state sets and Ext gives the set denoted by A; the P8.3 role is scoped away from the P6.3 forcing-name use. \\
\phantomsection\label{entry:symbol-registry-00530}\(a,b,\langle a\rangle A,[a]A\) & \texttt{\detokenize{a,}\allowbreak\detokenize{b,}\allowbreak\detokenize{\langle a\rangle,}\allowbreak\detokenize{[a]}} & Kozen program constants and indexed modalities & P8.3 & 3 & P8.3 & \thepage & Program constants index Kozen's primitive diamond; the box is its exact defined dual.  Automaton input letters use sigma to keep this role distinct. \\
\phantomsection\label{entry:symbol-registry-00531}\(\mathcal I\) & \texttt{\detokenize{\mathcal I}} & Kozen-model interpretation & P8.3 & n/a & P8.3 & \thepage & Maps propositional constants to state sets and program constants to binary relations in the exact L-mu semantics; disjoint from the P5.4 relative-interpretation datum. \\
\phantomsection\label{entry:symbol-registry-00532}\(\alpha,\beta\) & \texttt{\detokenize{\alpha,}\allowbreak\detokenize{\beta}} & CTL-star path-formula metavariables & P8.3 & n/a & P8.3 & \thepage & Range over Reynolds's path-formula category after the mutually generated CTL-star state/path syntax is declared; disjoint from ordinal and type roles in earlier parts. \\
\phantomsection\label{entry:symbol-registry-00533}\(L,Q;b,b_0\) & \texttt{\detokenize{L,}\allowbreak\detokenize{Q;b,}\allowbreak\detokenize{b_0}} & Reynolds auxiliary atom data & P8.3 & n/a & P8.3 & \thepage & In the exact CTL-star AA side condition \(L\) and \(Q\) are disjoint atom sets and \(b\) and \(b_0\) are auxiliary atoms; \(Q\) is locally typed as an atom set, unlike the Buchi state-set role. \\
\phantomsection\label{entry:symbol-registry-00534}\(B_iA\) & \texttt{\detokenize{B_i}} & agent-indexed belief operator & P8.3 & 3 & P8.3 & \thepage & Doxastic modality for agent i on the registered \(\mathsf{KD45}_{I}\) sheet; distinct from factive \(K_i\) knowledge. \\
\phantomsection\label{entry:symbol-registry-00535}\(D_GA\) & \texttt{\detokenize{D_G}} & distributed-knowledge operator & P8.3 & 3 & P8.3 & \thepage & Truth at every successor of the intersection of the group relations; distinct from everyone-knows and common knowledge. \\
\phantomsection\label{entry:symbol-registry-00536}\(R_G^{E},R_G^{C},R_G^{D}\) & \texttt{\detokenize{R_G^{E},}\allowbreak\detokenize{R_G^{C},}\allowbreak\detokenize{R_G^{D}}} & group epistemic relations & P8.3 & n/a & P8.3 & \thepage & Union relation its reflexive-transitive closure and relation intersection for everyone common and distributed knowledge. \\
\phantomsection\label{entry:symbol-registry-00537}\(U\,A\,B,S\,A\,B\) & \texttt{\detokenize{U,}\allowbreak\detokenize{S}} & Kamp strict tense connectives & P8.3 & 3 & P8.3 & \thepage & Kamp's prefix source syntax: the first operand holds at the future or past witness and the second at every strictly intervening point. \\
\phantomsection\label{entry:symbol-registry-00538}\(\mathsf P,\mathsf H\) & \texttt{\detokenize{\mathsf P,}\allowbreak\detokenize{\mathsf H}} & unary past tense operators & P8.3 & 3 & P8.3 & \thepage & Prior-style sometime-past and always-past operators; their temporal standing is fixed locally. \\
\phantomsection\label{entry:symbol-registry-00539}\([!\chi]A,\langle!\chi\rangle A\) & \texttt{\detokenize{[!\chi],}\allowbreak\detokenize{\langle!\chi\rangle}} & public-announcement modalities & P8.3 & 3 & P8.3 & \thepage & Truthful-update box and its executable dual; a false box precondition is vacuously true. \\
\phantomsection\label{entry:symbol-registry-00540}\(\mathcal A=(Q,\Sigma,\delta,q_0,F)\) & \texttt{\detokenize{\mathcal A,}\allowbreak\detokenize{Q,}\allowbreak\detokenize{\Sigma,}\allowbreak\detokenize{\delta,}\allowbreak\detokenize{q_0,}\allowbreak\detokenize{F}} & Buchi automaton datum & P8.3-P9.4 & n/a & P8.3 & \thepage & Finite state set alphabet set-valued transition function initial state and accepting set; a deterministic specialization has singleton transition sets. \\
\phantomsection\label{entry:symbol-registry-00541}\(\operatorname{Inf}(r)\) & \texttt{\detokenize{\operatorname{Inf}}} & infinitely recurring run states & P8.3-P9.4 & meta & P8.3 & \thepage & States occurring infinitely often in run r; Buchi acceptance intersects F while a Rabin pair separately tests its E and F components. \\
\phantomsection\label{entry:symbol-registry-00542}\(\mathbf K_3,\mathsf{LP},\mathbf{FDE}\) & \texttt{\detokenize{\mathbf K_3,}\allowbreak\detokenize{\mathsf{LP},}\allowbreak\detokenize{\mathbf{FDE}}} & three-valued and four-valued logic identifiers & P8.5 & n/a & P8.5 & \thepage & The inherited presentation controls bold K3 and sans-serif LP; Omori--Wansing controls bold upright FDE.  The weak-Kleene matrix remains unnamed because Kleene's presentation supplies no symbolic system identifier. \\
\phantomsection\label{entry:symbol-registry-00543}\(\text{\L}_n\) & \texttt{\detokenize{\text{\L}_n}} & finite Lukasiewicz chain & P8.5 & n/a & P8.5 & \thepage & Cignoli--D'Ottaviano--Mundici's stroked-L name for the n-element chain itself; the corresponding n-valued system is described in source-bound prose. \\
\phantomsection\label{entry:symbol-registry-00544}\(\text{\L}\) & \texttt{\detokenize{\text{\L}}} & standard infinite-valued Lukasiewicz logic identifier & P8.5 & n/a & P8.5 & \thepage & Hajek's exact stroked-L identifier for the infinite-valued logic; its standard value algebra is the unit interval. \\
\phantomsection\label{entry:symbol-registry-00545}\(\mathrm{MTL}\) & \texttt{\detokenize{\mathrm{MTL}}} & monoidal t-norm logic identifier & P8.5 & n/a & P8.5 & \thepage & Esteva--Godo's upright Roman system identifier; the source controls both its letters and typography. \\
\phantomsection\label{entry:symbol-registry-00546}\(\mathrm{BL},G,\Pi\) & \texttt{\detokenize{\mathrm{BL},}\allowbreak\detokenize{G,}\allowbreak\detokenize{\Pi}} & basic Goedel and Product fuzzy-logic identifiers & P8.5 & n/a & P8.5 & \thepage & Hajek's source identifiers: upright Roman BL, plain italic G, and plain Pi; local types distinguish G and Pi from metavariables. \\
\phantomsection\label{entry:symbol-registry-00547}\(\mathrm{BD}\) & \texttt{\detokenize{\mathrm{BD}}} & Belnap--Dunn four-valued algebra & P8.5 & n/a & P8.5 & \thepage & Omori--Wansing's upright source identifier for the four-valued algebra, represented here by the exact positive/negative support pairs. \\
\phantomsection\label{entry:symbol-registry-00548}\(x^+,x^-\) & \texttt{\detokenize{{}^+,}\allowbreak\detokenize{{}^-}} & positive and negative support coordinates & P8.5 & n/a & P8.5 & \thepage & Boolean support indicators of an FDE value; the two coordinates are required to distinguish both from neither. \\
\phantomsection\label{entry:symbol-registry-00549}\(A{:}a\) & \texttt{\detokenize{A{:}a}} & matrix value-labelled formula & P8.5 & meta & P8.5 & \thepage & Tableau occurrence asserting that formula A has exactly matrix value a; distinct from term typing and world labels. \\
\phantomsection\label{entry:symbol-registry-00550}\(f,n,b,t\) & \texttt{\detokenize{f,}\allowbreak\detokenize{n,}\allowbreak\detokenize{b,}\allowbreak\detokenize{t}} & Belnap--Dunn value names & P8.5 & n/a & P8.5 & \thepage & Omori--Wansing's source letters for false-only neither both and true-only values in the BD algebra. \\
\phantomsection\label{entry:symbol-registry-00551}\(\leq_t,\leq_k\) & \texttt{\detokenize{\leq_t,}\allowbreak\detokenize{\leq_k}} & Belnap--Dunn truth and knowledge orders & P8.5 & meta & P8.5 & \thepage & Omori--Wansing's source subscripts for the distinct truth and knowledge orders on BD; neither is a consequence relation. \\
\phantomsection\label{entry:symbol-registry-00552}\(x\Rightarrow y\) & \texttt{\detokenize{\Rightarrow}} & t-norm residuum and Lukasiewicz implication & P8.5 & 6 & P8.5 & \thepage & Hajek's plain source glyph.  For the Lukasiewicz t-norm it is min of one and one-minus-x-plus-y; numeric operands type it as a truth operation rather than a sequent arrow. \\
\phantomsection\label{entry:symbol-registry-00553}\(x\mathbin{\ast_{\text{\L}}}y,x\mathbin{\ast_{\Pi}}y\) & \texttt{\detokenize{\mathbin{\ast_{\text{\L}}},}\allowbreak\detokenize{\mathbin{\ast_{\Pi}}}} & Lukasiewicz and Product t-norms & P8.5 & 4 & P8.5 & \thepage & Hajek's source subscripts use stroked L for the Lukasiewicz max-zero-sum-minus-one t-norm and Pi for the Product multiplication t-norm. \\
\phantomsection\label{entry:symbol-registry-00554}\(x\ast y\) & \texttt{\detokenize{\ast}} & t-norm operation & P8.5 & 4 & P8.5 & \thepage & Hajek's generic associative commutative monotone binary operation on the unit interval with unit one. \\
\phantomsection\label{entry:symbol-registry-00555}\(\mathbf E,\mathbf R,\mathbf{RM}\) & \texttt{\detokenize{\mathbf E,}\allowbreak\detokenize{\mathbf R,}\allowbreak\detokenize{\mathbf{RM}}} & relevance-system identifiers & P8.6 & n/a & P8.6 & \thepage & Bold Roman system typography reproduced from Dunn--Restall's Anderson--Belnap convention; locally typed as system identifiers. \\
\phantomsection\label{entry:symbol-registry-00556}\(\mathbf E_{\to},\mathbf R_{\to}\) & \texttt{\detokenize{\mathbf E_{\to},}\allowbreak\detokenize{\mathbf R_{\to}}} & pure implicational relevance-fragment identifiers & P8.6 & n/a & P8.6 & \thepage & Dunn--Restall's source subscripts for the pure implicational fragments of bold Roman E and R. \\
\phantomsection\label{entry:symbol-registry-00557}\(\gamma\) & \texttt{\detokenize{\gamma}} & Ackermann theoremhood metarule & P8.6 & meta & P8.6 & \thepage & From theoremhood of A and of not-A-or-B infer theoremhood of B; admissible for E and R but not a primitive or local-consequence rule. \\
\phantomsection\label{entry:symbol-registry-00558}\(\mathfrak F\) & \texttt{\detokenize{\mathfrak F}} & declared frame datum & P8.2-P8.6 & n/a & P8.2 & \thepage & The locally declared tuple fixes whether this is a modal Kripke frame a broad Routley--Meyer frame or the pointed Routley--Meyer frame used by the cited presentation; disjoint from the P6.3 normal-subgroup-filter use. \\
\phantomsection\label{entry:symbol-registry-00559}\(a\leq b\) & \texttt{\detokenize{\leq}} & derived relevance-frame preorder & P8.6 & meta & P8.6 & \thepage & Metalinguistic abbreviation for R(0,a,b) in the displayed pointed frame; distinct from the algebraic orders later in the chapter. \\
\phantomsection\label{entry:symbol-registry-00560}\(R^2(ab)cd,R^2a(bc)d\) & \texttt{\detokenize{R^2}} & composite ternary-frame relations & P8.6 & meta & P8.6 & \thepage & The first two exact existential composites in the cited E R and RM frame presentation; parenthesization fixes the composition. \\
\phantomsection\label{entry:symbol-registry-00561}\(R^2b(ac)d,R^2abbc\) & \texttt{\detokenize{R^2}} & composite ternary-frame relations & P8.6 & meta & P8.6 & \thepage & The remaining exact existential composites in the cited E R and RM frame presentation; parenthesization fixes the composition. \\
\phantomsection\label{entry:symbol-registry-00562}\(\mathsf{IMLL},\mathsf{CLL}\) & \texttt{\detokenize{\mathsf{IMLL},}\allowbreak\detokenize{\mathsf{CLL}}} & linear-calculus identifiers & P8.6 & n/a & P8.6 & \thepage & Named multiplicative intuitionistic core and distinct full one-sided classical calculus.  The inherited four-rule bang extension is unnamed and therefore has no fabricated identifier. \\
\phantomsection\label{entry:symbol-registry-00563}\(A\bindnasrepma B\) & \texttt{\detokenize{\bindnasrepma}} & multiplicative disjunction (par) & P8.6 & 4 & P8.6 & \thepage & CLL connective dual to tensor; shares the disjunction rank only in this exact scope. \\
\phantomsection\label{entry:symbol-registry-00564}\(A\mathbin{\&}B\) & \texttt{\detokenize{\mathbin{\&}}} & additive conjunction (with) & P8.6 & 4 & P8.6 & \thepage & CLL additive conjunction; distinct from metalinguistic ampersands. \\
\phantomsection\label{entry:symbol-registry-00565}\(A\oplus B\) & \texttt{\detokenize{\oplus}} & additive disjunction (plus) & P8.6 & 5 & P8.6 & \thepage & CLL additive disjunction; exact P8.6 scope is disjoint from the P4.4 computable join role. \\
\phantomsection\label{entry:symbol-registry-00566}\(0\) & \texttt{\detokenize{0}} & CLL additive false unit & P8.6 & n/a & P8.6 & \thepage & Additive unit dual to top in the displayed CLL grammar; not the relevance base point or an ordinary numeral. \\
\phantomsection\label{entry:symbol-registry-00567}\(?A\) & \texttt{\detokenize{?}} & why-not exponential & P8.6 & 2.5 & P8.6 & \thepage & CLL exponential dual to bang; prefix unary and the only side supporting the displayed structural rules. \\
\phantomsection\label{entry:symbol-registry-00568}\(A^\perp\) & \texttt{\detokenize{{}^\perp}} & formula-level linear negation & P8.6 & 2.5 & P8.6 & \thepage & Involutive De Morgan operation on CLL formulas; distinct from phase orthogonality on subsets. \\
\phantomsection\label{entry:symbol-registry-00569}\(X^\perp\) & \texttt{\detokenize{{}^\perp}} & phase orthogonal of a subset & P8.6 & meta & P8.6 & \thepage & Girard's source notation for subset orthogonality relative to the phase-space pole.  Subset operands distinguish it from formula linear negation without altering the glyph. \\
\phantomsection\label{entry:symbol-registry-00570}\(\perp\subseteq M\) & \texttt{\detokenize{\perp}} & Girard phase-space pole & P8.6 & n/a & P8.6 & \thepage & Girard's exact perpendicular glyph, locally typed as the distinguished subset of the phase monoid rather than formula linear negation. \\
\phantomsection\label{entry:symbol-registry-00571}\(\sem A\) & \texttt{\detokenize{\sem}} & IMLL algebraic interpretation & P8.6 & meta & P8.6 & \thepage & The inherited source's semantic value of formula A in the declared commutative residuated partially ordered monoid. \\
\phantomsection\label{entry:symbol-registry-00572}\(A^*\) & \texttt{\detokenize{{}^*}} & Girard phase interpretation of a CLL formula & P8.6 & meta & P8.6 & \thepage & Fact assigned to formula A in Girard's source notation; phase validity is membership of the monoid unit 1 in A-star. \\
\phantomsection\label{entry:symbol-registry-00573}\(I\) & \texttt{\detokenize{I}} & idempotent exponential phase set & P8.6 & n/a & P8.6 & \thepage & Idempotent members of the tensor-unit fact used in the displayed full-CLL exponential clauses. \\
\phantomsection\label{entry:symbol-registry-00574}\(A\mathbin{\cdot}B\) & \texttt{\detokenize{\mathbin{\cdot}}} & Lambek product & P8.6 & 4 & P8.6 & \thepage & Associative unit-free product in Lambek L; operand type distinguishes it from phase and arithmetic multiplication. \\
\phantomsection\label{entry:symbol-registry-00575}\(A\backslash B,B/A\) & \texttt{\detokenize{\backslash,}\allowbreak\detokenize{/}} & Lambek residuals & P8.6 & 6 & P8.6 & \thepage & Left and right divisions of the ordered Lambek product; neither admits exchange weakening or contraction. \\
\phantomsection\label{entry:symbol-registry-00576}\((M,\perp);mn;1\) & \texttt{\detokenize{M,}\allowbreak\detokenize{\perp,}\allowbreak\detokenize{mn,}\allowbreak\detokenize{1}} & Girard phase-space datum and monoid notation & P8.6 & n/a & P8.6 & \thepage & Girard's phase space is the pair of a commutative monoid M and perpendicular pole; multiplication is juxtaposition and the monoid unit is 1. \\
\phantomsection\label{entry:symbol-registry-00577}\(X,Y\) & \texttt{\detokenize{X,}\allowbreak\detokenize{Y}} & phase subsets and facts & P8.6 & n/a & P8.6 & \thepage & Girard uses the same source metavariables for arbitrary phase subsets and for biorthogonally closed facts; the locally stated closure property supplies the fact role. \\
\phantomsection\label{entry:symbol-registry-00578}\(\mathsf L^*,\mathsf L_1\) & \texttt{\detokenize{\mathsf L^*,}\allowbreak\detokenize{\mathsf L_1}} & Lambek extension identifiers & P8.6 & n/a & P8.6 & \thepage & Contrast names for empty-antecedent and unit extensions; neither is the associative unit-free Lambek-L calculus displayed on its card. \\
\phantomsection\label{entry:symbol-registry-00579}\(\lbrack A\rbrack\) & \texttt{\detokenize{\lbrack\cdot\rbrack}} & IMLL Lindenbaum quotient class & P8.6 & meta & P8.6 & \thepage & Equivalence class of formula A in the annex sketch's syntactic algebra; distinct from semantic double brackets. \\
\phantomsection\label{entry:symbol-registry-00580}\(\Phi\) & \texttt{\detokenize{\Phi}} & ordered Lambek context metavariable & P8.6 & n/a & P8.6 & \thepage & Possibly empty context segment in rule notation subject to every displayed antecedent remaining nonempty. \\
\phantomsection\label{entry:symbol-registry-00581}\(i,j\) & \texttt{\detokenize{i,}\allowbreak\detokenize{j}} & hybrid-logic nominals & P8.3 & n/a & P8.3 & \thepage & Object-language atoms constrained to name exactly one world each; disjoint from metalinguistic world labels. \\
\phantomsection\label{entry:symbol-registry-00582}\(@_iA\) & \texttt{\detokenize{@_i}} & hybrid satisfaction operator & P8.3 & 3 & P8.3 & \thepage & Evaluates A at the unique world named by nominal i; prefix unary and internal to the object language. \\
\phantomsection\label{entry:symbol-registry-00583}\(@_i\Diamond j\) & \texttt{\detokenize{@_i\Diamond}} & hybrid tableau accessibility atom & P8.3 & meta & P8.3 & \thepage & Branch notation asserting that the world named by j is accessible from the world named by i. \\
\phantomsection\label{entry:symbol-registry-00584}\(\mathcal I(c)(\Gamma)\) & \texttt{\detokenize{\mathcal I(c)}} & world-relative interpretation of a quantified-modal constant & P8.2 & n/a & P8.2 & \thepage & The value of constant c at Fitting world Gamma; rigidity means equality of these values at every pair of worlds and is independent of inner-domain existence. \\
\phantomsection\label{entry:symbol-registry-00585}\(\mathsf f,\mathsf u,\mathsf t\) & \texttt{\detokenize{\mathsf f,}\allowbreak\detokenize{\mathsf u,}\allowbreak\detokenize{\mathsf t}} & neutral free-logic truth values & P8.4 & n/a & P8.4 & \thepage & Strong-Kleene false undefined and true values for the partial-term variant; only true is designated. \\
\phantomsection\label{entry:symbol-registry-00586}\(\models_{\mathrm{neg}},\models_{\mathrm{pos}},\models_{\mathrm{neu}}\) & \texttt{\detokenize{\models_{\mathrm{neg}},}\allowbreak\detokenize{\models_{\mathrm{pos}},}\allowbreak\detokenize{\models_{\mathrm{neu}}}} & free-logic variant consequence relations & P8.4 & meta & P8.4 & \thepage & Truth preservation in the negative and positive two-valued rows and designated-t preservation in the neutral strong-Kleene row; the subscript fixes the atomic and partial-term policy. \\
\phantomsection\label{entry:symbol-registry-00587}\(\mathsf{SDL},O A,\mathsf P_OA\) & \texttt{\detokenize{\mathsf{SDL},}\allowbreak\detokenize{O,}\allowbreak\detokenize{\mathsf P_O}} & standard deontic identifier obligation and permission & P8.4 & 3 & P8.4 & \thepage & SDL is KD under the deontic reading; permission is the eliminable abbreviation \(\mathsf P_OA\defeq\neg O\neg A\). \\
\phantomsection\label{entry:symbol-registry-00588}\(O(A\mid B)\) & \texttt{\detokenize{O(\cdot\mid\cdot)}} & dyadic obligation & P8.4 & 6 & P8.4 & \thepage & Binary modality evaluated at the best accessible B-worlds; it is not the monadic formula \(O(B\to A)\). \\
\phantomsection\label{entry:symbol-registry-00589}\(\preceq_w,\operatorname{Best}_w(B)\) & \texttt{\detokenize{\preceq_w,}\allowbreak\detokenize{\operatorname{Best}_w}} & world-relative ideality preorder and best-world selector & P8.4 & n/a & P8.4 & \thepage & A smooth total preorder on accessible worlds and the explicitly defined set of its minimal accessible B-worlds; the set is empty when there are no candidates. \\
\phantomsection\label{entry:symbol-registry-00590}\(A>B\) & \texttt{\detokenize{>}} & conditional connective & P8.4 & 6 & P8.4 & \thepage & Formula operands type this glyph as the primitive CK or Stalnaker conditional; numerical operands retain the strict-order role. \\
\phantomsection\label{entry:symbol-registry-00591}\(A\Box\!\!\to B\) & \texttt{\detokenize{\Box\!\!\to}} & Lewis variably strict conditional & P8.4 & 6 & P8.4 & \thepage & Evaluated by nested similarity spheres; distinct from material implication and the unique-selection conditional. \\
\phantomsection\label{entry:symbol-registry-00592}\(f(w,X),F_w(X)\) & \texttt{\detokenize{f,}\allowbreak\detokenize{F_w}} & singular and set-valued conditional selection functions & P8.4 & n/a & P8.4 & \thepage & Stalnaker f selects one member of a nonempty proposition; CK \(F_w\) returns an arbitrary set and need not be contained in X. \\
\phantomsection\label{entry:symbol-registry-00593}\(\mathcal S_w,U_w\) & \texttt{\detokenize{\mathcal S_w,}\allowbreak\detokenize{U_w}} & Lewis sphere system and accessible union & P8.4 & n/a & P8.4 & \thepage & The inclusion-nested union/intersection-closed sphere family at w and its union; no least antecedent-permitting sphere is assumed. \\
\phantomsection\label{entry:symbol-registry-00594}\(\mathsf{CK},\mathsf V,\mathsf{VW},\mathsf{VC},\mathsf{C1},\mathsf{C2}\) & \texttt{\detokenize{\mathsf{CK},}\allowbreak\detokenize{\mathsf V,}\allowbreak\detokenize{\mathsf{VW},}\allowbreak\detokenize{\mathsf{VC},}\allowbreak\detokenize{\mathsf{C1},}\allowbreak\detokenize{\mathsf{C2}}} & conditional-logic identifiers & P8.4 & n/a & P8.4 & \thepage & Basic conditional logic and the separately stated Lewis and Stalnaker extensions; VC is also called C1 in the cited nomenclature. \\
\phantomsection\label{entry:symbol-registry-00595}\(P(B\mid A)\) & \texttt{\detokenize{P(\cdot\mid\cdot)}} & conditional probability & P8.4 & meta & P8.4 & \thepage & Defined when \(P(A)>0\); probability P is distinct from the deontic permission abbreviation \(\mathsf P_O\). \\
\phantomsection\label{entry:symbol-registry-00596}\(A\mathbin{\&}B\) & \texttt{\detokenize{\mathbin{\&}}} & strong t-norm conjunction & P8.5 & 4 & P8.5 & \thepage & Primitive multiplicative conjunction of the MTL and BL-family sheets; distinct from weak lattice conjunction and CLL additive with. \\
\phantomsection\label{entry:symbol-registry-00597}\(\mathsf L_n,\mathsf L_{\infty}\) & \texttt{\detokenize{\mathsf L_n,}\allowbreak\detokenize{\mathsf L_{\infty}}} & finite and infinite Lukasiewicz system identifiers & P8.5 & n/a & P8.5 & \thepage & Editorial system-card identifiers; the source symbols \(\text{\L}_n\) and \(\text{\L}\) name the finite chain and standard infinite logic. \\
\phantomsection\label{entry:symbol-registry-00598}\(\derives_{\mathrm{MTL}},\derives_{\mathrm{BL}},\derives_G,\derives_{\Pi}\) & \texttt{\detokenize{\derives_{\mathrm{MTL}},}\allowbreak\detokenize{\derives_{\mathrm{BL}},}\allowbreak\detokenize{\derives_G,}\allowbreak\detokenize{\derives_{\Pi}}} & t-norm Hilbert derivability relations & P8.5 & meta & P8.5 & \thepage & The subscript fixes one exact sheet; their standard completeness strengths are stated separately. \\
\phantomsection\label{entry:symbol-registry-00599}\(\mathsf C_1,\derives_{\mathsf C_1}\) & \texttt{\detokenize{\mathsf C_1,}\allowbreak\detokenize{\derives_{\mathsf C_1}}} & da Costa C1 calculus and derivability & P8.7 & meta & P8.7 & \thepage & The source-fixed propositional calculus on sheet:hilbert-da-costa-c1; no claim is transferred to every member of the da Costa hierarchy. \\
\phantomsection\label{entry:symbol-registry-00600}\(A^\circ\) & \texttt{\detokenize{{}^\circ}} & da Costa C1 consistency formula & P8.7 & 2.5 & P8.7 & \thepage & Eliminable abbreviation \(\neg(A\land\neg A)\) expressing formula-level well-behavedness in C1; not a matrix truth value or a metalinguistic theory-consistency claim. \\
\phantomsection\label{entry:symbol-registry-00601}\(\mathrel{|\!\sim}\) & \texttt{\detokenize{|\!\sim}} & defeasible consequence & P8.7 & meta & P8.7 & \thepage & A relation between information states or formulas whose conclusions may be defeated by additional premises. \\
\phantomsection\label{entry:symbol-registry-00602}\((S,\ell,\prec),\mu(X)\) & \texttt{\detokenize{S,}\allowbreak\detokenize{\ell,}\allowbreak\detokenize{\prec,}\allowbreak\detokenize{\mu}} & preferential-model datum and minimal-state operator & P8.7 & meta & P8.7 & \thepage & The labeling map gives classical valuations; \(\mu(X)\) selects the \(\prec\)-minimal states of X. \\
\phantomsection\label{entry:symbol-registry-00603}\(K*A,K+A,K-A\) & \texttt{\detokenize{*,}\allowbreak\detokenize{+,}\allowbreak\detokenize{-}} & AGM revision expansion and contraction & P8.7 & meta & P8.7 & \thepage & Operations on the deductively closed belief set K; their standing is fixed by the AGM postulates and Levi and Harper identities. \\
\phantomsection\label{entry:symbol-registry-00604}\(\preceq_K,\min([A],\preceq_K)\) & \texttt{\detokenize{\preceq_K,}\allowbreak\detokenize{\min}} & AGM faithful preorder and minimal input valuations & P8.7 & meta & P8.7 & \thepage & A total preorder assigned to prior belief set K; the minimum selects its most preferred A-valuations. \\
\phantomsection\label{entry:symbol-registry-00605}\(\frac{\alpha:\beta_1,\ldots,\beta_n}{\gamma}\) & \texttt{\detokenize{:,}\allowbreak\detokenize{/}} & Reiter default & P8.7 & meta & P8.7 & \thepage & Prerequisite alpha justifications beta-i and conclusion gamma; this fraction is not an ordinary inference rule. \\
\phantomsection\label{entry:symbol-registry-00606}\((W,D),\Gamma_E,U_k\) & \texttt{\detokenize{W,}\allowbreak\detokenize{D,}\allowbreak\detokenize{\Gamma_E,}\allowbreak\detokenize{U_k}} & default theory grounded closure and stages & P8.7 & meta & P8.7 & \thepage & Gamma-E is the least closure grounded through the stages U-k while justification consistency is tested against candidate E. \\
\phantomsection\label{entry:symbol-registry-00607}\(T(\bar P,\bar Q,\bar Z),\bar P'\leq\bar P\) & \texttt{\detokenize{T,}\allowbreak\detokenize{\bar P,}\allowbreak\detokenize{\bar Q,}\allowbreak\detokenize{\bar Z,}\allowbreak\detokenize{\leq}} & circumscription theory and predicate-tuple inclusion & P8.7 & meta & P8.7 & \thepage & P is minimized Q is fixed and Z may vary; tuple inclusion abbreviates componentwise universal closure. \\
\phantomsection\label{entry:symbol-registry-00608}\(LA,\operatorname{Th}\) & \texttt{\detokenize{L,}\allowbreak\detokenize{\operatorname{Th}}} & autoepistemic belief and classical closure & P8.7 & 3 & P8.7 & \thepage & L is the belief operator; Th is ordinary classical propositional closure over the autoepistemic language and adds no modal axiom. \\
\end{longtable}
\endgroup
\Needspace{.4\textheight}
\subsection*{Part 9 scopes}
\begingroup\scriptsize\setlength{\tabcolsep}{3pt}
\begin{longtable}{@{}p{0.9in}p{1.3in}p{1.0in}p{0.5in}p{0.3in}p{0.4in}p{0.35in}p{1.65in}@{}}
\toprule
\textbf{Glyph} & \textbf{Macro} & \textbf{Role} & \textbf{Scope} & \textbf{Prec} & \textbf{First} & \textbf{Reg. page} & \textbf{Principal uses / notes} \\
\midrule
\endfirsthead
\toprule
\textbf{Glyph} & \textbf{Macro} & \textbf{Role} & \textbf{Scope} & \textbf{Prec} & \textbf{First} & \textbf{Reg. page} & \textbf{Principal uses / notes} \\
\midrule
\endhead
\midrule\multicolumn{8}{r}{\emph{continued}}\\\endfoot
\bottomrule\endlastfoot
\phantomsection\label{entry:symbol-registry-00635}\(\tau_{\mathrm{fin}}\) & \texttt{\detokenize{\tau_{\mathrm{fin}}}} & fixed finite relational signature & P9.1 & n/a & P9.1 & \thepage & The signature relative to which finite satisfiability and finite validity are defined; it is not the varying database signature of a combined-complexity input. \\
\phantomsection\label{entry:symbol-registry-00636}\(\mathcal{ALC}\) & \texttt{\detokenize{\mathcal{ALC}}} & description-logic language and system identifier & P9.3 & n/a & P9.3 & \thepage & The exact Boolean concept language with qualified role restrictions from Definition def:p09-alc; extensions in the description-logic table have different identifiers. \\
\phantomsection\label{entry:symbol-registry-00637}\(I=(\Delta^I,\mathord{\cdot}^I)\) & \texttt{\detokenize{I,}\allowbreak\detokenize{\Delta^I,}\allowbreak\detokenize{\mathord{\cdot}^I}} & description-logic interpretation datum & P9.3 & n/a & P9.3 & \thepage & Delta-I is the nonempty individual domain and dot-I interprets concept role and individual names with the types fixed in Definition def:p09-alc. \\
\phantomsection\label{entry:symbol-registry-00638}\(\mathcal T,\mathcal A\) & \texttt{\detokenize{\mathcal T,}\allowbreak\detokenize{\mathcal A}} & ALC TBox and ABox & P9.3 & n/a & P9.3 & \thepage & Finite sets respectively of concept inclusions and concept or role assertions; script A is locally typed apart from automata. \\
\phantomsection\label{entry:symbol-registry-00639}\(C\sqsubseteq D\) & \texttt{\detokenize{\sqsubseteq}} & description-logic concept inclusion & P9.3 & 6 & P9.3 & \thepage & An object-language TBox axiom interpreted by subset inclusion of concept extensions; not metalinguistic formula consequence. \\
\phantomsection\label{entry:symbol-registry-00640}\(\exists r.C,\forall r.C\) & \texttt{\detokenize{\exists r.,}\allowbreak\detokenize{\forall r.}} & ALC role restrictions & P9.3 & 3 & P9.3 & \thepage & Qualified existential and universal concept constructors interpreted over the declared role relation. \\
\phantomsection\label{entry:symbol-registry-00641}\(\operatorname{nnf},\operatorname{cl}\) & \texttt{\detokenize{\operatorname{nnf},}\allowbreak\detokenize{\operatorname{cl}}} & ALC normalization and finite closure operations & P9.3 & meta & P9.3 & \thepage & Metalinguistic operations used by blocking completion and type elimination; they add no concept constructors. \\
\phantomsection\label{entry:symbol-registry-00642}\(\mathsf{Glob}(\mathcal T)\) & \texttt{\detokenize{\mathsf{Glob}}} & ALC tableau global-GCI abbreviation & P9.3 & meta & P9.3 & \thepage & Metalinguistic abbreviation for the NNF internalizations of every GCI in the active TBox; used in body initialization and on sheet:tableaux-alc. \\
\phantomsection\label{entry:symbol-registry-00643}\(\mathcal L(x),\mathcal L(y)\subseteq\mathcal L(x)\) & \texttt{\detokenize{\mathcal L,}\allowbreak\detokenize{\subseteq}} & ALC completion label and subset-blocking comparison & P9.3 & meta & P9.3 & \thepage & Labels are subsets of the finite closure; a nonroot y is directly blocked by ancestor x exactly under the displayed subset comparison, re-evaluated after label growth. \\
\phantomsection\label{entry:symbol-registry-00644}\(\widehat y,b(y)\) & \texttt{\detokenize{\widehat y,}\allowbreak\detokenize{b}} & ALC retained-edge target and direct blocker & P9.3 & meta & P9.3 & \thepage & Defined only for a target y of an edge retained from an unblocked source; widehat-y is y when unblocked and otherwise a directly blocking unblocked ancestor b(y). \\
\phantomsection\label{entry:symbol-registry-00645}\(\mathsf Y A\) & \texttt{\detokenize{\mathsf Y}} & temporal previous operator & P9.4 & 3 & P9.4 & \thepage & Prefix unary past operator on pointed words; false at position zero. \\
\phantomsection\label{entry:symbol-registry-00646}\(A\mathbin{\mathsf S}B\) & \texttt{\detokenize{\mathbin{\mathsf S}}} & temporal since operator & P9.4 & 3.5 & P9.4 & \thepage & Past-time binary operator with a witness at or before the current position and A at every strictly intervening position. \\
\phantomsection\label{entry:symbol-registry-00647}\((Q,\Sigma,\delta,q_0,\Omega)\) & \texttt{\detokenize{Q,}\allowbreak\detokenize{\Sigma,}\allowbreak\detokenize{\delta,}\allowbreak\detokenize{q_0,}\allowbreak\detokenize{\Omega}} & deterministic parity word-automaton datum & P9.4 & n/a & P9.4 & \thepage & Omega assigns a finite priority to each state; the least priority occurring infinitely often determines acceptance. \\
\phantomsection\label{entry:symbol-registry-00648}\((E_i,F_i)_{i<k}\) & \texttt{\detokenize{E_i,}\allowbreak\detokenize{F_i}} & Rabin acceptance pairs & P9.4 & n/a & P9.4 & \thepage & A run accepts when one pair has no infinitely recurring E-state and at least one infinitely recurring F-state. \\
\phantomsection\label{entry:symbol-registry-00649}\(\mathsf{AP}(A),\ell_A(s)\) & \texttt{\detokenize{\mathsf{AP},}\allowbreak\detokenize{\ell_A}} & formula alphabet and induced state-label map & P9.4 & meta & P9.4 & \thepage & AP(A) is the finite set of proposition letters occurring in A; ell-A maps a state to the subset true there under the predicate-indexed valuation V. \\
\phantomsection\label{entry:symbol-registry-00650}\(\Sigma=\bigsqcup_{n\geq0}\Sigma_n,T_\Sigma\) & \texttt{\detokenize{\Sigma,}\allowbreak\detokenize{\Sigma_n,}\allowbreak\detokenize{T_\Sigma}} & finite ranked alphabet and finite tree set & P9.4 & meta & P9.4 & \thepage & Every symbol in Sigma-n has arity n and T-Sigma is the least finite-term set generated according to those ranks. \\
\phantomsection\label{entry:symbol-registry-00651}\(\mathcal A=(Q,\Sigma,\Delta,F),L(\mathcal A)\) & \texttt{\detokenize{\mathcal A,}\allowbreak\detokenize{Q,}\allowbreak\detokenize{\Sigma,}\allowbreak\detokenize{\Delta,}\allowbreak\detokenize{F,}\allowbreak\detokenize{L}} & bottom-up finite-tree automaton and accepted language & P9.4 & meta & P9.4 & \thepage & Delta contains ranked bottom-up transitions and acceptance requires a run whose root state lies in F; distinct from infinite-word and Rabin infinite-tree tuples. \\
\phantomsection\label{entry:symbol-registry-00652}\(\operatorname{Pre}_{\exists},\operatorname{Pre}_{\forall}\) & \texttt{\detokenize{\operatorname{Pre}_{\exists},}\allowbreak\detokenize{\operatorname{Pre}_{\forall}}} & existential and universal predecessor maps & P9.5 & meta & P9.5 & \thepage & Algorithmic maps from state sets to state sets in a finite transition structure; not CTL path quantifiers. \\
\phantomsection\label{entry:symbol-registry-00653}\(L(A)\) & \texttt{\detokenize{L}} & CTL labelling truth set & P9.5 & meta & P9.5 & \thepage & The set of states assigned to subformula A by bottom-up CTL labelling; locally typed apart from languages and well-founded bounds. \\
\phantomsection\label{entry:symbol-registry-00654}\(\operatorname{gfp}(F)\) & \texttt{\detokenize{\operatorname{gfp}}} & greatest fixed point of a monotone map & P9.5-P9.7 & meta & P9.5 & \thepage & The greatest set fixed by the displayed monotone operator; CTL formulas continue to use the object-language nu binder. \\
\phantomsection\label{entry:symbol-registry-00655}\(\Sigma_n^{N\mu},\Pi_n^{N\mu};\Sigma_n^{S\mu},\Pi_n^{S\mu};\Sigma_n^{EL\mu},\Pi_n^{EL\mu}\) & \texttt{\detokenize{\Sigma_n^{N\mu},}\allowbreak\detokenize{\Pi_n^{N\mu},}\allowbreak\detokenize{\Sigma_n^{S\mu},}\allowbreak\detokenize{\Pi_n^{S\mu},}\allowbreak\detokenize{\Sigma_n^{EL\mu},}\allowbreak\detokenize{\Pi_n^{EL\mu}}} & Niwiński simple and Emerson--Lei modal-mu alternation classes & P9.5 & meta & P9.5 & \thepage & Source-specific positive-form syntactic classes; the variants differ by capture-avoiding substitution closure, closed substitution only, or no substitution closure. \\
\phantomsection\label{entry:symbol-registry-00656}\(M=(I(V),T(V,V'),F(V))\) & \texttt{\detokenize{M,}\allowbreak\detokenize{I,}\allowbreak\detokenize{T,}\allowbreak\detokenize{F,}\allowbreak\detokenize{V,}\allowbreak\detokenize{V'}} & Boolean finite transition-system presentation & P9.5 & meta & P9.5 & \thepage & I and F denote initial and bad states and total T relates current to next state variables in the McMillan interpolation procedure. \\
\phantomsection\label{entry:symbol-registry-00657}\(W_i,R,P,\operatorname{rd}(M,F)\) & \texttt{\detokenize{W_i,}\allowbreak\detokenize{R,}\allowbreak\detokenize{P,}\allowbreak\detokenize{\operatorname{rd}}} & state copies reachable overapproximation interpolant and reverse depth & P9.5 & meta & P9.5 & \thepage & W-i are renamed state blocks R overapproximates reachable states P is the one-state interpolant and rd is the maximum finite shortest distance to the bad set. \\
\phantomsection\label{entry:symbol-registry-00658}\(\mathcal D=(S,P,L),\mathcal M=(S,\mathsf{Act},P,L)\) & \texttt{\detokenize{\mathcal D,}\allowbreak\detokenize{\mathcal M,}\allowbreak\detokenize{S,}\allowbreak\detokenize{P,}\allowbreak\detokenize{L,}\allowbreak\detokenize{\mathsf{Act}}} & finite DTMC and MDP data & P9.5 & n/a & P9.5 & \thepage & P is row-stochastic for the chain and action-indexed for the MDP; only a scheduler-resolved MDP induces one path measure. \\
\phantomsection\label{entry:symbol-registry-00659}\(\mathsf{Act}(s),\sigma:\operatorname{Paths}_{\mathrm{fin}}(\mathcal M)\to\mathsf{Act},\operatorname{Cyl}(\pi)\) & \texttt{\detokenize{\mathsf{Act},}\allowbreak\detokenize{\sigma,}\allowbreak\detokenize{\operatorname{Paths}_{\mathrm{fin}},}\allowbreak\detokenize{\operatorname{Cyl}}} & enabled actions scheduler and finite-path cylinder & P9.5 & meta & P9.5 & \thepage & Act-s is nonempty sigma chooses an enabled action from each finite history and Cyl-pi is the set of infinite paths extending pi; PCTL extrema range over exactly this scheduler class. \\
\phantomsection\label{entry:symbol-registry-00660}\(\Pr_s,\Pr_s^{\min},\Pr_s^{\max}\) & \texttt{\detokenize{\Pr_s,}\allowbreak\detokenize{\Pr_s^{\min},}\allowbreak\detokenize{\Pr_s^{\max}}} & path and extremal scheduler probabilities & P9.5 & meta & P9.5 & \thepage & The undecorated measure belongs to a DTMC or CTMC or to an MDP after fixing a scheduler; min and max quantify over MDP schedulers. \\
\phantomsection\label{entry:symbol-registry-00661}\(\mathsf P_{\bowtie q}[\psi],\mathsf P^{\min}_{\bowtie q}[\psi],\mathsf P^{\max}_{\bowtie q}[\psi]\) & \texttt{\detokenize{\mathsf P_{\bowtie q},}\allowbreak\detokenize{\mathsf P^{\min}_{\bowtie q},}\allowbreak\detokenize{\mathsf P^{\max}_{\bowtie q}}} & probability-threshold modalities & P9.5 & 3 & P9.5 & \thepage & The undecorated form uses the active DTMC or CTMC path measure; min and max forms take the MDP infimum and supremum over the declared scheduler class. \\
\phantomsection\label{entry:symbol-registry-00662}\(\mathsf U^{\leq k},\mathsf X^{I},\mathsf U^{I}\) & \texttt{\detokenize{\mathsf U^{\leq k},}\allowbreak\detokenize{\mathsf X^{I},}\allowbreak\detokenize{\mathsf U^{I}}} & bounded-step and real-time temporal operators & P9.5 & 3.5 & P9.5 & \thepage & The first bounds DTMC or MDP steps by natural k; X-I tests the first CTMC transition time and U-I bounds elapsed time by a declared real interval I. \\
\phantomsection\label{entry:symbol-registry-00663}\(\mathcal C=(S,R,L),R(s,t),E(s),P(s,t),\pi[i],\pi@t,\tau_i\) & \texttt{\detokenize{\mathcal C,}\allowbreak\detokenize{S,}\allowbreak\detokenize{R,}\allowbreak\detokenize{L,}\allowbreak\detokenize{E,}\allowbreak\detokenize{P,}\allowbreak\detokenize{\pi,}\allowbreak\detokenize{\tau}} & finite CTMC rates embedded jumps and timed paths & P9.5 & n/a & P9.5 & \thepage & The nonnegative rational rate map is typed on all state pairs and may contain positive self-loops; E sums the full row P is R divided by E at nonabsorbing states and tau-i is the residence time before path position i-plus-one. Pi-at-t uses the locally stated endpoint convention. \\
\phantomsection\label{entry:symbol-registry-00664}\(\mathsf S_{\bowtie q}[\Phi]\) & \texttt{\detokenize{\mathsf S_{\bowtie q}}} & CSL long-run threshold modality & P9.5 & 3 & P9.5 & \thepage & Compares the starting-state-dependent limiting probability of Phi-states with q; not the temporal since connective. \\
\phantomsection\label{entry:symbol-registry-00665}\((W,\mathcal F,\mu,\pi),w(\varphi)\) & \texttt{\detokenize{W,}\allowbreak\detokenize{\mathcal F,}\allowbreak\detokenize{\mu,}\allowbreak\detokenize{\pi,}\allowbreak\detokenize{w}} & atemporal probability structure and likelihood term & P9.5 & meta & P9.5 & \thepage & Mu is the event measure and w-phi its value on the phi-event; pi is locally the world-valuation map and not the temporal path metavariable. \\
\phantomsection\label{entry:symbol-registry-00666}\(\mathsf{PDL}\) & \texttt{\detokenize{\mathsf{PDL}}} & propositional dynamic logic and calculus identifier & P9.6 & n/a & P9.6 & \thepage & The regular-program modal system whose exact grammar semantics and Hilbert sheet are registered in Part IX. \\
\phantomsection\label{entry:symbol-registry-00667}\(\alpha;\beta,\alpha\cup\beta,\alpha^*,A?\) & \texttt{\detokenize{;,}\allowbreak\detokenize{\cup,}\allowbreak\detokenize{{}^*,}\allowbreak\detokenize{?}} & PDL program constructors & P9.6 & n/a & P9.6 & \thepage & Relational composition choice reflexive-transitive closure and identity test in the mutually generated PDL grammar. \\
\phantomsection\label{entry:symbol-registry-00668}\([\alpha]A\) & \texttt{\detokenize{[\alpha]}} & PDL universal program modality & P9.6 & 3 & P9.6 & \thepage & True when every alpha-successor satisfies A; the index is a regular program rather than a modal frame label. \\
\phantomsection\label{entry:symbol-registry-00669}\(\derives_{\mathsf{PDL}}A\) & \texttt{\detokenize{\derives_{\mathsf{PDL}}}} & PDL theoremhood judgment & P9.6 & meta & P9.6 & \thepage & Finite no-premise derivability on sheet:hilbert-pdl. \\
\phantomsection\label{entry:symbol-registry-00670}\(\{P\}\,C\,\{Q\}\) & \texttt{\detokenize{\{\mathord{\cdot}\}\,}\allowbreak\detokenize{\mathord{\cdot}\,}\allowbreak\detokenize{\{\mathord{\cdot}\}}} & Hoare partial or total correctness triple & P9.6 & meta & P9.6 & \thepage & The active correctness tag determines whether termination is required; braces do not denote sets in this role. \\
\phantomsection\label{entry:symbol-registry-00671}\(\langle C,s\rangle\Downarrow t\) & \texttt{\detokenize{\Downarrow}} & deterministic while-language big-step evaluation & P9.6 & meta & P9.6 & \thepage & Least evaluation relation from a command and initial integer state to a terminating final state. \\
\phantomsection\label{entry:symbol-registry-00672}\(\operatorname{wp}(C,Q),\operatorname{wlp}(C,Q)\) & \texttt{\detokenize{\operatorname{wp},}\allowbreak\detokenize{\operatorname{wlp}}} & weakest and weakest liberal preconditions & P9.6 & meta & P9.6 & \thepage & Predicate transformers distinguished by whether termination is required. \\
\phantomsection\label{entry:symbol-registry-00673}\(\derives_{\mathsf{par}},\derives_{\mathsf{tot}}\) & \texttt{\detokenize{\derives_{\mathsf{par}},}\allowbreak\detokenize{\derives_{\mathsf{tot}}}} & partial and total Hoare derivability tags & P9.6 & meta & P9.6 & \thepage & Judgments on sheet:hoare-while; only the total tag licenses the decreasing-variant loop rule. \\
\phantomsection\label{entry:symbol-registry-00674}\(h_1\perp h_2,h_1\uplus h_2\) & \texttt{\detokenize{\perp,}\allowbreak\detokenize{\uplus}} & disjoint finite heaps and their union & P9.6 & meta & P9.6 & \thepage & Perpendicular denotes disjoint domains and uplus the defined union of compatible finite partial heaps. \\
\phantomsection\label{entry:symbol-registry-00675}\(P*Q\) & \texttt{\detokenize{*}} & separating conjunction & P9.6 & 4 & P9.6 & \thepage & Assertion connective splitting one heap into disjoint subheaps satisfying P and Q; distinct from ordinary conjunction and arithmetic multiplication. \\
\phantomsection\label{entry:symbol-registry-00676}\(\mathit{HB}_P,T_P,M_P\) & \texttt{\detokenize{\mathit{HB}_P,}\allowbreak\detokenize{T_P,}\allowbreak\detokenize{M_P}} & Herbrand base immediate-consequence map and least model & P9.7 & meta & P9.7 & \thepage & HB-P is a set of ground atoms T-P is its monotone consequence operator and M-P is the least fixed point. \\
\phantomsection\label{entry:symbol-registry-00677}\(\leftarrow A_1,\ldots,A_n\) & \texttt{\detokenize{\leftarrow}} & SLD goal & P9.7 & meta & P9.7 & \thepage & Finite ordered goal sequence; an empty sequence is refutation success rather than a clause with an omitted head. \\
\phantomsection\label{entry:symbol-registry-00678}\(\xRightarrow{\theta}\) & \texttt{\detokenize{\xRightarrow{\theta}}} & SLD transition labelled by an mgu & P9.7 & meta & P9.7 & \thepage & Applies the displayed most general unifier to the entire resolvent on sheet:sld-resolution. \\
\phantomsection\label{entry:symbol-registry-00679}\(\mathsf{not}\ A\) & \texttt{\detokenize{\mathsf{not}}} & logic-programming default negation & P9.7 & 3 & P9.7 & \thepage & Ground operational negation by finite failure in SLDNF; not classical object-language negation. \\
\phantomsection\label{entry:symbol-registry-00680}\(\operatorname{Comp}(P)\) & \texttt{\detokenize{\operatorname{Comp}}} & Clark completion of a normal program & P9.7 & meta & P9.7 & \thepage & Predicate-defining biconditionals plus the explicitly stated signature-relative free-term equality axioms; distinct from least and stable models. \\
\phantomsection\label{entry:symbol-registry-00681}\(P^M\) & \texttt{\detokenize{P^M}} & Gelfond-Lifschitz reduct & P9.7 & meta & P9.7 & \thepage & Positive program obtained relative to atom set M by deleting defeated rules and remaining default-negated literals. \\
\phantomsection\label{entry:symbol-registry-00682}\(\Gamma_P,L,G\) & \texttt{\detokenize{\Gamma_P,}\allowbreak\detokenize{L,}\allowbreak\detokenize{G}} & stable operator and well-founded approximation bounds & P9.7 & meta & P9.7 & \thepage & Gamma-P maps M to the least model of P-superscript-M; L and G are the least and greatest fixed points of its square. \\
\phantomsection\label{entry:symbol-registry-00683}\(q(\bar x),D_q\) & \texttt{\detokenize{q,}\allowbreak\detokenize{D_q}} & conjunctive query and canonical database & P9.7 & meta & P9.7 & \thepage & The canonical finite structure D-q carries the query variables constants atoms and distinguished free tuple. \\
\phantomsection\label{entry:symbol-registry-00684}\(q_1\subseteq q_2\) & \texttt{\detokenize{\subseteq}} & conjunctive-query containment & P9.7 & 2 & P9.7 & \thepage & Every answer to q1 is an answer to q2 on every finite database; distinct from set inclusion inside one instance. \\
\phantomsection\label{entry:symbol-registry-00685}\(\operatorname{CSP}(B)\) & \texttt{\detokenize{\operatorname{CSP}}} & fixed-template constraint-satisfaction problem & P9.7 & n/a & P9.7 & \thepage & B is the fixed finite relational target and the input asks for a homomorphism from A to B. \\
\phantomsection\label{entry:symbol-registry-00686}\(\operatorname{Step}_{M,x},\operatorname{Eq},\operatorname{Sel}_b\) & \texttt{\detokenize{\operatorname{Step},}\allowbreak\detokenize{\operatorname{Eq},}\allowbreak\detokenize{\operatorname{Sel}}} & QBF reduction transition equality and selector formulas & P9.8 & meta & P9.8 & \thepage & Metalinguistic polynomial-size Boolean formulas on configuration blocks; the selector permits one textual recursive occurrence. \\
\phantomsection\label{entry:symbol-registry-00687}\(R_i(u,v)\) & \texttt{\detokenize{R_i}} & bounded configuration-reachability QBF family & P9.8 & meta & P9.8 & \thepage & True exactly when v is reachable from u within two-to-the-i steps; a metalinguistic formula family rather than a new object-language predicate. \\
\phantomsection\label{entry:symbol-registry-00688}\(c_{\mathrm{init}}(x),c_{\mathrm{acc}}\) & \texttt{\detokenize{c_{\mathrm{init}},}\allowbreak\detokenize{c_{\mathrm{acc}}}} & initial and canonical accepting configuration blocks & P9.8 & meta & P9.8 & \thepage & Boolean constant blocks substituted into the final closed QBF; not first-order individual constants. \\
\phantomsection\label{entry:symbol-registry-00689}\(\Phi_{M,x}\) & \texttt{\detokenize{\Phi_{M,x}}} & closed QBF produced by the PSPACE reduction & P9.8 & n/a & P9.8 & \thepage & The polynomial-size formula R-m applied to the initial and canonical accepting configuration blocks. \\
\phantomsection\label{entry:symbol-registry-00694}\(\bar E\) & \texttt{\detokenize{\bar}} & NNF complement of an ALC concept & P9.3 & meta & P9.3 & \thepage & The recursive involution complements literals and constants and exchanges Boolean and modal duals; its extension is the complement in the interpretation domain. \\
\phantomsection\label{entry:symbol-registry-00695}\(\mathcal C,m\) & \texttt{\detokenize{\mathcal C,}\allowbreak\detokenize{m}} & ALC input closure and its cardinality & P9.3 & meta & P9.3 & \thepage & Finite set of normalized input concepts their sub-concepts and NNF complements; m is linear in the total input length. \\
\phantomsection\label{entry:symbol-registry-00696}\(t,\operatorname{tp}_I(x)\) & \texttt{\detokenize{t,}\allowbreak\detokenize{\operatorname{tp}}} & complete ALC concept type and realized type & P9.3 & meta & P9.3 & \thepage & A type is a subset of the fixed closure satisfying Boolean completeness and every internalized GCI; the realized type records truth at x in I. \\
\phantomsection\label{entry:symbol-registry-00697}\(t\leadsto_r u\) & \texttt{\detokenize{\leadsto_r}} & role compatibility of ALC types & P9.3 & meta & P9.3 & \thepage & Every universal r-restriction in t has its filler in u; this relation supplies all role edges in the canonical model. \\
\phantomsection\label{entry:symbol-registry-00698}\(S_i,S\) & \texttt{\detokenize{S_i,}\allowbreak\detokenize{S}} & surviving ALC type sets & P9.3 & meta & P9.3 & \thepage & S-zero contains all complete types; each stage removes types lacking existential witnesses; S is the stable survivor set. \\
\phantomsection\label{entry:symbol-registry-00699}\(\mathsf Q_q,\mathsf H_i,\mathsf A_{i,\gamma}\) & \texttt{\detokenize{\mathsf Q,}\allowbreak\detokenize{\mathsf H,}\allowbreak\detokenize{\mathsf A}} & configuration concept names in the ALC hardness reduction & P9.3 & meta & P9.3 & \thepage & Unary concept names record machine control head position and one tape symbol per cell; their extensions are constrained by finite general inclusions. \\
\phantomsection\label{entry:symbol-registry-00700}\(\mathsf B_\ell,Z,L_\ell\) & \texttt{\detokenize{\mathsf B,}\allowbreak\detokenize{Z,}\allowbreak\detokenize{L}} & binary countdown concepts in the ALC hardness reduction & P9.3 & meta & P9.3 & \thepage & B-ell records one rank bit; Z and L-ell abbreviate conjunctions saying all bits or all lower bits are zero. \\
\phantomsection\label{entry:symbol-registry-00701}\(\rho(x)\) & \texttt{\detokenize{\rho}} & nonnegative rank of an ALC reduction model element & P9.3 & meta & P9.3 & \thepage & Integer obtained from the concept-membership bits at x; every transition role edge decreases this value by exactly one. \\
\phantomsection\label{entry:symbol-registry-00702}\(C_{M,w},\mathcal T_{M,s}\) & \texttt{\detokenize{C_{M,w},}\allowbreak\detokenize{\mathcal T_{M,s}}} & ALC concept and TBox encoding alternating acceptance & P9.3 & meta & P9.3 & \thepage & The concept fixes the input configuration and maximal binary rank; the finite TBox enforces configurations branching transitions and countdown. \\
\end{longtable}
\endgroup
\Needspace{.4\textheight}
\subsection*{Part 10 scopes}
\begingroup\scriptsize\setlength{\tabcolsep}{3pt}
\begin{longtable}{@{}p{0.9in}p{1.3in}p{1.0in}p{0.5in}p{0.3in}p{0.4in}p{0.35in}p{1.65in}@{}}
\toprule
\textbf{Glyph} & \textbf{Macro} & \textbf{Role} & \textbf{Scope} & \textbf{Prec} & \textbf{First} & \textbf{Reg. page} & \textbf{Principal uses / notes} \\
\midrule
\endfirsthead
\toprule
\textbf{Glyph} & \textbf{Macro} & \textbf{Role} & \textbf{Scope} & \textbf{Prec} & \textbf{First} & \textbf{Reg. page} & \textbf{Principal uses / notes} \\
\midrule
\endhead
\midrule\multicolumn{8}{r}{\emph{continued}}\\\endfoot
\bottomrule\endlastfoot
\phantomsection\label{entry:symbol-registry-00703}\(T_\tau(X)\) & \texttt{\detokenize{T_\tau}} & term algebra on a generator set & P10.1 & meta & P10.1 & \thepage & Finite terms over the finitary signature tau and variables X; a nonempty algebra only when variables or nullary symbols exist. \\
\phantomsection\label{entry:symbol-registry-00704}\(t^A(v)\) & \texttt{\detokenize{t^A(v)}} & evaluation of an algebraic term & P10.1 & meta & P10.1 & \thepage & Term value in A under the variable assignment v; every assignment extends uniquely to a term-algebra homomorphism. \\
\phantomsection\label{entry:symbol-registry-00705}\(t\sigma,\widehat\sigma\) & \texttt{\detokenize{\widehat\sigma}} & simultaneous algebraic term substitution & P10.1 & meta & P10.1 & \thepage & The map on variables extends homomorphically to terms; the extension may have a different generator alphabet as codomain. \\
\phantomsection\label{entry:symbol-registry-00706}\(\Gamma\models_{\mathrm{alg}}s\approx t\) & \texttt{\detokenize{\models_{\mathrm{alg}}}} & universal identity consequence & P10.1 & meta & P10.1 & \thepage & Every nonempty algebra validating all premise identities under all valuations validates the conclusion identity under all valuations. \\
\phantomsection\label{entry:symbol-registry-00707}\(\Gamma\vdash_{\mathrm{eq}}s\approx t\) & \texttt{\detokenize{\vdash_{\mathrm{eq}}}} & finite equational derivability & P10.1 & meta & P10.1 & \thepage & Derivability from identity premises by the complete registered reflexivity symmetry transitivity compatibility and substitution rules. \\
\phantomsection\label{entry:symbol-registry-00708}\(\theta_\Gamma\) & \texttt{\detokenize{\theta_\Gamma}} & congruence of provable identities & P10.1 & meta & P10.1 & \thepage & The least fully invariant congruence on the nonempty term algebra containing all premise equation pairs. \\
\phantomsection\label{entry:symbol-registry-00709}\(F_{\mathcal V}(X),\eta_X\) & \texttt{\detokenize{F_{\mathcal V},}\allowbreak\detokenize{\eta_X}} & free algebra and its generator map & P10.1 & meta & P10.1 & \thepage & The map from X has a unique homomorphic extension into every member of the variety; it need not be injective in a trivial variety. \\
\phantomsection\label{entry:symbol-registry-00710}\(\theta_{\mathcal V,X}\) & \texttt{\detokenize{\theta_{\mathcal V,X}}} & semantic congruence defining a relative free algebra & P10.1 & meta & P10.1 & \thepage & Two X-terms are related exactly when all evaluations into all algebras of the variety identify them. \\
\phantomsection\label{entry:symbol-registry-00711}\(\operatorname{Con}A,\Delta_A,\nabla_A\) & \texttt{\detokenize{\operatorname{Con},}\allowbreak\detokenize{\Delta,}\allowbreak\detokenize{\nabla}} & congruence lattice and its bounds & P10.1 & meta & P10.1 & \thepage & Congruences ordered by inclusion; Delta is equality and nabla is the universal relation; meets are intersections. \\
\phantomsection\label{entry:symbol-registry-00712}\(\operatorname{Cg}_A(a,b)\) & \texttt{\detokenize{\operatorname{Cg}}} & principal congruence generated by a pair & P10.1 & meta & P10.1 & \thepage & Intersection of every congruence on A containing the pair (a,b). \\
\phantomsection\label{entry:symbol-registry-00713}\(\ker h\) & \texttt{\detokenize{\ker}} & kernel congruence of an algebra homomorphism & P10.1 & meta & P10.1 & \thepage & Pairs of domain elements with equal h-images; the kernel of a product map is the intersection of its coordinate kernels. \\
\phantomsection\label{entry:symbol-registry-00714}\(\phi/\theta\) & \texttt{\detokenize{/}} & quotient congruence above a kernel & P10.1 & meta & P10.1 & \thepage & For theta contained in phi this relation on A/theta relates exactly the classes of phi-related representatives. \\
\phantomsection\label{entry:symbol-registry-00715}\(\mu\) & \texttt{\detokenize{\mu}} & monolith of a nontrivial algebra & P10.1 & meta & P10.1 & \thepage & Least congruence strictly above equality when it exists; this condition characterizes nontrivial subdirect irreducibility. \\
\phantomsection\label{entry:symbol-registry-00716}\(\theta_{a,b}\) & \texttt{\detokenize{\theta_{a,b}}} & maximal congruence separating a chosen pair & P10.1 & meta & P10.1 & \thepage & Maximal among congruences omitting (a,b); the quotient is subdirectly irreducible and still separates the chosen elements. \\
\phantomsection\label{entry:symbol-registry-00717}\(F[a]\) & \texttt{\detokenize{F[a]}} & Boolean filter generated by adjoining one element & P10.2 & meta & P10.2 & \thepage & Elements above f meet a for some f in F; used in ultrafilter dichotomy and avoiding-element extension. \\
\phantomsection\label{entry:symbol-registry-00718}\(X_B,\widehat a\) & \texttt{\detokenize{X_B,}\allowbreak\detokenize{\widehat}} & Stone space and its basic clopen sets & P10.2 & meta & P10.2 & \thepage & X-B consists of the ultrafilters of B; hat-a consists of those containing a; trivial B gives the empty space. \\
\phantomsection\label{entry:symbol-registry-00719}\(\operatorname{Clop}(X)\) & \texttt{\detokenize{\operatorname{Clop}}} & Boolean algebra of clopen subsets & P10.2 & meta & P10.2 & \thepage & Intersection union and relative complement act on the clopens of X; the bounds are the empty set and X. \\
\phantomsection\label{entry:symbol-registry-00720}\(\beta_B\) & \texttt{\detokenize{\beta_B}} & Boolean representation isomorphism & P10.2 & meta & P10.2 & \thepage & Maps a Boolean element a to its basic clopen hat-a in the Stone space X-B. \\
\phantomsection\label{entry:symbol-registry-00721}\(\varepsilon_X\) & \texttt{\detokenize{\varepsilon_X}} & Stone-space reconstruction homeomorphism & P10.2 & meta & P10.2 & \thepage & Maps x to the ultrafilter of all clopen subsets containing x. \\
\phantomsection\label{entry:symbol-registry-00722}\(X(h)\) & \texttt{\detokenize{X(h)}} & continuous dual of a Boolean homomorphism & P10.2 & meta & P10.2 & \thepage & For h from B to C sends an ultrafilter of C to its inverse image under h, an ultrafilter of B. \\
\phantomsection\label{entry:symbol-registry-00723}\(\operatorname{Clop}(g)\) & \texttt{\detokenize{\operatorname{Clop}(g)}} & Boolean dual of a continuous map & P10.2 & meta & P10.2 & \thepage & For g from X to Y sends a clopen subset of Y to its inverse image in X. \\
\phantomsection\label{entry:symbol-registry-00724}\(\mathsf{BA},\mathsf{Stone}\) & \texttt{\detokenize{\mathsf{BA},}\allowbreak\detokenize{\mathsf{Stone}}} & categories in Stone duality & P10.2 & meta & P10.2 & \thepage & Boolean algebras with homomorphisms preserving both bounds; compact Hausdorff spaces with clopen bases and continuous maps; trivial and empty objects are allowed. \\
\phantomsection\label{entry:symbol-registry-00725}\(K=2^{\mathbb N}\) & \texttt{\detokenize{K,}\allowbreak\detokenize{2^{\mathbb N}}} & Cantor space of binary sequences & P10.2 & meta & P10.2 & \thepage & The countable product of the discrete two-element space with coordinates starting at zero. \\
\phantomsection\label{entry:symbol-registry-00726}\([s],x|_n\) & \texttt{\detokenize{[s],}\allowbreak\detokenize{|_n}} & binary cylinder and length-n initial segment & P10.2 & meta & P10.2 & \thepage & The cylinder contains precisely the sequences extending the finite word s; the empty word gives all Cantor space. \\
\phantomsection\label{entry:symbol-registry-00727}\(\mathcal B_n,\mathcal B_{\rm cyl}\) & \texttt{\detokenize{\mathcal B_n,}\allowbreak\detokenize{\mathcal B_{\rm cyl}}} & finite cylinder algebras and their union & P10.2 & meta & P10.2 & \thepage & B-n contains all unions of length-n cylinders; the union is the clopen algebra of Cantor space and the direct limit of the inclusions. \\
\phantomsection\label{entry:symbol-registry-00728}\(H:K\to X_B\) & \texttt{\detokenize{H}} & explicit Cantor-to-Stone homeomorphism & P10.2 & meta & P10.2 & \thepage & Using a Boolean isomorphism alpha from B to the cylinder algebra sends x to the elements b for which x belongs to alpha(b). \\
\phantomsection\label{entry:symbol-registry-00729}\(U_p\) & \texttt{\detokenize{U_p}} & principal ultrafilter of a Boolean atom & P10.2 & meta & P10.2 & \thepage & All elements above atom p; its singleton is the basic clopen hat-p and hence an isolated Stone point. \\
\phantomsection\label{entry:symbol-registry-00730}\(\Lambda,n_f\) & \texttt{\detokenize{\Lambda,}\allowbreak\detokenize{n_f}} & finitary Boolean-operator signature and arity & P10.3 & meta & P10.3 & \thepage & A set of additional symbols with one finite arity per symbol; nullary symbols designate arbitrary Boolean elements. \\
\phantomsection\label{entry:symbol-registry-00731}\(f^B\) & \texttt{\detokenize{f^B}} & abstract Boolean operator & P10.3 & meta & P10.3 & \thepage & For positive arity preserves zero and finite joins in each coordinate; the conditions are vacuous for nullary symbols. \\
\phantomsection\label{entry:symbol-registry-00732}\(\mathfrak F^+\) & \texttt{\detokenize{\mathfrak F^+}} & full complex algebra of a relational frame & P10.3 & meta & P10.3 & \thepage & Power-set Boolean algebra with existential successor-tuple operations; the carrier may be the power set of the empty frame. \\
\phantomsection\label{entry:symbol-registry-00733}\(R_f(U,V_1,\ldots,V_n)\) & \texttt{\detokenize{R_f}} & canonical relation of a finitary Boolean operator & P10.3 & meta & P10.3 & \thepage & Every tuple of elements from the successor ultrafilters has its operator value in U; for arity zero this is a unary predicate. \\
\phantomsection\label{entry:symbol-registry-00734}\(\mathbf B^\bullet\) & \texttt{\detokenize{\mathbf B^\bullet}} & ultrafilter frame of a Boolean algebra with operators & P10.3 & meta & P10.3 & \thepage & The set of Boolean ultrafilters with all canonical operator relations. \\
\phantomsection\label{entry:symbol-registry-00735}\(\mathbf B^\sigma\) & \texttt{\detokenize{\mathbf B^\sigma}} & canonical extension of a Boolean algebra with operators & P10.3 & meta & P10.3 & \thepage & The full complex algebra of the ultrafilter frame; the Boolean embedding is dense and compact. \\
\phantomsection\label{entry:symbol-registry-00736}\(\mathbf B_*\) & \texttt{\detokenize{\mathbf B_*}} & descriptive general frame of a Boolean algebra with operators & P10.3 & meta & P10.3 & \thepage & The canonical relations on the Stone space together with precisely the basic-clopen image as admissible algebra. \\
\phantomsection\label{entry:symbol-registry-00737}\(\mathcal G^+,\mathcal A\) & \texttt{\detokenize{\mathcal G^+,}\allowbreak\detokenize{\mathcal A}} & admissible algebra of a general frame & P10.3 & meta & P10.3 & \thepage & The carrier is the Boolean family of admissible sets closed under complex operators. \\
\phantomsection\label{entry:symbol-registry-00738}\(h_*\) & \texttt{\detokenize{h_*}} & bounded morphism dual to a Boolean-operator homomorphism & P10.3 & meta & P10.3 & \thepage & Inverse image on ultrafilters; its back condition follows from the simultaneous filter-extension lemma. \\
\phantomsection\label{entry:symbol-registry-00739}\(p_U=\bigwedge e[U]\) & \texttt{\detokenize{p_U}} & atom in a dense compact Boolean completion & P10.3 & meta & P10.3 & \thepage & The meet of the images of all elements of an ultrafilter U; compactness makes it nonzero and density makes it an atom. \\
\phantomsection\label{entry:symbol-registry-00740}\(\Box_A,\Diamond_A\) & \texttt{\detokenize{\Box_A,}\allowbreak\detokenize{\Diamond_A}} & dual operations of a normal modal algebra & P10.3 & meta & P10.3 & \thepage & Box preserves finite meets and one; diamond is its Boolean dual and preserves finite joins and zero. \\
\phantomsection\label{entry:symbol-registry-00741}\(\varphi\equiv_L\psi,[\varphi]_L\) & \texttt{\detokenize{\equiv_L,}\allowbreak\detokenize{[\cdot]_L}} & normal-modal provable equivalence and its classes & P10.3 & meta & P10.3 & \thepage & Equivalence means theoremhood of the biconditional in L and is a congruence for all formula constructors. \\
\phantomsection\label{entry:symbol-registry-00742}\(F_L\) & \texttt{\detokenize{F_L}} & normal modal Lindenbaum algebra & P10.3 & meta & P10.3 & \thepage & The formula algebra modulo provable equivalence; arbitrary valuations are handled with finite representative selections per formula. \\
\phantomsection\label{entry:symbol-registry-00743}\(\operatorname{Alg}(L)\) & \texttt{\detokenize{\operatorname{Alg}}} & algebra class validating a normal modal logic & P10.3 & meta & P10.3 & \thepage & Normal modal algebras validating every member of L under all algebra valuations. \\
\phantomsection\label{entry:symbol-registry-00744}\(\mathfrak G_L,\mathcal A_L\) & \texttt{\detokenize{\mathfrak G_L,}\allowbreak\detokenize{\mathcal A_L}} & canonical general frame and its admissible sets & P10.3 & meta & P10.3 & \thepage & Ultrafilters of the modal Lindenbaum algebra with the canonical relation and Stone clopens as admissible variable values. \\
\phantomsection\label{entry:symbol-registry-00745}\(\mathfrak F_L,X_L,R_L\) & \texttt{\detokenize{\mathfrak F_L,}\allowbreak\detokenize{X_L,}\allowbreak\detokenize{R_L}} & underlying canonical relational frame & P10.3 & meta & P10.3 & \thepage & Same worlds and relation as the canonical general frame; ordinary valuations may use every subset of X-L. \\
\phantomsection\label{entry:symbol-registry-00746}\(\operatorname{Gen}(L),\operatorname{Fr}(L)\) & \texttt{\detokenize{\operatorname{Gen},}\allowbreak\detokenize{\operatorname{Fr}}} & nonempty validating general-frame and ordinary-frame classes & P10.3 & meta & P10.3 & \thepage & General frames validate L using admissible valuations; ordinary frames validate it using all subset valuations. \\
\phantomsection\label{entry:symbol-registry-00747}\(L(Y)\) & \texttt{\detokenize{L(Y)}} & schematic extension of a modal logic to another variable set & P10.3 & meta & P10.3 & \thepage & The least normal logic over Y containing all substituted instances of the original logic; used to state all-variable-set canonicity. \\
\phantomsection\label{entry:symbol-registry-00748}\(\operatorname{Fm}_\Lambda(X)\) & \texttt{\detokenize{\operatorname{Fm}_\Lambda}} & formulas for a finitary modal signature & P10.3 & meta & P10.3 & \thepage & Finite Boolean formulas over X with one constructor of arity n-f for every symbol f in Lambda. \\
\phantomsection\label{entry:symbol-registry-00749}\(f_i(\chi)\) & \texttt{\detokenize{f_i}} & one-coordinate modal formula context & P10.3 & meta & P10.3 & \thepage & All other arguments are fixed; used in theorem-level monotonicity and the normality and additivity axiom schemes. \\
\phantomsection\label{entry:symbol-registry-00750}\(F_{\Lambda,L}\) & \texttt{\detokenize{F_{\Lambda,L}}} & Lindenbaum algebra of a finitary normal logic & P10.3 & meta & P10.3 & \thepage & The formula quotient by L-provable equivalence with all finite-arity and nullary operators interpreted on classes. \\
\phantomsection\label{entry:symbol-registry-00751}\(\operatorname{Alg}_\Lambda(L)\) & \texttt{\detokenize{\operatorname{Alg}_\Lambda}} & Boolean-operator algebras validating a finitary logic & P10.3 & meta & P10.3 & \thepage & All Lambda-BAOs validating L under every algebra valuation, including trivial algebras. \\
\phantomsection\label{entry:symbol-registry-00752}\(\operatorname{Ult}(A)\) & \texttt{\detokenize{\operatorname{Ult}}} & ultrafilters of the Boolean reduct & P10.3 & meta & P10.3 & \thepage & The set of proper maximal Boolean filters, used as the canonical frame carrier of A. \\
\phantomsection\label{entry:symbol-registry-00753}\(\mathcal D_L,W_L\) & \texttt{\detokenize{\mathcal D_L,}\allowbreak\detokenize{W_L}} & descriptive canonical frame of a finitary normal logic & P10.3 & meta & P10.3 & \thepage & Worlds are ultrafilters of F-Lambda-L; the relations are the full canonical operator relations and the admissible sets are Stone clopens. \\
\phantomsection\label{entry:symbol-registry-00754}\(\operatorname{DGen}_\Lambda(L)\) & \texttt{\detokenize{\operatorname{DGen}_\Lambda}} & descriptive general frames validating a finitary logic & P10.3 & meta & P10.3 & \thepage & The nonempty descriptive Lambda-general frames validating every theorem under all admissible valuations. \\
\phantomsection\label{entry:symbol-registry-00755}\(v_0(p)=\widehat{[p]_L}\) & \texttt{\detokenize{v_0}} & canonical modal valuation & P10.3 & meta & P10.3 & \thepage & Assigns each proposition variable the basic clopen determined by its formula class. \\
\phantomsection\label{entry:symbol-registry-00756}\(X_D,\widehat a\) & \texttt{\detokenize{X_D,}\allowbreak\detokenize{\widehat}} & Priestley spectrum and its basic clopen upsets & P10.3 & meta & P10.3 & \thepage & Proper prime lattice filters ordered by inclusion; hat-a consists of those containing a, and the topology also uses their complements. \\
\phantomsection\label{entry:symbol-registry-00757}\(\operatorname{ClopUp}(X)\) & \texttt{\detokenize{\operatorname{ClopUp}}} & lattice of clopen upsets & P10.3 & meta & P10.3 & \thepage & Bounded distributive lattice for a Priestley space; carries the residual complement-of-downclosure operation for an Esakia space. \\
\phantomsection\label{entry:symbol-registry-00758}\(\eta_D\) & \texttt{\detokenize{\eta_D}} & Priestley lattice representation & P10.3 & meta & P10.3 & \thepage & Maps each lattice element to the clopen upset of prime filters containing it. \\
\phantomsection\label{entry:symbol-registry-00759}\(\varepsilon_X:X\to X_{\operatorname{ClopUp}(X)}\) & \texttt{\detokenize{\varepsilon_X}} & ordered-space reconstruction by prime filters & P10.3 & meta & P10.3 & \thepage & Maps a point to its membership filter of clopen upsets; an order isomorphism and homeomorphism. \\
\phantomsection\label{entry:symbol-registry-00760}\(\Pi_D,p_a\) & \texttt{\detokenize{\Pi_D,}\allowbreak\detokenize{p_a}} & classical encoding of prime lattice filters & P10.3 & meta & P10.3 & \thepage & The proposition p-a records membership of a; Pi-D imposes only bounded-lattice homomorphism laws to two, not Heyting implication preservation. \\
\phantomsection\label{entry:symbol-registry-00761}\(\mathsf{DC}\) & \texttt{\detokenize{\mathsf{DC}}} & dependent choice for a prescribed start & P10.3 & meta & P10.3 & \thepage & A serial relation on a nonempty set admits an infinite path beginning at a prescribed point; used in the arbitrary-poset dense-set lemma. \\
\phantomsection\label{entry:symbol-registry-00762}\((p_n),\mathcal D\) & \texttt{\detokenize{p_n,}\allowbreak\detokenize{\mathcal D}} & descending conditions and dense requirements & P10.3 & meta & P10.3 & \thepage & Rasiowa-Sikorski construction has p-n-plus-one stronger than p-n and in D-n; its upward closure is generic for the given family. \\
\phantomsection\label{entry:symbol-registry-00763}\(a\Rightarrow b\) & \texttt{\detokenize{\Rightarrow}} & Heyting residual on algebra elements & P10.3 & meta & P10.3 & \thepage & The greatest c with c meet a below b; source notation a-arrow-b denotes this same operation when operands are algebra elements. \\
\phantomsection\label{entry:symbol-registry-00764}\(\mathcal O(X)\) & \texttt{\detokenize{\mathcal O}} & Heyting algebra of open subsets & P10.3 & meta & P10.3 & \thepage & Finite meets are intersections and implication is the interior of the union of the first operand's complement with the second. \\
\phantomsection\label{entry:symbol-registry-00765}\(\operatorname{Up}(W)\) & \texttt{\detokenize{\operatorname{Up}}} & Heyting algebra of upsets of a preorder & P10.3 & meta & P10.3 & \thepage & Implication holds at w when every larger point in the antecedent upset belongs to the consequent upset. \\
\phantomsection\label{entry:symbol-registry-00766}\(H_{\mathsf{IPC}},[\varphi]\) & \texttt{\detokenize{H_{\mathsf{IPC}},}\allowbreak\detokenize{[\cdot]}} & IPC Lindenbaum Heyting algebra & P10.3 & meta & P10.3 & \thepage & Formula classes under mutual IPL derivability, ordered by one-premise derivability. \\
\phantomsection\label{entry:symbol-registry-00767}\(\operatorname{Pf}(H)\) & \texttt{\detokenize{\operatorname{Pf}}} & prime-filter order of a Heyting algebra & P10.3 & meta & P10.3 & \thepage & All proper prime lattice filters ordered by inclusion; filters need not be maximal. \\
\phantomsection\label{entry:symbol-registry-00768}\(\beta_H(a)\) & \texttt{\detokenize{\beta_H}} & Heyting embedding into upsets of prime filters & P10.3 & meta & P10.3 & \thepage & The set of prime filters containing a; the implication witness proves preservation of the residual. \\
\phantomsection\label{entry:symbol-registry-00769}\(\downarrow S,\uparrow x\) & \texttt{\detokenize{\downarrow,}\allowbreak\detokenize{\uparrow}} & downward closure and principal upset & P10.3 & meta & P10.3 & \thepage & Order closures in a Priestley or Esakia space; the downclosure condition applies to every clopen subset. \\
\phantomsection\label{entry:symbol-registry-00770}\(g[\uparrow x]=\uparrow g(x)\) & \texttt{\detokenize{g[\uparrow x]}} & Esakia order back condition & P10.3 & meta & P10.3 & \thepage & Together with continuity and monotonicity this characterizes Esakia morphisms and ensures implication-preserving inverse images. \\
\phantomsection\label{entry:symbol-registry-00771}\(\mathsf{HA},\mathsf{BA}\) & \texttt{\detokenize{\mathsf{HA},}\allowbreak\detokenize{\mathsf{BA}}} & Heyting and Boolean classes in the common residual signature & P10.3 & meta & P10.3 & \thepage & Classes of algebras for the intermediate-logic correspondence; HA here is not Heyting arithmetic from Part VII. \\
\phantomsection\label{entry:symbol-registry-00772}\(H_L\) & \texttt{\detokenize{H_L}} & Lindenbaum Heyting algebra of an intermediate logic & P10.3 & meta & P10.3 & \thepage & The quotient by L-provable biconditional, with formula implication inducing the algebraic residual. \\
\phantomsection\label{entry:symbol-registry-00773}\(\mathsf V(L),\operatorname{Log}(\mathcal V)\) & \texttt{\detokenize{\mathsf V,}\allowbreak\detokenize{\operatorname{Log}}} & order-reversing logic-variety assignments & P10.3 & meta & P10.3 & \thepage & Heyting algebras validating L and formulas valid in V; inverse on the intervals IPC-to-CPL and BA-to-HA. \\
\phantomsection\label{entry:symbol-registry-00774}\(s^\flat,\varphi^{\rm alg}\) & \texttt{\detokenize{{}^\flat,}\allowbreak\detokenize{{}^{\rm alg}}} & translation between residual terms and formulas & P10.3 & meta & P10.3 & \thepage & Converts algebraic zero one and residual into bottom top and implication or conversely; identities correspond to biconditionals with value one. \\
\phantomsection\label{entry:symbol-registry-00775}\(\mathcal C(A,B)\) & \texttt{\detokenize{\mathcal C(A,}\allowbreak\detokenize{B)}} & categorical hom-set & P10.7-P10.9 & meta & P10.7 & \thepage & Arrows from object A to object B; local smallness places each such set in the fixed universe. \\
\phantomsection\label{entry:symbol-registry-00776}\(\mathcal C^{\mathrm{op}}\) & \texttt{\detokenize{{}^{\mathrm{op}}}} & opposite category & P10.7-P10.9 & meta & P10.7 & \thepage & Reverses the arrows and composition order; a contravariant functor is an ordinary functor out of this category. \\
\phantomsection\label{entry:symbol-registry-00777}\(1_A,1_{\mathcal C}\) & \texttt{\detokenize{1_A,}\allowbreak\detokenize{1_{\mathcal C}}} & identity arrow and identity functor & P10.7-P10.9 & meta & P10.7 & \thepage & The subscript states the categorical level: an object for an arrow and a category for its identity functor. \\
\phantomsection\label{entry:symbol-registry-00778}\(\operatorname{Nat}(F,G)\) & \texttt{\detokenize{\operatorname{Nat}}} & family of natural transformations & P10.7-P10.9 & meta & P10.7 & \thepage & Component families in the containing ambient universe; Yoneda gives a small set of parameters for transformations from a representable. \\
\phantomsection\label{entry:symbol-registry-00779}\(F\dashv G,\Phi_{A,B}\) & \texttt{\detokenize{\dashv,}\allowbreak\detokenize{\Phi_{A,B}}} & adjunction and hom-set transposition & P10.7-P10.9 & meta & P10.7 & \thepage & Phi sends an arrow from F(A) to B to an arrow from A to G(B); its bijections are natural in both arguments. \\
\phantomsection\label{entry:symbol-registry-00780}\(\eta,\varepsilon\) & \texttt{\detokenize{\eta,}\allowbreak\detokenize{\varepsilon}} & unit and counit of an adjunction & P10.7-P10.9 & meta & P10.7 & \thepage & Natural transformations from the identity to GF and from FG to the identity; their components satisfy the two triangle identities. \\
\phantomsection\label{entry:symbol-registry-00781}\(\pi_i,\iota_i,\langle f_i\rangle_i\) & \texttt{\detokenize{\pi_i,}\allowbreak\detokenize{\iota_i,}\allowbreak\detokenize{\langle\rangle}} & product projections coproduct injections and pairing & P10.7-P10.9 & meta & P10.7 & \thepage & Domains and codomains distinguish projections from injections; pairing is the unique arrow with its specified product components. \\
\phantomsection\label{entry:symbol-registry-00782}\(f\times g\) & \texttt{\detokenize{\times}} & product of categorical arrows & P10.7-P10.9 & meta & P10.7 & \thepage & The arrow with components f after the first projection and g after the second; respects identity and composition. \\
\phantomsection\label{entry:symbol-registry-00783}\(B^A,\operatorname{ev}_{A,B},\Lambda(f)\) & \texttt{\detokenize{{}^{A},}\allowbreak\detokenize{\operatorname{ev},}\allowbreak\detokenize{\Lambda}} & exponential object evaluation and transpose & P10.7-P10.9 & meta & P10.7 & \thepage & Evaluation has domain \(B^A\times A\); transposition sends \(f:X\times A\to B\) to the unique corresponding arrow \(X\to B^A\). \\
\phantomsection\label{entry:symbol-registry-00784}\(\mathcal S_{\beta\eta}\) & \texttt{\detokenize{\mathcal S_{\beta\eta}}} & category of arrow-only typed contexts & P10.8 & meta & P10.8 & \thepage & Objects are finite type lists and arrows are typed term tuples modulo componentwise beta-eta equality. \\
\phantomsection\label{entry:symbol-registry-00785}\(\Theta\vdash_{\beta\eta}M=N:\alpha\) & \texttt{\detokenize{\vdash_{\beta\eta}}} & typed beta-eta equational consequence & P10.8 & meta & P10.8 & \thepage & The equational metatheory includes alpha-renaming and typed beta and eta schemas; the registered reduction relation remains beta reduction. \\
\phantomsection\label{entry:symbol-registry-00786}\(\sigma:\Theta\to\Psi,M[\sigma]\) & \texttt{\detokenize{\sigma,}\allowbreak\detokenize{M[\sigma]}} & typed simultaneous substitution and its action & P10.8 & meta & P10.8 & \thepage & A tuple of Theta-terms with the types declared in Psi; substitutes into a Psi-term without resubstituting into inserted terms. \\
\phantomsection\label{entry:symbol-registry-00787}\(c_\Gamma(\beta),\Delta^\Gamma\) & \texttt{\detokenize{c_\Gamma,}\allowbreak\detokenize{\Delta^\Gamma}} & curried type and context exponential & P10.8 & meta & P10.8 & \thepage & Right-associated arrows over Gamma's type list; the exponential list has one such curried type per component of Delta. \\
\phantomsection\label{entry:symbol-registry-00788}\(\llbracket\Gamma\rrbracket,\llbracket M\rrbracket\) & \texttt{\detokenize{\llbracket\rrbracket}} & categorical interpretation of a context or term & P10.8 & meta & P10.8 & \thepage & Context products begin with the chosen terminal object; a term denotes an arrow from its context object to its interpreted type. \\
\phantomsection\label{entry:symbol-registry-00789}\(\underline A,A_T,P_\Gamma\) & \texttt{\detokenize{\underline,}\allowbreak\detokenize{A_T,}\allowbreak\detokenize{P_\Gamma}} & named type interpreted type and context object & P10.8 & meta & P10.8 & \thepage & The named base type denotes \(A\); \(A_T\) interprets the type \(T\), and \(P_\Gamma\) is the recursively chosen product with leading terminal factor. \\
\phantomsection\label{entry:symbol-registry-00790}\(c_f,u,\mathsf{pair}_{A,B}\) & \texttt{\detokenize{c_f,}\allowbreak\detokenize{u,}\allowbreak\detokenize{\mathsf{pair}}} & named arrow terminal and pairing constants & P10.8 & meta & P10.8 & \thepage & Closed constants in the expanded internal language; their declared applications denote the corresponding arrows and product pairing. \\
\phantomsection\label{entry:symbol-registry-00791}\(e_{A,B},l_{A,B}\) & \texttt{\detokenize{e_{A,B},}\allowbreak\detokenize{l_{A,B}}} & comparison constants for a named exponential & P10.8 & meta & P10.8 & \thepage & Compare the base type naming \(B^A\) with the arrow type between the named types for \(A\) and \(B\); both unary interpretations are identities. \\
\phantomsection\label{entry:symbol-registry-00792}\(\operatorname{Diag}(\mathcal C),=_{\mathcal C}\) & \texttt{\detokenize{\operatorname{Diag},}\allowbreak\detokenize{=_{\mathcal C}}} & full true diagram and its equational theory & P10.8 & meta & P10.8 & \thepage & Contains every true typed equation in every finite context of the named language; no effective enumeration is asserted. \\
\phantomsection\label{entry:symbol-registry-00793}\(\kappa_T,\iota_T\) & \texttt{\detokenize{\kappa_T,}\allowbreak\detokenize{\iota_T}} & recursive type comparison terms & P10.8 & meta & P10.8 & \thepage & Closed inverse comparisons between T and the named type for its interpreted object; applications preserve denotation. \\
\phantomsection\label{entry:symbol-registry-00794}\(\operatorname{pack}_\Gamma,\operatorname{unpack}_\Gamma\) & \texttt{\detokenize{\operatorname{pack},}\allowbreak\detokenize{\operatorname{unpack}}} & inverse context packing arrows & P10.8 & meta & P10.8 & \thepage & Compare \(\Gamma\) with the singleton named context for \(P_\Gamma\); the packing term and its context arrow have different codomains. \\
\phantomsection\label{entry:symbol-registry-00795}\(\mathcal S_{\mathcal C},E,J\) & \texttt{\detokenize{\mathcal S_{\mathcal C},}\allowbreak\detokenize{E,}\allowbreak\detokenize{J}} & internal context category evaluation and quasi-inverse & P10.8 & meta & P10.8 & \thepage & Evaluation sends \(\Gamma\) to \(P_\Gamma\); \(J\) names objects and arrows, with \(EJ(A)=1\times A\) for the chosen product. \\
\phantomsection\label{entry:symbol-registry-00796}\(j_{\Gamma,\Delta},\chi_{\Gamma,\Delta}\) & \texttt{\detokenize{j_{\Gamma,\Delta},}\allowbreak\detokenize{\chi_{\Gamma,\Delta}}} & product and exponential comparison isomorphisms & P10.8 & meta & P10.8 & \thepage & The product comparison uses interpreted block projections; the exponential comparison transposes interpreted evaluation after the inverse product comparison. \\
\phantomsection\label{entry:symbol-registry-00797}\(\mathbb F,j:\mathbb F^{\mathrm{op}}\to T\) & \texttt{\detokenize{\mathbb F,}\allowbreak\detokenize{j}} & finite ordinals and the structural theory functor & P10.7 & meta & P10.7 & \thepage & The identity-on-objects functor preserves finite products; a finite-set function \(m\to n\) selects an arrow \(n\to m\) in the theory. \\
\phantomsection\label{entry:symbol-registry-00798}\(p_i^n,\langle f_i\rangle_{i<n},!_n\) & \texttt{\detokenize{p_i^n,}\allowbreak\detokenize{\langle\rangle,}\allowbreak\detokenize{!_n}} & chosen powers and tuples in a Lawvere theory & P10.7 & meta & P10.7 & \thepage & The arrows \(p_i^n:n\to1\) exhibit \(n=1^n\), with \(p_0^1=1_1\); the empty tuple is the unique arrow \(n\to0\). \\
\phantomsection\label{entry:symbol-registry-00799}\((k,s,t),E\) & \texttt{\detokenize{(k,}\allowbreak\detokenize{s,}\allowbreak\detokenize{t),}\allowbreak\detokenize{E}} & equations with explicit finite contexts & P10.7 & meta & P10.7 & \thepage & Each triple asserts equality of its two terms under every assignment to the declared \(k\)-variable context. \\
\phantomsection\label{entry:symbol-registry-00800}\(\operatorname{Alg}_{\mathrm{all}}(\Sigma,E),\operatorname{Alg}_{\mathrm{ne}}(\Sigma,E)\) & \texttt{\detokenize{\operatorname{Alg}_{\mathrm{all}},}\allowbreak\detokenize{\operatorname{Alg}_{\mathrm{ne}}}} & all-carrier and nonempty-carrier equational categories & P10.7 & meta & P10.7 & \thepage & Both use carriers in the fixed universe and all operation-preserving functions; the latter is the full subcategory with nonempty carrier. \\
\phantomsection\label{entry:symbol-registry-00801}\(\equiv_E^n\) & \texttt{\detokenize{\equiv_E^n}} & equational congruence in a finite context & P10.7 & meta & P10.7 & \thepage & The least family containing the stated equations and closed under operation compatibility and simultaneous substitutions between all finite contexts. \\
\phantomsection\label{entry:symbol-registry-00802}\(T_{\Sigma,E}(n,m)\) & \texttt{\detokenize{T_{\Sigma,E}}} & term theory of an equational presentation & P10.7 & meta & P10.7 & \thepage & Its arrows are \(m\)-tuples of congruence classes of \(n\)-variable terms; composition is substitution and object \(0\) is terminal. \\
\phantomsection\label{entry:symbol-registry-00803}\(M_A,M_h\) & \texttt{\detokenize{M_A,}\allowbreak\detokenize{M_h}} & evaluation model and homomorphism transformation & P10.7 & meta & P10.7 & \thepage & The model assigns \(A^n\) to \(n\) and evaluates term tuples; a homomorphism \(h\) gives components \(h^n\). \\
\phantomsection\label{entry:symbol-registry-00804}\(\operatorname{Mod}(T),U\) & \texttt{\detokenize{\operatorname{Mod},}\allowbreak\detokenize{U}} & Lawvere models and the underlying-set functor & P10.7 & meta & P10.7 & \thepage & Finite-product-preserving functors to Set and all natural transformations; \(U(M)=M(1)\) and \(U(\alpha)=\alpha_1\). \\
\phantomsection\label{entry:symbol-registry-00805}\(c_n^M\) & \texttt{\detokenize{c_n^M}} & canonical model power comparison & P10.7 & meta & P10.7 & \thepage & The bijection \(M(n)\to M(1)^n\) has coordinates \(M(p_i^n)\); the comparison at one uses singleton tuples. \\
\phantomsection\label{entry:symbol-registry-00806}\(\bar t,\bar\sigma\) & \texttt{\detokenize{\bar t,}\allowbreak\detokenize{\bar\sigma}} & one-component term and basic-operation arrows & P10.7 & meta & P10.7 & \thepage & An \(n\)-variable term determines an arrow \(n\to1\); a symbol of arity \(k\) gives its basic term arrow \(k\to1\). \\
\phantomsection\label{entry:symbol-registry-00807}\(\Sigma_{\mathcal T},E_{\mathcal T},\sigma_f\) & \texttt{\detokenize{\Sigma_{\mathcal T},}\allowbreak\detokenize{E_{\mathcal T},}\allowbreak\detokenize{\sigma_f}} & named equational presentation of a theory & P10.7 & meta & P10.7 & \thepage & One operation symbol for each arrow to the generator; projection and substitution equations recover precisely the original arrows. \\
\phantomsection\label{entry:symbol-registry-00808}\(\llbracket t\rrbracket_n\) & \texttt{\detokenize{\llbracket\rrbracket_n}} & categorical evaluation of a named algebraic term & P10.7 & meta & P10.7 & \thepage & An arrow \(n\to1\) in the given theory; substitution of terms is composition with their tuple of interpreted arrows. \\
\phantomsection\label{entry:symbol-registry-00809}\(\varepsilon_{\mathcal T},\rho\) & \texttt{\detokenize{\varepsilon_{\mathcal T},}\allowbreak\detokenize{\rho}} & isomorphism recovering a Lawvere theory & P10.7 & meta & P10.7 & \thepage & Evaluation from the named term theory is identity on objects; \(\rho\) denotes this fixed isomorphism when transporting free models. \\
\phantomsection\label{entry:symbol-registry-00810}\(H^*,\widehat H\) & \texttt{\detokenize{H^*,}\allowbreak\detokenize{\widehat H}} & restriction of models and translation of named terms & P10.7 & meta & P10.7 & \thepage & A theory morphism induces precomposition on models; its named-symbol translation commutes with the recovery isomorphisms. \\
\phantomsection\label{entry:symbol-registry-00811}\(R_n,\eta^n\) & \texttt{\detokenize{R_n,}\allowbreak\detokenize{\eta^n}} & representable finite free model and its generators & P10.7 & meta & P10.7 & \thepage & The model \(R_n=T(n,-)\) is free on \(n\); its generator indexed by \(i\) is \(p_i^n:n\to1\). \\
\phantomsection\label{entry:symbol-registry-00812}\(\theta_X,Q_X\) & \texttt{\detokenize{\theta_X,}\allowbreak\detokenize{Q_X}} & semantic term congruence and free equational algebra & P10.7 & meta & P10.7 & \thepage & The congruence requires equality under every evaluation into an all-carrier model; its quotient includes the empty term case. \\
\phantomsection\label{entry:symbol-registry-00813}\(\mathcal F_{\mathbb T},\eta_X\) & \texttt{\detokenize{\mathcal F_{\mathbb T},}\allowbreak\detokenize{\eta_X}} & free Lawvere-model functor and generator map & P10.7 & meta & P10.7 & \thepage & The free model has values \(Q_X^n\) after the fixed theory isomorphism; the generator maps \(x\) to the singleton tuple \(([x])\). \\
\phantomsection\label{entry:symbol-registry-00814}\(P(I),\top_I,\bot_I\) & \texttt{\detokenize{P,}\allowbreak\detokenize{\top_I,}\allowbreak\detokenize{\bot_I}} & predicate fibre and its bounds & P10.9 & meta & P10.9 & \thepage & A small Heyting algebra of predicates in context object \(I\); its order is semantic entailment in that context. \\
\phantomsection\label{entry:symbol-registry-00815}\(\mathsf{Heyt}\) & \texttt{\detokenize{\mathsf{Heyt}}} & category of Heyting algebras and homomorphisms & P10.9 & meta & P10.9 & \thepage & Maps preserve both bounds, finite meets and joins, and Heyting implication; reindexing takes values in this category. \\
\phantomsection\label{entry:symbol-registry-00816}\(f^*:P(I)\to P(J)\) & \texttt{\detokenize{f^*}} & predicate reindexing along a context map & P10.9 & meta & P10.9 & \thepage & For \(f:J\to I\), the contravariant functor sends predicates in \(I\) to predicates in \(J\), preserving every Heyting operation. \\
\phantomsection\label{entry:symbol-registry-00817}\(\pi_{I,A},f\times1_A\) & \texttt{\detokenize{\pi_{I,A},}\allowbreak\detokenize{f\times1_A}} & context projection and parameter substitution & P10.9 & meta & P10.9 & \thepage & The projection \(I\times A\to I\) forgets one variable; a parameter map lifts to the displayed product substitution square. \\
\phantomsection\label{entry:symbol-registry-00818}\(\exists_\pi,\forall_\pi\) & \texttt{\detokenize{\exists_\pi,}\allowbreak\detokenize{\forall_\pi}} & projection quantifier adjoints & P10.9 & meta & P10.9 & \thepage & The maps from \(P(I\times A)\) to \(P(I)\) are respectively left and right adjoints to projection reindexing. \\
\phantomsection\label{entry:symbol-registry-00819}\(\delta_A\) & \texttt{\detokenize{\delta_A}} & internal equality predicate & P10.9 & meta & P10.9 & \thepage & An element of \(P(A\times A)\) satisfying the full parameterized diagonal adjunction; no subobject representation is assumed. \\
\phantomsection\label{entry:symbol-registry-00820}\(e_{I,A},p_{I,A},q_{I,A},r_{I,A}\) & \texttt{\detokenize{e_{I,A},}\allowbreak\detokenize{p_{I,A},}\allowbreak\detokenize{q_{I,A},}\allowbreak\detokenize{r_{I,A}}} & maps for equality with parameters & P10.9 & meta & P10.9 & \thepage & Use \((I\times A)\times A\): duplicate the final variable, select either copy with its parameters, or select the equality pair. \\
\phantomsection\label{entry:symbol-registry-00821}\(L_e(\alpha)\) & \texttt{\detokenize{L_e}} & left adjoint for parameterized equality & P10.9 & meta & P10.9 & \thepage & Defined by \(p^*\alpha\wedge r^*\delta_A\), with \(L_e\dashv e^*\); its full order adjunction gives equality elimination. \\
\phantomsection\label{entry:symbol-registry-00822}\(\Delta_A,\tau_A\) & \texttt{\detokenize{\Delta_A,}\allowbreak\detokenize{\tau_A}} & diagonal and exchange of an equality pair & P10.9 & meta & P10.9 & \thepage & The arrows \(A\to A\times A\) and \(A\times A\to A\times A\) used to state reflexivity and symmetry. \\
\phantomsection\label{entry:symbol-registry-00823}\(\exists_{\pi_{1,A}}\top_{1\times A}=\top_1\) & \texttt{\detokenize{\exists_{\pi_{1,A}}}} & internal inhabitance of a sort & P10.9 & meta & P10.9 & \thepage & Equivalent by Beck--Chevalley to the corresponding condition in every parameter context; it implies reflection of predicate order, not of base-arrow equality. \\
\phantomsection\label{entry:symbol-registry-00824}\(P(X)=\mathcal P(X)\) & \texttt{\detokenize{\mathcal P}} & powerset instance of a predicate fibre & P10.9 & meta & P10.9 & \thepage & Boolean algebra of all subsets of a small set \(X\), including the one-element fibre for empty \(X\). \\
\phantomsection\label{entry:symbol-registry-00825}\(U\Rightarrow_X V\) & \texttt{\detokenize{\Rightarrow_X}} & Boolean implication between subsets & P10.9 & meta & P10.9 & \thepage & The subset \((X\setminus U)\cup V\) is the Heyting residual in the powerset fibre. \\
\phantomsection\label{entry:symbol-registry-00826}\(\exists_f,\forall_f\) & \texttt{\detokenize{\exists_f,}\allowbreak\detokenize{\forall_f}} & powerset quantifiers along an arbitrary function & P10.9 & meta & P10.9 & \thepage & Direct image and the condition that the entire fibre belongs to the predicate; empty fibres satisfy the universal condition. \\
\phantomsection\label{entry:symbol-registry-00827}\(\Gamma,\Phi\) & \texttt{\detokenize{\Gamma,}\allowbreak\detokenize{\Phi}} & variable context and hypothesis list in Pitts's calculus & P10.9 & meta & P10.9 & \thepage & The first is a finite list of distinct sorted variables; the second is a finite list of formulas in that context. \\
\phantomsection\label{entry:symbol-registry-00828}\(M:\sigma[\Gamma],M=N:\sigma[\Gamma]\) & \texttt{\detokenize{M:\sigma,}\allowbreak\detokenize{M=N:\sigma}} & term formation and external term-equation judgments & P10.9 & meta & P10.9 & \thepage & Formation is well-sortedness in the retained context; an equation is interpreted as equality of term arrows. \\
\phantomsection\label{entry:symbol-registry-00829}\(\Phi\vdash\psi[\Gamma]\) & \texttt{\detokenize{\vdash}} & contextual predicate sequent & P10.9 & meta & P10.9 & \thepage & The conjunction of the finite assumptions is below the conclusion in \(P(I_\Gamma)\); unused variable declarations are retained. \\
\phantomsection\label{entry:symbol-registry-00830}\(M=_\sigma N\) & \texttt{\detokenize{=_\sigma}} & equality as a first-order formula & P10.9 & meta & P10.9 & \thepage & Reindexes \(\delta_{A_\sigma}\) along the pair of term arrows; its truth need not reflect external equality of those arrows. \\
\phantomsection\label{entry:symbol-registry-00831}\(\varphi\mathbin{\&}\psi,\varphi\Rightarrow\psi,\mathsf{false}\) & \texttt{\detokenize{\&,}\allowbreak\detokenize{\Rightarrow,}\allowbreak\detokenize{\mathsf{false}}} & Pitts's propositional formula constructors & P10.9 & meta & P10.9 & \thepage & Conjunction, implication, and falsity at the formula level; their denotations are the corresponding Heyting operations. \\
\phantomsection\label{entry:symbol-registry-00832}\(\forall_\sigma(\varphi),\exists_\sigma(\varphi)\) & \texttt{\detokenize{\forall_\sigma,}\allowbreak\detokenize{\exists_\sigma}} & quantifiers on unary formula abstractions & P10.9 & meta & P10.9 & \thepage & Pitts's printed quantifier notation; the body \(\varphi(x)\) is interpreted in the context extended by \(x:\sigma\). \\
\phantomsection\label{entry:symbol-registry-00833}\((x)e,\varphi(M)\) & \texttt{\detokenize{(x)e,}\allowbreak\detokenize{\varphi(M)}} & metalevel formula abstraction and instantiation & P10.9 & meta & P10.9 & \thepage & If \(\varphi=(x)e\), then \(\varphi(M)=e[M/x]\); this is binding notation, not an additional first-order connective. \\
\phantomsection\label{entry:symbol-registry-00834}\(I_\Gamma\) & \texttt{\detokenize{I_\Gamma}} & interpreted finite variable context & P10.9 & meta & P10.9 & \thepage & A chosen product starting with terminal object \(1\), extended by one sort object for each declaration. \\
\phantomsection\label{entry:symbol-registry-00835}\(\llbracket M\rrbracket_\Gamma,\llbracket\varphi\rrbracket_\Gamma\) & \texttt{\detokenize{\llbracket\rrbracket_\Gamma}} & contextual term and predicate interpretations & P10.9 & meta & P10.9 & \thepage & A term denotes an arrow from \(I_\Gamma\) to its sort object; a formula denotes an element of \(P(I_\Gamma)\). \\
\phantomsection\label{entry:symbol-registry-00836}\(s:I_\Delta\to I_\Gamma\) & \texttt{\detokenize{s}} & interpreted simultaneous term substitution & P10.9 & meta & P10.9 & \thepage & Its components interpret one substituting term per variable of \(\Gamma\); formula substitution is reindexing along this map. \\
\phantomsection\label{entry:symbol-registry-00837}\(\mathcal U\) & \texttt{\detokenize{\mathcal U}} & universe used for categorical size bounds & P10.7-P10.9 & meta & P10.7 & \thepage & May be the ambient collection of all sets; a set-sized relativization has the explicitly stated Grothendieck closures, with existence treated as a hypothesis. \\
\phantomsection\label{entry:symbol-registry-00838}\(\mathrm{ND},\mathrm{ADJ}\) & \texttt{\detokenize{\mathrm{ND},}\allowbreak\detokenize{\mathrm{ADJ}}} & local abbreviations for two contextual rule presentations & P10.9 & meta & P10.9 & \thepage & The same theory and structural/equational basis, with either the natural-deduction logical rules or every direction of Pitts's adjoint bi-rules. \\
\phantomsection\label{entry:symbol-registry-00839}\(\mathrm{Eq0}\) & \texttt{\detokenize{\mathrm{Eq0}}} & parameter-free equality bi-rule & P10.9 & meta & P10.9 & \thepage & Both directions of Pitts's Remark 5.6.2; the displayed variable contexts have no hidden parameter prefix. \\
\phantomsection\label{entry:symbol-registry-00840}\((x)(x')e,\psi(M,N)\) & \texttt{\detokenize{(x)(x')e,}\allowbreak\detokenize{\psi(M,}\allowbreak\detokenize{N)}} & binary formula abstraction and simultaneous instantiation & P10.9 & meta & P10.9 & \thepage & Application replaces the two displayed argument variables simultaneously in the body, avoiding capture. \\
\phantomsection\label{entry:symbol-registry-00841}\(\widehat{\mathcal C}\) & \texttt{\detokenize{\widehat{\mathcal C}}} & presheaf category on a small base & P10.9 & meta & P10.9 & \thepage & The category of functors \(\mathcal C^{\mathrm{op}}\to\mathsf{Set}\) and natural transformations, under the declared size convention. \\
\phantomsection\label{entry:symbol-registry-00842}\(\mathrm y c,\mathrm y(u)\) & \texttt{\detokenize{\mathrm y}} & representable presheaf and the Yoneda map & P10.9 & meta & P10.9 & \thepage & The value at \(d\) is \(\mathcal C(d,c)\); an arrow between representing objects acts by postcomposition. \\
\phantomsection\label{entry:symbol-registry-00843}\(X(u),\alpha_c\) & \texttt{\detokenize{X(u),}\allowbreak\detokenize{\alpha_c}} & presheaf restriction and natural-transformation component & P10.9 & meta & P10.9 & \thepage & For \(u:d\to c\), restriction maps \(X(c)\) to \(X(d)\); the component \(\alpha_c\) maps \(X(c)\) to \(Y(c)\). \\
\phantomsection\label{entry:symbol-registry-00844}\(\operatorname{Sub}(X)\) & \texttt{\detokenize{\operatorname{Sub}}} & ordered set of subpresheaves & P10.9 & meta & P10.9 & \thepage & Restriction-stable families of component subsets; canonical images represent categorical subobjects and their factorization order. \\
\phantomsection\label{entry:symbol-registry-00845}\(0(c)=\varnothing,\ 1(c)=\{*\}\) & \texttt{\detokenize{0,}\allowbreak\detokenize{1}} & constant initial and terminal presheaves & P10.9 & meta & P10.9 & \thepage & These denote presheaves; numeric object names of a particular base category have their separately declared base-object role. \\
\phantomsection\label{entry:symbol-registry-00846}\(\operatorname{Im}(m),q,i\) & \texttt{\detokenize{\operatorname{Im},}\allowbreak\detokenize{q,}\allowbreak\detokenize{i}} & componentwise image and its factorization & P10.9 & meta & P10.9 & \thepage & The natural map factors through its image by a componentwise surjection followed by an inclusion; a monic map is isomorphic to that inclusion. \\
\phantomsection\label{entry:symbol-registry-00847}\(\Omega,\mathsf{true}:1\to\Omega\) & \texttt{\detokenize{\Omega,}\allowbreak\detokenize{\mathsf{true}}} & subobject classifier and truth map & P10.9 & meta & P10.9 & \thepage & Every monomorphism is a pullback along its unique characteristic arrow; the definition retains the full universal property. \\
\phantomsection\label{entry:symbol-registry-00848}\((Y^X)(c)\) & \texttt{\detokenize{Y^X}} & component of a presheaf exponential & P10.9 & meta & P10.9 & \thepage & The set \(\operatorname{Nat}(\mathrm y c\times X,Y)\), with restriction by precomposition; it need not be the pointwise external function set. \\
\phantomsection\label{entry:symbol-registry-00849}\(\operatorname{ev}_c,\Lambda h\) & \texttt{\detokenize{\operatorname{ev}_c,}\allowbreak\detokenize{\Lambda}} & presheaf evaluation and transposition & P10.9 & meta & P10.9 & \thepage & Evaluation uses the representing identity arrow; transposition uses every restriction of the parameter element. \\
\phantomsection\label{entry:symbol-registry-00850}\(\Omega(c),\Omega(u),t(c)\) & \texttt{\detokenize{\Omega(c),}\allowbreak\detokenize{\Omega(u),}\allowbreak\detokenize{t(c)}} & sieves their restrictions and the maximal sieve & P10.9 & meta & P10.9 & \thepage & A sieve is a set of arrows into \(c\) closed under precomposition; the maximal sieve contains the identity. \\
\phantomsection\label{entry:symbol-registry-00851}\(\chi_A:X\to\Omega\) & \texttt{\detokenize{\chi_A}} & characteristic map of a subpresheaf & P10.9 & meta & P10.9 & \thepage & At \(x\in X(c)\), it is the sieve of arrows \(u:d\to c\) whose restriction \(X(u)(x)\) belongs to \(A(d)\). \\
\phantomsection\label{entry:symbol-registry-00852}\(P(X)=\Omega^X\) & \texttt{\detokenize{P(X)}} & power object of a presheaf & P10.9 & meta & P10.9 & \thepage & Here \(X\) is a presheaf and \(P(X)\) is an object of \(\widehat{\mathcal C}\), distinct from the earlier powerset-fibre usage for sets. \\
\phantomsection\label{entry:symbol-registry-00853}\(\in_X\hookrightarrow P(X)\times X\) & \texttt{\detokenize{\in_X}} & membership subobject of a power object & P10.9 & meta & P10.9 & \thepage & The pullback of truth along evaluation; under stage classification membership tests whether \((1_c,x)\) lies in the represented subpresheaf. \\
\phantomsection\label{entry:symbol-registry-00854}\(S|u\) & \texttt{\detokenize{S|u}} & restriction of a classified power-object element & P10.9 & meta & P10.9 & \thepage & The pullback of \(S\subseteq\mathrm y c\times X\) along \(\mathrm y u\times1_X\). \\
\phantomsection\label{entry:symbol-registry-00855}\(S_z\) & \texttt{\detokenize{S_z}} & subpresheaf determined by a relation and a parameter element & P10.9 & meta & P10.9 & \thepage & For \(z\in Z(c)\), its component consists of pairs \((u,x)\) with \((Z(u)(z),x)\) in the original relation. \\
\end{longtable}
\endgroup
\Needspace{.4\textheight}
\subsection*{Part 12 scopes}
\begingroup\scriptsize\setlength{\tabcolsep}{3pt}
\begin{longtable}{@{}p{0.9in}p{1.3in}p{1.0in}p{0.5in}p{0.3in}p{0.4in}p{0.35in}p{1.65in}@{}}
\toprule
\textbf{Glyph} & \textbf{Macro} & \textbf{Role} & \textbf{Scope} & \textbf{Prec} & \textbf{First} & \textbf{Reg. page} & \textbf{Principal uses / notes} \\
\midrule
\endfirsthead
\toprule
\textbf{Glyph} & \textbf{Macro} & \textbf{Role} & \textbf{Scope} & \textbf{Prec} & \textbf{First} & \textbf{Reg. page} & \textbf{Principal uses / notes} \\
\midrule
\endhead
\midrule\multicolumn{8}{r}{\emph{continued}}\\\endfoot
\bottomrule\endlastfoot
\phantomsection\label{entry:symbol-registry-00609}\(\mathsf P\) & \texttt{\detokenize{\mathsf P}} & deterministic polynomial-time language class & P12.1 & n/a & P12.1 & \thepage & Complexity class reused in the P12 SAT and certificate comparison; disjoint from the P4.3 pairing-graph formula. \\
\phantomsection\label{entry:symbol-registry-00610}\(\mathsf{SAT}\) & \texttt{\detokenize{\mathsf{SAT}}} & satisfiable propositional encodings & P12.1 & n/a & P12.1 & \thepage & SAT instances under the exact clause-set and encoding conventions fixed in the automation chapter. \\
\phantomsection\label{entry:symbol-registry-00611}\(\mathsf{DPLL}\) & \texttt{\detokenize{\mathsf{DPLL}}} & Davis--Putnam--Logemann--Loveland method identifier & P12.1 & n/a & P12.1 & \thepage & Exact recursive finite-clause search method defined before its correctness theorem and trace. \\
\phantomsection\label{entry:symbol-registry-00612}\(\#\mathsf{SAT}\) & \texttt{\detokenize{\#\mathsf{SAT}}} & satisfying-assignment counting function & P12.1 & n/a & P12.1 & \thepage & Maps a propositional encoding to its number of satisfying total assignments; not a decision language. \\
\phantomsection\label{entry:symbol-registry-00613}\(\#\mathsf P\) & \texttt{\detokenize{\#\mathsf P}} & counting complexity class & P12.1 & n/a & P12.1 & \thepage & Functions counting accepting witnesses of a polynomial-time decidable relation. \\
\phantomsection\label{entry:symbol-registry-00614}\(\mathsf{MaxSAT}\) & \texttt{\detokenize{\mathsf{MaxSAT}}} & maximum satisfiability optimization problem & P12.1 & n/a & P12.1 & \thepage & Optimization problem maximizing the number or declared weight of satisfied clauses under the local variant. \\
\phantomsection\label{entry:symbol-registry-00615}\(\mathsf{SMT}\) & \texttt{\detokenize{\mathsf{SMT}}} & satisfiability-modulo-theories identifier & P12.2 & n/a & P12.2 & \thepage & Decision problem for quantifier-free formulas interpreted in a fixed background theory. \\
\phantomsection\label{entry:symbol-registry-00616}\(\mathsf{EUF}\) & \texttt{\detokenize{\mathsf{EUF}}} & equality with uninterpreted functions theory identifier & P12.2 & n/a & P12.2 & \thepage & Quantifier-free equality theory whose function symbols have no axioms beyond equality congruence. \\
\phantomsection\label{entry:symbol-registry-00617}\(M\) & \texttt{\detokenize{M}} & Slagle ground-interpretation metavariable & P12.3 & n/a & P12.3 & \thepage & The fixed ground interpretation in semantic resolution; exact source role, disjoint from the P7 term-metavariable scope. \\
\phantomsection\label{entry:symbol-registry-00618}\(N,E_1,\ldots,E_q\) & \texttt{\detokenize{N,}\allowbreak\detokenize{E_i}} & semantic-clash nucleus and electron clauses & P12.3 & n/a & P12.3 & \thepage & Exact Slagle roles: \(N\) is the nucleus and the \(E_i\) are electrons tested against the fixed ground interpretation \(M\). \\
\phantomsection\label{entry:symbol-registry-00619}\(R_1,\ldots,R_{q+1}\) & \texttt{\detokenize{R_i}} & semantic-clash intermediate resolvents & P12.3 & n/a & P12.3 & \thepage & Exact Slagle resolvent chain with \(R_1=N\) and final false \(M\)-resolvent \(R_{q+1}\); distinct from frame relations. \\
\phantomsection\label{entry:symbol-registry-00620}\(\mathsf{LCF}\) & \texttt{\detokenize{\mathsf{LCF}}} & LCF trusted-kernel architecture identifier & P12.5 & n/a & P12.5 & \thepage & Interactive-prover architecture in which an abstract theorem type can be constructed only by a small trusted kernel. \\
\phantomsection\label{entry:symbol-registry-00621}\(\mathsf{PTS}\) & \texttt{\detokenize{\mathsf{PTS}}} & pure type system identifier & P12.5 & n/a & P12.5 & \thepage & Parameterized typed calculus determined by its sorts axioms and allowed product triples. \\
\phantomsection\label{entry:symbol-registry-00622}\(S=(\mathcal S,\mathcal A,\mathcal R);\lambda S\) & \texttt{\detokenize{S,}\allowbreak\detokenize{\mathcal S,}\allowbreak\detokenize{\mathcal A,}\allowbreak\detokenize{\mathcal R,}\allowbreak\detokenize{\lambda S}} & Barendregt pure-type-system specification and calculus & P12.5 & n/a & P12.5 & \thepage & The exact source datum consists of sorts sort-typing axioms and dependent-product triples; lambda S names the calculus it determines. \\
\phantomsection\label{entry:symbol-registry-00623}\(\ast,\Box;(s_1,s_2)\) & \texttt{\detokenize{\ast,}\allowbreak\detokenize{\Box;(s_1,}\allowbreak\detokenize{s_2)}} & Barendregt lambda-cube sorts and product-sort pair & P12.5 & n/a & P12.5 & \thepage & Exact source notation: the pair abbreviates the triple (s1,s2,s2); the locally declared PTS role distinguishes Box and star from modal and algebraic uses. \\
\phantomsection\label{entry:symbol-registry-00624}\(\lambda{\to},\lambda\mathrm P\) & \texttt{\detokenize{\lambda{\to},}\allowbreak\detokenize{\lambda\mathrm P}} & Barendregt lambda-cube calculus identifiers & P12.5 & n/a & P12.5 & \thepage & Exact source identifiers for the base simply typed vertex and the types-on-terms vertex. \\
\phantomsection\label{entry:symbol-registry-00625}\(\lambda 2,\lambda\mathrm P2\) & \texttt{\detokenize{\lambda 2,}\allowbreak\detokenize{\lambda\mathrm P2}} & Barendregt lambda-cube calculus identifiers & P12.5 & n/a & P12.5 & \thepage & Exact source identifiers for the terms-on-types vertex and its types-on-terms extension. \\
\phantomsection\label{entry:symbol-registry-00626}\(\lambda\underline{\omega},\lambda\mathrm P\underline{\omega}\) & \texttt{\detokenize{\lambda\underline{\omega},}\allowbreak\detokenize{\lambda\mathrm P\underline{\omega}}} & Barendregt lambda-cube calculus identifiers & P12.5 & n/a & P12.5 & \thepage & Exact source identifiers for the types-on-types vertex and its types-on-terms extension. \\
\phantomsection\label{entry:symbol-registry-00627}\(\lambda\omega,\lambda\mathrm P\omega=\lambda\mathrm C\) & \texttt{\detokenize{\lambda\omega,}\allowbreak\detokenize{\lambda\mathrm P\omega,}\allowbreak\detokenize{\lambda\mathrm C}} & Barendregt lambda-cube calculus identifiers & P12.5 & n/a & P12.5 & \thepage & Exact source identifiers for the terms-on-types plus types-on-types vertex and the all-three-axes vertex; lambda P omega equals lambda C. \\
\phantomsection\label{entry:symbol-registry-00628}\(\mathsf{CIC}\) & \texttt{\detokenize{\mathsf{CIC}}} & Calculus of Inductive Constructions identifier & P12.5 & n/a & P12.5 & \thepage & Calculus of Constructions extended by the separately described inductive families and eliminators. \\
\phantomsection\label{entry:symbol-registry-00629}\(\alpha,\beta,\gamma\) & \texttt{\detokenize{\alpha,}\allowbreak\detokenize{\beta,}\allowbreak\detokenize{\gamma}} & type metavariables & P12.5 & n/a & P12.5 & \thepage & Types in the reused STLC and lambda-cube examples; this scope is disjoint from ordinal and path-formula roles. \\
\phantomsection\label{entry:symbol-registry-00630}\(\alpha\to_{\mathrm{ty}}\beta\) & \texttt{\detokenize{\to_{\mathrm{ty}}}} & object-language function-type former & P12.5 & 6 & P12.5 & \thepage & Right-associative arrow of the exact reused STLC syntax. \\
\phantomsection\label{entry:symbol-registry-00631}\(M{:}\alpha\) & \texttt{\detokenize{:}} & type ascription & P12.5 & meta & P12.5 & \thepage & The local STLC example fixes the term and type categories; distinct from labelled formulas. \\
\phantomsection\label{entry:symbol-registry-00632}\(\Theta_{\mathrm{ty}}\) & \texttt{\detokenize{\Theta_{\mathrm{ty}}}} & term-variable typing context & P12.5 & n/a & P12.5 & \thepage & Finite list of distinct term-variable declarations in the reused STLC example. \\
\phantomsection\label{entry:symbol-registry-00633}\(\mathsf{STLC}\) & \texttt{\detokenize{\mathsf{STLC}}} & arrow-only simply typed lambda-calculus identifier & P12.5 & n/a & P12.5 & \thepage & Exact calculus reused from sheet:stlc in the mechanization example. \\
\phantomsection\label{entry:symbol-registry-00634}\(\Theta_{\mathrm{ty}}\derives_{\mathsf{STLC}}M{:}\alpha\) & \texttt{\detokenize{\derives_{\mathsf{STLC}}}} & STLC typing judgment & P12.5 & meta & P12.5 & \thepage & Exact reused arrow-only typing judgment; no new P12 calculus is introduced. \\
\end{longtable}
\endgroup
\end{landscape}


\subsection{Classical-logic core}
\label{sec:p01-classical-symbols}

For direct lookup, Table~\ref{tab:p01-classical-symbols} extracts the
classical core from the same CSV register.  The table summarizes
existing rows; it does not declare a second notation layer.

\begingroup\small
\begin{longtable}{@{}p{0.20\textwidth}p{0.34\textwidth}p{0.38\textwidth}@{}}
\caption{Classical-logic sub-register, derived from
\texttt{registry/symbols.csv}.}
\label{tab:p01-classical-symbols}\\
\toprule
Forms & Standing & Scope and binding \\
\midrule
\endfirsthead
\toprule
Forms & Standing & Scope and binding \\
\midrule
\endhead
\(\mathcal L;\ x,y,z;\ r,t\) & signature or language;
object-language variables; term metavariables & Global after a language is
fixed; \(r,t\in\operatorname{Term}(\mathcal L)\). \\
\(\varphi,\psi,\chi;\ \sigma,\tau\) & formula; sentence metavariables &
Global after a language is fixed; no operator precedence. \\
\(\Gamma,\Delta,\Sigma,\Pi;\ T,U\) & formula contexts; theories &
Global; the active calculus declares whether a context is a set, multiset,
or sequence. \\
\(\mathcal M,\mathcal N;\ v;\ s\) & structures; propositional valuation;
variable assignment & Global semantic metavariables; their types are fixed
by the active language and semantics. \\
\(\bot,\top;\ \neg,\land,\lor,\to,\leftrightarrow\) & truth constants;
propositional connectives & Global classical core; the constants have no
rank, while the five connectives have ranks 3, 4, 5, 6, 7 respectively;
implication is right-associative. \\
\(\forall x,\exists x\) & object-language quantifiers & Global; the active
language fixes the bound variable's sort or type.  Prefix rank 3. \\
\(v\entails\varphi;\ \mathcal M,s\entails\varphi;
\ \mathcal M\entails\sigma;\ \Gamma\entails\varphi\) & propositional
satisfaction; satisfaction under an assignment; truth in a structure;
semantic consequence & Metalevel relations; the sentence-only third form
suppresses assignments only after assignment-independence. \\
\(\Gamma\derives_{S}\varphi;\ \Gamma\Rightarrow\Delta\) & derivability in
calculus \(S\); sequent & Metalevel; the system subscript is mandatory
outside a declared local abbreviation. \\
\bottomrule
\end{longtable}
\endgroup

\section{Binding and precedence}
\label{sec:p01-binding-precedence}
\index{precedence}\index{binding}

The grammar declared by the active language is always primary.  After unique
readability has been established, Table~\ref{tab:p01-precedence} licenses the
parentheses omitted in expository formulas.  Smaller numerical ranks bind
more tightly.  Operators sharing a rank are governed by the association
column; unrelated subsystem symbols share a rank only where their scopes are
disjoint.

\begingroup\small
\begin{longtable}{@{}p{0.08\textwidth}p{0.31\textwidth}p{0.21\textwidth}p{0.31\textwidth}@{}}
\caption{Handbook binding and precedence, derived from the symbol registry.}
\label{tab:p01-precedence}\\
\toprule
Rank & Symbols or forms & Association & Reading rule \\
\midrule
\endfirsthead
\toprule
Rank & Symbols or forms & Association & Reading rule \\
\midrule
\endhead
1 & function and predicate application & fixed by arity & Arguments belong to
the displayed function or predicate symbol. \\
2 & \(=,\leq,>\) and other declared atomic relations & nonassociative & A
relation combines terms into an atomic formula. \\
2.5 & \(!A\) in the linear-logic scope & prefix & The exponential governs
the immediately following formula. \\
3 & \(\neg,\Box,\Diamond,\mathsf X,\mathsf F,\mathsf G\);
\(K_i,E_G,C_G,\forall x,\exists x\) & prefix & The operator or quantifier governs the
immediately following formula unless parentheses extend its scope. \\
3.5 & \(A\mathbin{\mathsf U}B\) & nonassociative & Temporal until binds
below unary temporal operators and above binary propositional connectives. \\
4 & \(\land\), and \(\otimes\) in its linear scope & left associative &
Conjunction-level operators bind above disjunction. \\
5 & \(\lor\) & left associative & Disjunction binds above implication. \\
6 & \(\to\), and \(\multimap\) in its linear scope & right associative &
Implication-level operators associate to the right. \\
7 & \(\leftrightarrow\) & nonassociative & Repeated biconditionals require
parentheses. \\
meta & \(\defeq,\Leftrightarrow,\entails,\derives\) and their declared
subscripted forms & not object-language operators & These signs relate
expressions or statements at the metalinguistic level. \\
\bottomrule
\end{longtable}
\endgroup

Worked disambiguations fix the convention rather than relying on visual
guesswork:
\[
\begin{aligned}
 \neg P\land Q\to R
   &\quad\text{means}\quad ((\neg P)\land Q)\to R,\\
 P\to Q\to R
   &\quad\text{means}\quad P\to(Q\to R),\\
 \forall x\,P(x)\to Q(x)
   &\quad\text{means}\quad (\forall x\,P(x))\to Q(x),\\
 !A\otimes B\multimap C
   &\quad\text{means}\quad ((!A)\otimes B)\multimap C,\\
 \mathsf X A\mathbin{\mathsf U}B\lor C
   &\quad\text{means}\quad (\mathsf X A\mathbin{\mathsf U}B)\lor C.
\end{aligned}
\]
Write \(\forall x\,(P(x)\to Q(x))\) when the implication is inside the
quantifier's scope.  Restore parentheses whenever a reader would otherwise
need to reconstruct a local grammar or a scope declaration.

\section{The sixteen binary Boolean truth functions}
\label{sec:p01-sixteen-functions}
\index{truth function!binary}\index{Sheffer function}\index{Peirce arrow}

For inputs \((p,q)=(0,0),(0,1),(1,0),(1,1)\), a binary Boolean truth
function is determined by its four output bits.  Table~\ref{tab:p01-sixteen}
lists all \(2^{2^2}=16\) functions.  The formulas in the fourth column are
definitions using the handbook's registered connectives; historical glyphs
mentioned in the last column are literature variants, not additional active
handbook operators.

\begingroup\scriptsize\setlength{\tabcolsep}{3pt}
\begin{longtable}{@{}c p{0.20\textwidth} cccc p{0.23\textwidth}p{0.22\textwidth}@{}}
\caption{All binary Boolean truth functions.}
\label{tab:p01-sixteen}\\
\toprule
Code & Standard name & 00 & 01 & 10 & 11 & Defining formula & Notes \\
\midrule
\endfirsthead
\toprule
Code & Standard name & 00 & 01 & 10 & 11 & Defining formula & Notes \\
\midrule
\endhead
0000 & contradiction / constant false & 0&0&0&0 & \(\bot\) & independent of both inputs \\
0001 & conjunction & 0&0&0&1 & \(p\land q\) & --- \\
0010 & material nonimplication & 0&0&1&0 & \(p\land\neg q\) & negation of \(p\to q\) \\
0011 & first projection & 0&0&1&1 & \(p\) & independent of \(q\) \\
0100 & converse nonimplication & 0&1&0&0 & \(\neg p\land q\) & negation of \(q\to p\) \\
0101 & second projection & 0&1&0&1 & \(q\) & independent of \(p\) \\
0110 & exclusive disjunction & 0&1&1&0 & \(\neg(p\leftrightarrow q)\) & often XOR or \(\oplus\) \\
0111 & disjunction & 0&1&1&1 & \(p\lor q\) & inclusive disjunction \\
1000 & joint denial / NOR & 1&0&0&0 & \(\neg(p\lor q)\) & Peirce arrow; often \(p\downarrow q\) \\
1001 & biconditional / equivalence & 1&0&0&1 & \(p\leftrightarrow q\) & Boolean equality \\
1010 & negation of second input & 1&0&1&0 & \(\neg q\) & independent of \(p\) \\
1011 & converse implication & 1&0&1&1 & \(q\to p\) & false only at \((0,1)\) \\
1100 & negation of first input & 1&1&0&0 & \(\neg p\) & independent of \(q\) \\
1101 & material implication & 1&1&0&1 & \(p\to q\) & false only at \((1,0)\) \\
1110 & alternative denial / NAND & 1&1&1&0 & \(\neg(p\land q)\) & Sheffer function; often \(p\uparrow q\) or \(p\mid q\) \\
1111 & tautology / constant true & 1&1&1&1 & \(\top\) & independent of both inputs \\
\bottomrule
\end{longtable}
\endgroup

The Sheffer stroke (NAND) and Peirce arrow (NOR) are each functionally
complete.
\index{Sheffer stroke}\index{NAND}\index{NOR}
The vertical-bar glyph is unavailable as active handbook notation, so this
definability note uses eliminable metalinguistic names:
\(\operatorname{NAND}(p,q)\defeq\neg(p\land q)\) and
\(\operatorname{NOR}(p,q)\defeq\neg(p\lor q)\).  These names designate
exactly rows 1110 and 1000, are not object-language connectives, and occur
only here.  Thus
\[
\begin{aligned}
\neg p&\Leftrightarrow\operatorname{NAND}(p,p),\\
p\land q&\Leftrightarrow
 \operatorname{NAND}(\operatorname{NAND}(p,q),
                     \operatorname{NAND}(p,q));\\
\neg p&\Leftrightarrow\operatorname{NOR}(p,p),\\
p\lor q&\Leftrightarrow
 \operatorname{NOR}(\operatorname{NOR}(p,q),
                    \operatorname{NOR}(p,q)).
\end{aligned}
\]
The arbitrary-connective criterion appears with Post's five maximal clones
in Part~2.

\section{Representative notation in the literature}
\label{sec:p01-notation-variants}
\index{notation!historical variants}

Table~\ref{tab:p01-notation-variants} records representative printed usage
from identified source traditions.  It is bibliographic rather than
normative: authors within a tradition and later editions vary, and the table
alone changes no live presentation.  When a chapter explicitly presents one
of those named sources, however, that source's own syntax governs locally.

\begingroup\scriptsize\setlength{\tabcolsep}{3pt}
\begin{longtable}{@{}p{0.16\textwidth}p{0.18\textwidth}p{0.18\textwidth}p{0.18\textwidth}p{0.18\textwidth}@{}}
\caption{Representative notation in four named primary sources.}
\label{tab:p01-notation-variants}\\
\toprule
Role (handbook key) & Whitehead--Russell, \emph{Principia Mathematica} I (1910) & {\L}ukasiewicz, \emph{Elementy} (1929) & Hilbert--Ackermann, \emph{Grundz\"uge}, 1st ed. (1928) & Gentzen, \emph{Untersuchungen I} (1935) \\
\midrule
\endfirsthead
\toprule
Role (handbook key) & Whitehead--Russell (1910) & {\L}ukasiewicz (1929) & Hilbert--Ackermann (1928) & Gentzen (1935) \\
\midrule
\endhead
negation (\(\neg A\)) & \(\mathord{\sim}p\) & \(Np\) & \(\overline X\) & \(\neg A\) \\
conjunction (\(A\land B\)) & \(p\mathbin{\boldsymbol{\cdot}}q\) & \(Kpq\) & \(X\mathbin{\&}Y\) & \(A\mathbin{\&}B\) \\
inclusive disjunction (\(A\lor B\)) & \(p\lor q\) & \(Apq\) & \(X\lor Y\) & \(A\lor B\) \\
implication (\(A\to B\)) & \(p\supset q\) & \(Cpq\) & \(X\to Y\) & \(A\supset B\) \\
biconditional (\(A\leftrightarrow B\)) & \(p\equiv q\) & \(Epq\) & \(X\sim Y\) & \(A\supset\subset B\), defined abbreviation \\
universal quantification (\(\forall x\,A\)) & \((x).\,\phi x\) & \(\Pi p\,\varphi(p)\) & \((x)A(x)\) & \(\forall x\,A\) \\
existential quantification (\(\exists x\,A\)) & \((\exists x).\,\phi x\) & \(\Sigma p\,\varphi(p)\) & \((E x)A(x)\) & \(\exists x\,A\) \\
\bottomrule
\end{longtable}
\endgroup

The entries reproduce the displayed forms in the cited works and editions;
they compare glyphs and argument placement, not the types of bound variables,
primitive bases, or every notation used by the authors elsewhere.  In the
cited {\L}ukasiewicz passage, \(p\) is a propositional variable.  Gentzen's
\(A\supset\subset B\) is a defined abbreviation, not a primitive formula
former, and his sequent \(\to\) is an auxiliary sign rather than implication.
Primary-source locators are recorded in
\cite{WhiteheadRussell1910,Lukasiewicz1929,HilbertAckermann1928,Gentzen1934a}.

\section{Letter-family conventions}
\label{sec:p01-letter-conventions}
\index{Greek letters}\index{Fraktur letters}\index{blackboard bold}

Typeface alone never assigns standing.  Table~\ref{tab:p01-letter-families}
summarizes the active family-level convention; the generated registry above
remains authoritative for individual symbols and scopes.

\begingroup\small
\begin{longtable}{@{}p{0.16\textwidth}p{0.27\textwidth}p{0.21\textwidth}p{0.27\textwidth}@{}}
\caption{Greek, Fraktur, calligraphic, and blackboard-letter conventions.}
\label{tab:p01-letter-families}\\
\toprule
Family & Registered examples & Scope & Convention \\
\midrule
\endfirsthead
\toprule
Family & Registered examples & Scope & Convention \\
\midrule
\endhead
Greek lowercase & \(\varphi,\psi,\chi;\ \sigma,\tau\) & global after a language is fixed & formula and sentence metavariables, respectively \\
Greek lowercase & \(\alpha,\beta\) & P6 and locally P7; separately P7 type passages & ordinal metavariables; type metavariables only under an explicit disjoint local declaration \\
Greek lowercase & \(\theta,\rho,\eta\) & P12 resolution & simultaneous term substitutions, not formula metavariables \\
Greek uppercase & \(\Gamma,\Delta,\Sigma,\Pi\) & global after a calculus is fixed & formula contexts; the active sheet declares set, multiset, or sequence standing \\
Greek uppercase & \(\Sigma^0_n,\Pi^0_n;\ \boldsymbol\Sigma^0_n,
\boldsymbol\Pi^0_n,\boldsymbol\Delta^0_n\) & P4--P5 & syntactic
formula classes; definability classes of relations, respectively; boldface
has no unscoped Borel-pointclass reading here \\
Calligraphic & \(\mathcal L;\ \mathcal M,\mathcal N;\ \mathcal A_{\mathrm{in}},\mathcal A_{\mathrm{tape}}\) & registered global or part-local rows & languages, structures, and machine alphabets as declared; the typeface by itself has no role \\
Blackboard bold & \(\mathbb N;\ \mathbb P\) & global natural-number convention; P6 forcing & standard natural numbers and a forcing poset, respectively \\
Fraktur & none globally assigned & all parts & no default standing; a future use requires a scoped registry row before substantive use \\
Roman or sans serif & \(\operatorname{FV},\operatorname{dom};\ \mathsf{PA},\mathsf X\) & the registered operation or system scope & named operations and fixed systems/operators, never generic variables \\
\bottomrule
\end{longtable}
\endgroup

The table deliberately records an unassigned Fraktur family rather than
inventing a role.  That absence is a usable convention: typography cannot
create a new mathematical category without a definition and a scoped
registry entry.
